\documentclass[11pt,a4paper]{article}
\usepackage[T1]{fontenc}
\usepackage[utf8]{inputenc}
\usepackage{lmodern,amsmath,amssymb,amsthm,mathtools,microtype,mathrsfs}
\usepackage[margin=25mm,top=23mm,bottom=25mm]{geometry}
\usepackage{booktabs,array,longtable,enumitem,xcolor,fancyhdr}
\usepackage[colorlinks=true,linkcolor=blue!45!black,urlcolor=blue!45!black,citecolor=blue!45!black]{hyperref}
\hypersetup{pdftitle={Reading the Alpoge--Claude RH and Dirichlet-GRH proof: working version 0.55},pdfauthor={ChatGPT},pdfsubject={Partial human-readable reconstruction with explicit arithmetic inputs},pdfkeywords={Riemann hypothesis, Lean, theta lifts, prime sums, Arakelov slopes}}
\setlength{\parindent}{0pt}
\setlength{\parskip}{5pt plus 1pt}
\setlist{topsep=3pt,itemsep=3pt,leftmargin=*}
\setlength{\headheight}{14pt}
\pagestyle{fancy}\fancyhf{}
\fancyhead[L]{\small RH and Dirichlet GRH: a reconstruction}
\fancyhead[R]{\small Working version 0.55}
\fancyfoot[C]{\thepage}
\renewcommand{\headrulewidth}{0.3pt}
\numberwithin{equation}{section}
\newtheorem{theorem}{Theorem}[section]
\newtheorem{proposition}[theorem]{Proposition}
\newtheorem{lemma}[theorem]{Lemma}
\newtheorem{corollary}[theorem]{Corollary}
\theoremstyle{remark}\newtheorem{remark}[theorem]{Remark}
\newcommand{\Rea}{\operatorname{Re}}
\newcommand{\HS}{\mathrm{HS}}
\newcommand{\tr}{\operatorname{tr}}
\newcommand{\codim}{\operatorname{codim}}
\newcommand{\covol}{\operatorname{covol}}
\newcommand{\OK}{\mathcal O_K}
\newcommand{\src}[2]{\href{https://github.com/CocoSgt/riemann-hypothesis-lean/blob/7cbf7a0c32fca1badf05931879a37c3f8a2b03b4/src/#1}{\texttt{#2}}}
\newcommand{\status}[1]{\textbf{Status:} #1}
\begin{document}
\begin{center}
{\LARGE\bfseries RH and Dirichlet GRH: new human proofs}\\[7pt]
{\large Version 0.55 --- 23 November 2026}\\[5pt]
{\large High twists, arithmetic local charts, stabilizer volumes and invariant matrix descent}
\end{center}
\textbf{Scope.} This edition contains the current audit and all new chapters.
It is a partial human reconstruction of the accepted formal proof.
Earlier named claims are in the cumulative manuscript and preserved components.
Each result below states its original inputs and exact remaining producers.
\setcounter{tocdepth}{1}
\tableofcontents
\clearpage
\begingroup\emergencystretch=3em
\section{Current proof obligations and completed suppliers: version 0.55}
\label{sec:dependency-audit-v55}

This continues the human reconstruction with the formal proof accepted
as correct. The selected sufficient route has 121 nodes and 161 edges.
There are 24 open whole obligations: 17 class C and 7 class B.
Of these, 17 are shared by ordinary RH and unrestricted Dirichlet
GRH; 7 are additional Dirichlet-GRH inputs, and none are RH-only.
The complete version-0.49 \texttt{GLOBAL\_\allowbreak LOCAL\_\allowbreak GAIN} closure remains
in force.

The four chapters contain 87 new numbered claims, bringing the
cumulative inventory to 3538. They prove substantial original
mathematical statements with their objects and assumptions preserved.
The geometry chapter gives both all-cover Serre first-cohomology
vanishing and the high-power homogeneous-polynomial statement.
The arithmetic chapter constructs the actual local objects and proves
their completion, flatness, residual types and representation charts
in ranks two and three. The volume chapter proves the complete
stabilizer comparison from the original smooth Siegel--Weil premises,
then supplies all three actual volume fields on the GMR19 setting.
The period chapter proves the full inhomogeneous fixed-matrix descent,
including the generic additive and matrix bijectivity arguments.

The precise remaining inputs have changed. The independent geometric
polarization and multiplication comparisons, invariant-function
descent, principal rank and square-law statements remain; Serre
vanishing itself is now supplied. Arithmetic dimension bounds,
component chains, density and seed incidence, the rank-three corner,
propagation and \texttt{Off\allowbreak Fin} remain; the local objects and their
flat representation charts are now supplied. The arithmetic
\texttt{In\allowbreak Volumes} fields are now supplied at the original
Siegel--Weil premise, while the distinct reference unipotent
discriminant calculation, \texttt{Jac\allowbreak Gauge}/\texttt{Place\allowbreak Split},
\texttt{Prod\allowbreak Exact} and the other mass/counting fields remain.
The actual Wach prism is now the only period-ring supplier retained
before conditional full \texttt{Wach\allowbreak Ex}. Original tied period data,
crystalline realization and the later faithful-invariant, etale-basis,
Herr and quantitative comparison work remain explicit.

The proofs also make the order of choices visible. The finite graded
module and syzygies are fixed before the large degree. The local
arithmetic maps return to the same marked frame and residue point.
The positive theta test is used with the same character and measures
in both identities, and finite stabilizer mass is obtained before
cancellation. The period approximation uses the original sharp-root
coordinates and completed fixed field. Its shorter argument has a
full return to the source's generator subring, so an omitted
cohomology detour is not a new assumption.

A class-B claim is a written human deduction at its stated scope.
The classifications and counts are inventories of this selected route,
not a minimality assertion or a completion percentage. Earlier proof
records and historical text are retained unchanged. The following
register governs the current remaining scope.

\subsection{Complete current whole-obligation register}
\paragraph{1. GLOBAL\_\allowbreak DENOMINATOR — C — Actual LinkN local table and binding producer}

Closed(Des $\mathcal C$ 3) LinkN: the actual etaSum is bounded with exact dZ0 and same-host vvTerm, subject only to the literal ShaFinOdd and summability guards.

\emph{Shared objects and choice order.} one FamP; one class $\mathcal A$ and its target-compatible subclass; one coefficient field per family datum; one x and i at each application; original selected rows, words and levels; thresholds chosen before the common family; actual D.nd; actual etaJ and Aj; sha\_\allowbreak eq

\emph{Completed inputs (version 0.55).} IFD1–IFD26 (full original inhomogeneous invariant matrix descent; exact one-period-supplier Wach return), WIS1–WIS26 (full four-clause actual inclusion and fixed/near-parameter completeness; exact two-supplier Wach return), WLD1-WLD26 (full generic Wach lattice constructors; exact conditional three-supplier all-layer return), SSC1-SSC17 (full original Selmer SignatureChoice and family setup retaining only coherent Wach comparison), DFD1-DFD15, DNE1-DNE14 (full original ConclF-to-LinkN consumer)

\emph{Remaining producer scope.} IFD supplies the full original inhomogeneous cyclotomic-kernel-fixed Frobenius matrix descent, not only the homogeneous case required by WLD26. It proves full generic const\_\allowbreak dvd/near\_\allowbreak closed, exact actual invariant tilt approximation, eventual Frobenius approximation, and full abstract additive and matrix quotient bijectivity. Exact fixed sharp-root coordinates and LS1 completed fixed fields avoid an unnecessary cyclotomic Tate cohomology detour; the original coefficient field, action, parameter, finite basis and p=2 case are preserved. WIS inclusion, saturation and both completeness structures are completed inputs. Exactly one period-ring supplier remains before conditional full WachEx: the actual Wach prism with all prescribed maps, scalar identity and both coefficient decompositions. IFD26/WLD26 retains original LocTied/IsStd and Bhatt-Scholze crystalline realization. Full all-layer WachEx is not yet unconditional. Separate FaithH, EtBasis, Herr class maps and quantitative comparison remain for HW/Fon/Cmp; no extra outer-layer transition system is introduced. Then every original local row/table bound must hold in one common ConclF family, retaining ShaFinOdd, exceptional sets, vvTerm summability and uniform thresholds. SSC, DFD12 and DNE remain available, with DNE K=2000(|Ca|+1)\textasciicircum{}3 fixed before the family. Whole GLOBAL\_\allowbreak DENOMINATOR remains open.

\emph{Source.} \href{https://github.com/CocoSgt/riemann-hypothesis-lean/blob/7cbf7a0c32fca1badf05931879a37c3f8a2b03b4/src/Formalism/SharpRestriction/GlobalTheoremInputs.lean#L45-L48}{LinkN}.

\emph{Source.} \href{https://github.com/CocoSgt/riemann-hypothesis-lean/blob/7cbf7a0c32fca1badf05931879a37c3f8a2b03b4/src/Formalism/SharpRestriction/GlobalTheoremInputs.lean#L40-L44}{LinkN\_\allowbreak at}.

\emph{Source.} \href{https://github.com/CocoSgt/riemann-hypothesis-lean/blob/7cbf7a0c32fca1badf05931879a37c3f8a2b03b4/src/Development/AutomorphicArithmetic/CharacterBoxes/DivinNOfHypV.lean#L117-L134}{divinN\_\allowbreak of\_\allowbreak hypV}.



\paragraph{2. GLOBAL\_\allowbreak RETAINED\_\allowbreak DEFECT — C — Original retained-row defect bound LinkIRR}

Rows.DelpNonneg and delp$\le$exp(-x\textasciicircum{}(ez/7)) uniformly for every late actual row datum. CURRENT VERSION0.47 QUALIFICATION: Finish the same-D/common-family adelic setting after ASG: normalized nonzero Bessel W and decay, measured-field block, original Siegel-Weil rescaling and measure/character ties, Bessel uniqueness/local ties, then the original AutTie/RowDictTie/LabTie. RowsFinDim is now proved universally by RFD. Supply the actual CountIdempotentsClass theorem or its three sufficient arithmetic cylinder counts, with one simultaneous family threshold; reuse the completed ARO structure. Completed new original producers: ARO1-ARO30, RFD1-RFD9, ASG1-ASG18. Any earlier narrative naming these completed producers as outstanding is historical; this qualification is the current scope authority. CURRENT VERSION0.48 QUALIFICATION: Finish the same-D/common-family adelic setting after ASG and AMD: normalized nonzero Bessel W and decay, original orbital/fiber and Siegel-Weil measure/character ties, Bessel uniqueness/local ties, and the full original transport theorem, then AutTie/RowDictTie/LabTie. The complete MExists measured-data block and unipotent covolume normalization are now supplied, and universal RowsFinDim remains available. Supply the actual CountIdempotentsClass theorem or its three sufficient arithmetic cylinder counts with one simultaneous family threshold. Completed current producers or conditional consumers: ARO1-ARO30, RFD1-RFD9, ASG1-ASG18, AMD1-AMD17. Earlier frontier descriptions are historical; this qualification is the current scope authority. Original whole-contract premises and quantifiers are retained. CURRENT VERSION0.49 QUALIFICATION: Prove the exact original BesselNVGap nonzero integral on its supplied discrete-series, compact-representation and character data. CoefWeightedLoc, WlocExists with its original nonvanishing antecedent, WFromLoc, the normalized global product and all-later-measured-block decay are complete. On the inherited Forms0 scope this leaves BesselNVGap for ConcreteSettingExists. Continue original orbital/fiber and Siegel-Weil measure/character ties, the required Bessel uniqueness/local ties and the full original transport theorem, then AutTie/RowDictTie/LabTie. Supply the actual CountIdempotentsClass theorem or its three sufficient arithmetic cylinder counts with one simultaneous family threshold. Completed current producers or conditional consumers: ARO1-ARO30, RFD1-RFD9, ASG1-ASG18, AMD1-AMD17, BWP1-BWP24 at their exact BesselNVGap/Forms0 scope. Earlier frontier descriptions are historical; this qualification is the current scope authority. Original whole-contract premises and quantifiers are retained.

\emph{Shared objects and choice order.} one FamP; one class $\mathcal A$ and its target-compatible subclass; one coefficient field per family datum; one x and i at each application; original selected rows, words and levels; thresholds chosen before the common family; AllRows A; LabelCountGoodInjective; AllJoinA; CountAllRows

\emph{Completed inputs (version 0.55).} SSC full Selmer-signature comparison, DFD common-family subresults, DNE full separate denominator consumer, WRS1-WRS10 full selected-class local gain, Inherited first DivDef3 and row-pair integrability

\emph{Remaining producer scope.} Complete the original F3 local coefficient bridges and F4 common-overfield readings, with valuation/norm/GoodCoef/iq and the entire coefficient field, and actual complex Hecke standard/no-Jordan plus arithmetic character data on the same form datum. Construct the actual base rows, host/frame lengths, nonzero chains and group ties, MultEq support, eigensystem/carrier/parts, all-embedding logLam, HD/d3b and NEb. Keep the stronger second DivRows3 scope distinct from the completed first DivDef3 and row-pair integrability. The full arithmetic entry and conditional common-frame construction do not by themselves install HOST\_\allowbreak NUMBER\_\allowbreak S, HOST\_\allowbreak FRAME\_\allowbreak LENGTH, TASK\_\allowbreak HOST\_\allowbreak VAN\_\allowbreak E\_\allowbreak INSTALL, TASK\_\allowbreak HOST\_\allowbreak XPP\_\allowbreak INSTALL or TASK\_\allowbreak COUNT\_\allowbreak V\_\allowbreak INSTALL. Continue the distinct retained-defect, period-height and normalized spectral producers, and the original compatible second-layer selection on one common arithmetic/geometric family. GLOBAL\_\allowbreak LOCAL\_\allowbreak GAIN remains fully closed at its original v0.49 scope. Complete local comparison theorems do not themselves construct the whole common family. Common design existence, NBUILT compatibility, original true-column records, shared fields/embeddings, Fits and eta agreement, simultaneous rows and eventual-family thresholds remain explicit. A stronger historical WPt/TSel construction may serve other producers but is not needed to reopen the proved selected-class local-gain estimate.

\emph{Source.} \href{https://github.com/CocoSgt/riemann-hypothesis-lean/blob/7cbf7a0c32fca1badf05931879a37c3f8a2b03b4/src/Formalism/SharpRestriction/GlobalTheoremInputs.lean#L55-L58}{LinkIRR}.

\emph{Source.} \href{https://github.com/CocoSgt/riemann-hypothesis-lean/blob/7cbf7a0c32fca1badf05931879a37c3f8a2b03b4/src/Development/AbelianVarietyTheory/Proofs/GlobalInterpolation/LinkIRROfJoinEx.lean#L17-L38}{linkIRR\_\allowbreak of\_\allowbreak join\_\allowbreak ex}.

\emph{Source.} \href{https://github.com/CocoSgt/riemann-hypothesis-lean/blob/7cbf7a0c32fca1badf05931879a37c3f8a2b03b4/src/Development/AbelianVarietyTheory/Proofs/GlobalInterpolation/LinkIRROfJoinEx.lean#L143-L166}{link\_\allowbreak irr}.



\paragraph{3. GLOBAL\_\allowbreak PERIOD\_\allowbreak HEIGHT — C — Full actual period and height comparison LinkDIVt}

Closed conditional TuSize -> DGamSize -> (LinkEta -> DIVinN -> LinkDIV), with every original LinkDIV\_\allowbreak at clause.

\emph{Shared objects and choice order.} one FamP; one class $\mathcal A$ and its target-compatible subclass; one coefficient field per family datum; one x and i at each application; original selected rows, words and levels; thresholds chosen before the common family; same RR target; same Bu/besS; same Nu and Bk

\emph{Source.} \href{https://github.com/CocoSgt/riemann-hypothesis-lean/blob/7cbf7a0c32fca1badf05931879a37c3f8a2b03b4/src/Formalism/SharpRestriction/GlobalTheoremInputs.lean#L77-L90}{LinkDIV\_\allowbreak at}.

\emph{Source.} \href{https://github.com/CocoSgt/riemann-hypothesis-lean/blob/7cbf7a0c32fca1badf05931879a37c3f8a2b03b4/src/Formalism/SharpRestriction/GlobalTheoremInputs.lean#L94-L96}{LinkDIVt}.

\emph{Source.} \href{https://github.com/CocoSgt/riemann-hypothesis-lean/blob/7cbf7a0c32fca1badf05931879a37c3f8a2b03b4/src/Development/AbelianVarietyTheory/Proofs/GlobalInterpolation/LinkBrest.lean#L7-L25}{GlobalTheory.Top.link\_\allowbreak brest}.

\emph{Source.} \href{https://github.com/CocoSgt/riemann-hypothesis-lean/blob/7cbf7a0c32fca1badf05931879a37c3f8a2b03b4/src/Development/AbelianVarietyTheory/Proofs/GlobalInterpolation/WireSel.lean#L11-L41}{GlobalTheory.Top.wire\_\allowbreak sel}.



\paragraph{4. GLOBAL\_\allowbreak SPECTRAL — C — Original normalized kernel/Gram estimates LinkS}

Closed TuSize -> LinkS with exact Gamma normalization, off-diagonal Hilbert-Schmidt bound, diagonal error and Dall/Dstar estimate, errors exp(-x\textasciicircum{}(epsilon/8)). CURRENT VERSION0.48 QUALIFICATION: Complete the original same-family spectral/Bessel and Siegel-Weil inputs with all measure, character, transport and common-object ties. The complete MExists measured block and prescribed-measure covolume normalization are supplied; these are not the full normalized spectral theorem. Completed current producers or conditional consumers: ASG1-ASG18, AMD1-AMD17, GCF1-GCF18 at their stated scopes. Earlier frontier descriptions are historical; this qualification is the current scope authority. Original whole-contract premises and quantifiers are retained. CURRENT VERSION0.49 QUALIFICATION: Complete the original same-family spectral/Bessel and Siegel-Weil inputs with every measure, character, transport and common-object tie. The normalized packet and decay are now constructed conditional on the exact BesselNVGap, and the analytic setting uses the original supplied Forms0 scope. These do not imply the full normalized spectral theorem. Completed current producers or conditional consumers: ASG1-ASG18, AMD1-AMD17, GCF1-GCF18 at their stated scopes, BWP1-BWP24 at their exact BesselNVGap/Forms0 scope, WRS1-WRS10 original local gain. Earlier frontier descriptions are historical; this qualification is the current scope authority. Original whole-contract premises and quantifiers are retained.

\emph{Shared objects and choice order.} one FamP; one class $\mathcal A$ and its target-compatible subclass; one coefficient field per family datum; one x and i at each application; original selected rows, words and levels; thresholds chosen before the common family; actual KD; same S.Cu spaces; GamOf; Dstar; Dall

\emph{Completed inputs (version 0.55).} SVW1–SVW18 (full conditional stabilizer volume and actual InVolumes join on unchanged GMR19 output), AOS1–AOS18 (full exact archimedean orbital scalar, gauge, compact volume and all-section integration clause), GMR1-GMR19 (full global/local measure suppliers and RescaleStep3 on the actual setting), BQS1-BQS23 (full QuotP/QuotM, StepStatement and BesselNVGap; setting at Forms0), BST1-BST11, BSE1-BSE9, BNV1-BNV23, BSP1-BSP10, BWP1-BWP24 at original Forms0 scope, AMD1-AMD24

\emph{Remaining producer scope.} SVW supplies full original VolStabFromSWS under exactly SWa1 and StabilizerFirstChernCompatibility, including every IsOmT, self-dual/fiber, character, Haar and fundamental-domain witness. It proves finite theta sums, actual rational Witt orbit transitivity, Borel stabilizer domains, both unimodularities, fixed-parameter density identification, nonnegative unfolding before finiteness, one common positive admissible smooth test, and finite stabilizer mass before scalar cancellation. It supplies the literal generic InVolumesAt and the actual bOf.InVolumes return on the GMR19 rescaled setting. That setting already contains both unipotent covolume-one equations; no new normalization is required for this join. AOS exact archimedean scalar and all-section integration remain available on the same bridge and infinity measures. The distinct RealBracket still needs reference unipotent covolume/discriminant calculation, broad JacGauge/PlaceSplit and ProdExact at every supplied factorization and section. The original Siegel-Weil arithmetic density premise itself, finite-place factors and remaining mass/index/column/form/counting fields are not proved by the conditional volume comparison. BQS full BesselNVGap is available; Forms0, all-family Witt/genus/level/intertwiner Transport/LevelMatch, Bessel uniqueness, AutTie/RowDictTie/LabTie, simultaneous arithmetic cylinder counts and the normalized GLOBAL\_\allowbreak SPECTRAL estimate remain.

\emph{Source.} \href{https://github.com/CocoSgt/riemann-hypothesis-lean/blob/7cbf7a0c32fca1badf05931879a37c3f8a2b03b4/src/Formalism/SharpRestriction/GlobalTheoremInputs.lean#L145-L149}{LinkS\_\allowbreak at}.

\emph{Source.} \href{https://github.com/CocoSgt/riemann-hypothesis-lean/blob/7cbf7a0c32fca1badf05931879a37c3f8a2b03b4/src/Formalism/SharpRestriction/GlobalTheoremInputs.lean#L152-L157}{LinkS}.

\emph{Source.} \href{https://github.com/CocoSgt/riemann-hypothesis-lean/blob/7cbf7a0c32fca1badf05931879a37c3f8a2b03b4/src/Development/AbelianVarietyTheory/Proofs/GlobalInterpolation/LinkAlt.lean#L16-L50}{GlobalTheory.Top.link\_\allowbreak alt}.



\paragraph{5. AUT\_\allowbreak UNRAM\_\allowbreak COEF — C — Actual unramified-coefficient automorphy}

For each actual Dat with coefficient ramification one and fixed iota, choose one marked g1 supplying FinA, AutB and LemW together. CURRENT VERSION0.48 QUALIFICATION: Construct the actual AutomorphyUnramifiedCoefficients package: for every original Dat with coefficient ramification one and fixed iota, choose one marked g1 that supplies FinA, AutB and LemW together. The original InpZ record, full RigII, CNR3, finite PropGK consumer and now the entire OneChar implication are supplied. Continue RH Cases I/V and the independent additional Dirichlet-GRH CaseIII/promodularity, SeedIV, tensor-factor and rank-four-chain tasks. Completed current producers or conditional consumers: OCH1-OCH30 (the full consumer implication), KIZ original InpZ, KTR1-KTR13, RNC16, PGK1-PGK6. Earlier frontier descriptions are historical; this qualification is the current scope authority. Original whole-contract premises and quantifiers are retained. CURRENT VERSION0.49 QUALIFICATION: Construct one common GdEnd for each original datum and fixed iota: the same marked g1 must supply FinEnd, AutEnd3 and AutEnd2. UCE now proves all ordinary local/curve/coefficient consumers and the exact implication to FinA/AutB/LemW, including arbitrary entry maps and every original outer-corner extension. The producer still needs the actual common abelian-CM GdFrameEx and LinkedEndE with both rank-two and rank-three runs, full ModUAlt/PfromAut and the two actual seeds, followed by the original local/automorphic descent. At this pin the alternative upstream route gives GdEndAlt; its additional local/descent fields must be matched to the actual original points, so no name-only GdEndAlt-to-GdEnd identification is allowed. Keep the supplied InpZ, RigII, CNR3, finite PropGK and whole OneChar implication. Continue RH Cases I/V and the independent additional Dirichlet-GRH CaseIII/promodularity, SeedIV, tensor-factor and rank-four-chain tasks. Completed current producers or conditional consumers: UCE1-UCE19 (complete GdEnd consumer), OCH1-OCH30, KIZ original InpZ, KTR1-KTR13, RNC16, PGK1-PGK6. Earlier frontier descriptions are historical; this qualification is the current scope authority. Original whole-contract premises and quantifiers are retained.

\emph{Completed inputs (version 0.55).} ALC1–ALC22 (full actual local existence, completion, flatness, residual types and representation/component charts), RTC1–RTC25 (upstream universal coefficient rings, full preconnected\_\allowbreak rank\_\allowbreak two and FinPtDown), ACP1-ACP19 (full FinPtOf and abstract propagation principles; rank-specific suppliers retained), AFE1-AFE21 (full Gen entries, conditional frame joins, NRep descent and FinPt-to-FinEnd), LFC1-LFC23, CAF1-CAF11, MFC1-MFC14, UCE1-UCE19 (full original GdEnd consumer), OCH1-OCH30, KIZ original InpZ, KTR1-KTR13, RNC16, PGK1-PGK6

\emph{Remaining producer scope.} ALC supplies full original cpl\_\allowbreak of\_\allowbreak univ, loc\_\allowbreak cpl, c3a, completion construction/uniqueness, LocFlat, idem, types, LocRho, unramified coefficient extension and loc\_\allowbreak ex. Every original closed point over the coefficient maximal ideal now has an actual local object with the required completion, prime contraction, flatness, representation and component-chart fields in both ranks two and three. The finite quotients are identified by the original universal property; Artin-Rees and inverse limits prove completion flatness on the original algebra map. Rank-three projectors are explicit polynomials in the marked diagonal and do not borrow the rank-two triangular argument. The local coefficient extension is constructed from a monic lift of a finite residue-field polynomial. All RTC upstream wpres/co\_\allowbreak B/gr\_\allowbreak A/cont\_\allowbreak u, preconnected\_\allowbreak rank\_\allowbreak two and FinPtDown are completed inputs; no FinA conclusion is used upstream. Remaining are sharp component chains, dimension bounds, nice-prime properties, global density and seed incidence, complete rank-two propagation, the rank-two-to-three corner and exceptional ideal, rank-three propagation and OffFin. Existence of the originally supplied universal spaces, actual automorphic seeds, patching/modularity, classicality and stronger Alt WitA/WitAQ outputs remain separate. Preserve the common AFE marked frame, original arithmetic P where used, rank-specific rings and ideals, all coefficient maps and the Host ell=3 case; do not identify the two ranks local coefficient extensions without a proved comparison.

\emph{Source.} \href{https://github.com/CocoSgt/riemann-hypothesis-lean/blob/7cbf7a0c32fca1badf05931879a37c3f8a2b03b4/src/Development/AutomorphicArithmetic/AutomorphicVarieties/AutomorphicCases/AutomorphyUnramifiedCoefficients.lean#L24-L27}{AutomorphyUnramifiedCoefficients}.

\emph{Source.} \href{https://github.com/CocoSgt/riemann-hypothesis-lean/blob/7cbf7a0c32fca1badf05931879a37c3f8a2b03b4/src/Development/AutomorphicArithmetic/AutomorphicVarieties/AutomorphicCases/AutomorphyUnramifiedCoefficients.lean#L12-L13}{FinA}.

\emph{Source.} \href{https://github.com/CocoSgt/riemann-hypothesis-lean/blob/7cbf7a0c32fca1badf05931879a37c3f8a2b03b4/src/Development/AutomorphicArithmetic/AutomorphicVarieties/AutomorphicCases/AutomorphyUnramifiedCoefficients.lean#L15-L17}{AutB}.

\emph{Source.} \href{https://github.com/CocoSgt/riemann-hypothesis-lean/blob/7cbf7a0c32fca1badf05931879a37c3f8a2b03b4/src/Development/AutomorphicArithmetic/AutomorphicVarieties/AutomorphicCases/AutomorphyUnramifiedCoefficients.lean#L19-L22}{LemW}.



\paragraph{6. TASK\_\allowbreak CASE3\_\allowbreak AUT\_\allowbreak JOIN — B — Complete the original Case-III automorphy/image join}

Install the actual Case-III ProModAlt, SeedIVAlt and shared PMExAlt with the inherited CohomologyHyp/Ferrero-Washington/HFit/SeedII facts into the original whole Case-III automorphy/image conclusion.

\emph{Source.} \href{https://github.com/CocoSgt/riemann-hypothesis-lean/blob/7cbf7a0c32fca1badf05931879a37c3f8a2b03b4/src/Development/AutomorphicArithmetic/AutomorphicVarieties/ModularityArgument/ThresholdArguments/DenseThresholdCases/PromodularityCaseThree.lean#L43-L44}{AutomorphyRankFourDense}.

\emph{Source.} \href{https://github.com/CocoSgt/riemann-hypothesis-lean/blob/7cbf7a0c32fca1badf05931879a37c3f8a2b03b4/src/Development/AutomorphicArithmetic/AutomorphicVarieties/ModularityArgument/ThresholdArguments/DenseThresholdCases/CuFinal.lean#L44-L44}{main4\_\allowbreak alt}.



\paragraph{7. TASK\_\allowbreak CASE3\_\allowbreak PROMOD\_\allowbreak JOIN — C — Reconstruct the actual Case-III chain-to-promodularity proof}

Use the original tensor-factor tower, ChainRankFour, ChainDesc4, seeded rank-four chain reduction, actual PMExAlt and SeedIVAlt to construct ProModAlt on every original Ctx3 star point at the same fixed iota, preserving all finite-coefficient, inertia and support maps.

\emph{Source.} \href{https://github.com/CocoSgt/riemann-hypothesis-lean/blob/7cbf7a0c32fca1badf05931879a37c3f8a2b03b4/src/Development/AutomorphicArithmetic/AutomorphicVarieties/ModularityArgument/ThresholdArguments/DenseThresholdCases/PromodularityCaseThreeOfLemmas.lean#L131-L153}{UpperTheory.V.Cu.promodularity\_\allowbreak case\_\allowbreak three\_\allowbreak of\_\allowbreak chain\_\allowbreak reduction}.

\emph{Source.} \href{https://github.com/CocoSgt/riemann-hypothesis-lean/blob/7cbf7a0c32fca1badf05931879a37c3f8a2b03b4/src/Development/AutomorphicArithmetic/AutomorphicVarieties/ModularityArgument/ThresholdArguments/DenseThresholdCases/PromodularityCaseThree.lean#L40-L42}{PromodularityCaseThreeDense}.



\paragraph{8. SEED\_\allowbreak IV\_\allowbreak ALT — C — Actual Case-III seed with prechosen residual datum}

Base Ctx3/Four tower, iota and Amb -> w,s before every later Dies tower -> compatible finite-coefficient seed with exact residual support/root entries.

\emph{Source.} \href{https://github.com/CocoSgt/riemann-hypothesis-lean/blob/7cbf7a0c32fca1badf05931879a37c3f8a2b03b4/src/Development/AutomorphicArithmetic/AutomorphicVarieties/ModularityArgument/ThresholdArguments/DenseThresholdCases/PromodularityCaseThree.lean#L19-L25}{SeedIVAlt}.



\paragraph{9. TENSOR\_\allowbreak FACTOR\_\allowbreak TOWER — B — Actual tensor-factor tower and characters}

Original finite-field, tame, factor and character constraints; source-level construction is sizeable and has not been reconstructed here.

\emph{Source.} \href{https://github.com/CocoSgt/riemann-hypothesis-lean/blob/7cbf7a0c32fca1badf05931879a37c3f8a2b03b4/src/Development/AutomorphicArithmetic/AutomorphicVarieties/ModularityArgument/ThresholdArguments/DenseThresholdCases/PromodularityCaseThreeOfLemmas.lean#L131-L153}{UpperTheory.V.Cu.promodularity\_\allowbreak case\_\allowbreak three\_\allowbreak of\_\allowbreak chain\_\allowbreak reduction}.



\paragraph{10. CHAIN\_\allowbreak RANK\_\allowbreak FOUR — C — Actual SomeIn-to-AllIn propagation}

Same source Space and PM ideal, residual point Good4, original bound 30(C+1)\textasciicircum{}2 <= d; geometry must preserve all identities.

\emph{Source.} \href{https://github.com/CocoSgt/riemann-hypothesis-lean/blob/7cbf7a0c32fca1badf05931879a37c3f8a2b03b4/src/Development/AutomorphicArithmetic/AutomorphicVarieties/ModularityArgument/ThresholdArguments/DenseThresholdCases/PromodularityCaseThree.lean#L26-L32}{ChainRankFour}.



\paragraph{11. CHAIN\_\allowbreak DESC4 — B — Original residual chain descent}

Finite residual descent retaining source Chain4 compatibility.

\emph{Source.} \href{https://github.com/CocoSgt/riemann-hypothesis-lean/blob/7cbf7a0c32fca1badf05931879a37c3f8a2b03b4/src/Development/AutomorphicArithmetic/AutomorphicVarieties/ModularityArgument/ThresholdArguments/DenseThresholdCases/PromodularityCaseThreeOfLemmas.lean#L131-L153}{UpperTheory.V.Cu.promodularity\_\allowbreak case\_\allowbreak three\_\allowbreak of\_\allowbreak chain\_\allowbreak reduction}.



\paragraph{12. RANK4\_\allowbreak CHAIN\_\allowbreak REDUCTION — C — Seeded reduction to the original residual chain}

Select T1 before every later Deep tower; retain one compatible seed and Chain4 for every original residual lift over Q.

\emph{Source.} \href{https://github.com/CocoSgt/riemann-hypothesis-lean/blob/7cbf7a0c32fca1badf05931879a37c3f8a2b03b4/src/Development/AutomorphicArithmetic/AutomorphicVarieties/ModularityArgument/ThresholdArguments/DenseThresholdCases/PromodularityCaseThree.lean#L33-L39}{RankFourChainReductionDense}.

\emph{Source.} \href{https://github.com/CocoSgt/riemann-hypothesis-lean/blob/7cbf7a0c32fca1badf05931879a37c3f8a2b03b4/src/Development/AutomorphicArithmetic/AutomorphicVarieties/ModularityArgument/ThresholdArguments/DenseThresholdCases/CuPChainIV.lean#L16-L82}{UpperTheory.V.Cu.chainIV\_\allowbreak alt}.



\paragraph{13. TASK\_\allowbreak COUNT\_\allowbreak V\_\allowbreak INSTALL — B — Complete the original CountV finite assembly}

From exactly the listed source host/cohomology inputs, complete countV\_\allowbreak of\_\allowbreak finiteness\_\allowbreak duality\_\allowbreak euler\_\allowbreak host, including original host/frame formation, normalized length, the CaseV guard and R1<=2\textasciicircum{}c2.

\emph{Completed inputs (version 0.55).} ALC actual local objects and charts; full dimensions/seeds/propagation retained, RTC full rank-two preconnected core and coefficient descent; actual propagation/counting inputs retained, ACP full finite-point and abstract propagation implications; actual rank-specific and counting inputs retained, AFE full Host entry/auxiliary frame/local witnesses including ell=3, LFC1-LFC23, CAF9-CAF11, MFC1-MFC14, Inherited CountV consumers

\emph{Remaining producer scope.} ALC supplies full original cpl\_\allowbreak of\_\allowbreak univ, loc\_\allowbreak cpl, c3a, completion construction/uniqueness, LocFlat, idem, types, LocRho, unramified coefficient extension and loc\_\allowbreak ex. Every original closed point over the coefficient maximal ideal now has an actual local object with the required completion, prime contraction, flatness, representation and component-chart fields in both ranks two and three. The finite quotients are identified by the original universal property; Artin-Rees and inverse limits prove completion flatness on the original algebra map. Rank-three projectors are explicit polynomials in the marked diagonal and do not borrow the rank-two triangular argument. The local coefficient extension is constructed from a monic lift of a finite residue-field polynomial. All RTC upstream wpres/co\_\allowbreak B/gr\_\allowbreak A/cont\_\allowbreak u, preconnected\_\allowbreak rank\_\allowbreak two and FinPtDown are completed inputs; no FinA conclusion is used upstream. Remaining are sharp component chains, dimension bounds, nice-prime properties, global density and seed incidence, complete rank-two propagation, the rank-two-to-three corner and exceptional ideal, rank-three propagation and OffFin. Existence of the originally supplied universal spaces, actual automorphic seeds, patching/modularity, classicality and stronger Alt WitA/WitAQ outputs remain separate. Preserve the common AFE marked frame, original arithmetic P where used, rank-specific rings and ideals, all coefficient maps and the Host ell=3 case; do not identify the two ranks local coefficient extensions without a proved comparison. Complete the original F3 local coefficient bridges and F4 common-overfield readings, with valuation/norm/GoodCoef/iq and the entire coefficient field, and actual complex Hecke standard/no-Jordan plus arithmetic character data on the same form datum. Construct the actual base rows, host/frame lengths, nonzero chains and group ties, MultEq support, eigensystem/carrier/parts, all-embedding logLam, HD/d3b and NEb. Keep the stronger second DivRows3 scope distinct from the completed first DivDef3 and row-pair integrability. The full arithmetic entry and conditional common-frame construction do not by themselves install HOST\_\allowbreak NUMBER\_\allowbreak S, HOST\_\allowbreak FRAME\_\allowbreak LENGTH, TASK\_\allowbreak HOST\_\allowbreak VAN\_\allowbreak E\_\allowbreak INSTALL, TASK\_\allowbreak HOST\_\allowbreak XPP\_\allowbreak INSTALL or TASK\_\allowbreak COUNT\_\allowbreak V\_\allowbreak INSTALL.

\emph{Source.} \href{https://github.com/CocoSgt/riemann-hypothesis-lean/blob/7cbf7a0c32fca1badf05931879a37c3f8a2b03b4/src/Development/AutomorphicArithmetic/AutomorphicVarieties/HeightAndVanishing/Counts.lean#L69-L69}{CountV}.

\emph{Source.} \href{https://github.com/CocoSgt/riemann-hypothesis-lean/blob/7cbf7a0c32fca1badf05931879a37c3f8a2b03b4/src/Development/AutomorphicArithmetic/AutomorphicVarieties/ModularityArgument/ThresholdArguments/DenseThresholdCases/CountVOf.lean#L33-L60}{UpperTheory.V.UnitaryThreeInertia.countV\_\allowbreak of\_\allowbreak finiteness\_\allowbreak duality\_\allowbreak euler\_\allowbreak host}.



\paragraph{14. THM\_\allowbreak S6\_\allowbreak S — C — Original unitary-three ordinary support and automorphy}

Original host/frame and fixed iota -> one marked g1 -> every universal R,u -> weight-zero finiteness and every Type-I BPrime\_\allowbreak s point.

\emph{Source.} \href{https://github.com/CocoSgt/riemann-hypothesis-lean/blob/7cbf7a0c32fca1badf05931879a37c3f8a2b03b4/src/Development/AutomorphicArithmetic/AutomorphicVarieties/ModularityArgument/ThresholdArguments/DenseThresholdCases/ThmS6CoreS.lean#L17-L26}{ThmS6\_\allowbreak s}.

\emph{Source.} \href{https://github.com/CocoSgt/riemann-hypothesis-lean/blob/7cbf7a0c32fca1badf05931879a37c3f8a2b03b4/src/Development/AutomorphicArithmetic/AutomorphicVarieties/ModularityArgument/ThresholdArguments/DenseThresholdCases/ThmS6CoreS.lean#L12-L16}{BPrime\_\allowbreak s}.



\paragraph{15. HOST\_\allowbreak FRAME\_\allowbreak LENGTH — B — Original Case-V host/frame and normalized-length installation}

Actual HostData, frame, local-condition equality, abstract count and normalized original Selmer length must all refer to the same objects.

\emph{Completed inputs (version 0.55).} AFE/LFC full entry and conditional common frame, CAF and MFC, DNE weight/host-degree dictionary at supplied tables

\emph{Remaining producer scope.} Complete the original F3 local coefficient bridges and F4 common-overfield readings, with valuation/norm/GoodCoef/iq and the entire coefficient field, and actual complex Hecke standard/no-Jordan plus arithmetic character data on the same form datum. Construct the actual base rows, host/frame lengths, nonzero chains and group ties, MultEq support, eigensystem/carrier/parts, all-embedding logLam, HD/d3b and NEb. Keep the stronger second DivRows3 scope distinct from the completed first DivDef3 and row-pair integrability. The full arithmetic entry and conditional common-frame construction do not by themselves install HOST\_\allowbreak NUMBER\_\allowbreak S, HOST\_\allowbreak FRAME\_\allowbreak LENGTH, TASK\_\allowbreak HOST\_\allowbreak VAN\_\allowbreak E\_\allowbreak INSTALL, TASK\_\allowbreak HOST\_\allowbreak XPP\_\allowbreak INSTALL or TASK\_\allowbreak COUNT\_\allowbreak V\_\allowbreak INSTALL.

\emph{Source.} \href{https://github.com/CocoSgt/riemann-hypothesis-lean/blob/7cbf7a0c32fca1badf05931879a37c3f8a2b03b4/src/Development/AutomorphicArithmetic/AutomorphicVarieties/ModularityArgument/ThresholdArguments/DenseThresholdCases/CountVOf.lean#L33-L60}{UpperTheory.V.UnitaryThreeInertia.countV\_\allowbreak of\_\allowbreak finiteness\_\allowbreak duality\_\allowbreak euler\_\allowbreak host}.



\paragraph{16. HOST\_\allowbreak NUMBER\_\allowbreak S — C — Uniform original-host rank-three catalogue}

One exponent before all closures, coefficient rings, levels, characters, frames, embeddings and marked points; one separating trace function controls all later finite pairwise nonisomorphic collections.

\emph{Completed inputs (version 0.55).} AFE full Host entry at original scope, LFC full conditional common frame, CAF and MFC, Inherited host-count consumers

\emph{Remaining producer scope.} Complete the original F3 local coefficient bridges and F4 common-overfield readings, with valuation/norm/GoodCoef/iq and the entire coefficient field, and actual complex Hecke standard/no-Jordan plus arithmetic character data on the same form datum. Construct the actual base rows, host/frame lengths, nonzero chains and group ties, MultEq support, eigensystem/carrier/parts, all-embedding logLam, HD/d3b and NEb. Keep the stronger second DivRows3 scope distinct from the completed first DivDef3 and row-pair integrability. The full arithmetic entry and conditional common-frame construction do not by themselves install HOST\_\allowbreak NUMBER\_\allowbreak S, HOST\_\allowbreak FRAME\_\allowbreak LENGTH, TASK\_\allowbreak HOST\_\allowbreak VAN\_\allowbreak E\_\allowbreak INSTALL, TASK\_\allowbreak HOST\_\allowbreak XPP\_\allowbreak INSTALL or TASK\_\allowbreak COUNT\_\allowbreak V\_\allowbreak INSTALL.

\emph{Source.} \href{https://github.com/CocoSgt/riemann-hypothesis-lean/blob/7cbf7a0c32fca1badf05931879a37c3f8a2b03b4/src/Development/AutomorphicArithmetic/AutomorphicVarieties/ModularityArgument/ThresholdArguments/DenseThresholdCases/HostNumberS.lean#L18-L28}{HostNumber\_\allowbreak s}.



\paragraph{17. TASK\_\allowbreak HOST\_\allowbreak VAN\_\allowbreak E\_\allowbreak INSTALL — B — Transport K/duality to literal HostVanE}

Use the already supplied LemmaK/local-global duality and finiteness on the actual CountV host/frame and exact original coefficient modules.

\emph{Source.} \href{https://github.com/CocoSgt/riemann-hypothesis-lean/blob/7cbf7a0c32fca1badf05931879a37c3f8a2b03b4/src/Development/AutomorphicArithmetic/AutomorphicVarieties/UnitaryThreeDeformations/HostNumber.lean#L9-L12}{HostVanE}.

\emph{Source.} \href{https://github.com/CocoSgt/riemann-hypothesis-lean/blob/7cbf7a0c32fca1badf05931879a37c3f8a2b03b4/src/Development/AutomorphicArithmetic/AutomorphicVarieties/LevelCohomology/SoHostVanE.lean#L11-L18}{UpperTheory.V.X.So.hostVanE}.



\paragraph{18. TASK\_\allowbreak HOST\_\allowbreak XPP\_\allowbreak INSTALL — B — Transport CorX to literal HostXpp}

Preserve Psi\textasciicircum{}-2 eta, the prime-to-ell quotient, actual local conditions and the original CountV host while installing BR7.

\emph{Source.} \href{https://github.com/CocoSgt/riemann-hypothesis-lean/blob/7cbf7a0c32fca1badf05931879a37c3f8a2b03b4/src/Development/AutomorphicArithmetic/AutomorphicVarieties/UnitaryThreeDeformations/HostNumber.lean#L14-L17}{HostXpp}.

\emph{Source.} \href{https://github.com/CocoSgt/riemann-hypothesis-lean/blob/7cbf7a0c32fca1badf05931879a37c3f8a2b03b4/src/Development/AutomorphicArithmetic/AutomorphicVarieties/LevelCohomology/Xpp.lean#L11-L84}{UpperTheory.V.X.xpp}.



\paragraph{19. DESIGN\_\allowbreak CHOICES\_\allowbreak COMMON — C — Construct the original common design family and bundle}

For the one target HostExW, construct the original compatible DesignChoices on Cdes=WithTw(WithSm Cprime), where Cprime=WithChN(A intersect target class), A is the one Built construction class, and cT(eps,y)=cTmin(y). Retain every original field, tuple, word, metric and common bundle duty. DSZ10 and RMP11 give size/rank properties of class members, not member existence. The original RecExP/RecAtG and common object/row route remains coupled. Construct ONE original common design family R and its recursive/geometric half with original Inputs/Pins/Holds/ExAt/InGeo and required CoreP. PKT now reconstructs full exists\_\allowbreak pkT, rowPk\_\allowbreak nonempty and exists\_\allowbreak moved with Copy/Coh on those original given inputs and a supplied fitting later stage: one catalogue, both ChT, complete LfT, actual integral lattices and all metric/measure/period constants are retained; one numerical record precedes every eta and V/S/N/G. RDR supplies one Side0 before all readers, ReadPk/pins and full ReadL with exact o.eta=q.eta. The initial actual geometry, Fits/division/central conditions, primitive rational tests and eta agreement must still be produced. Inherited CSR supplies finite columns, one ReadC before later twists, period/evaluation and scalar/error transports, RowSideDetermined and same-R Bundle/BTw/BWit/BunU consumers. RPI supplies exact RowPairInt for actual pairs and finite translates with original mu\_\allowbreak b/fd0; first literal DivDef3 at the same updated R,F,Y,Z,C,sEE follows from its retained original readings, Pins/Holds, stored nonzero IsItW chain, five real-group predicates and group tie. Remaining common fields include host/frame, MultEq nonzero-part support, that original nonzero chain/group tie, eigensystem/carrier/parts and all-embedding logLam, HD/d3b, NEb, actual rational tests and base rows. Retain every simultaneous family hypothesis and eventual threshold. Stronger second DivRows3 remains separate; universal quotient-height CuspRapidDecay and actual common member existence are not inferred from the comparison constructors. Version 0.41 supplies the finite phase-correct rational theta route: exact coefficient extension, split representatives, square-law/PhiIs, Pins from CoversAt, class-sum/point-identity transfer to OneMapOverL, actual-good-projector RatG and both-level LiftTestB. Its original higher covers, all-tuple class-sum uniqueness including phase rescaling, one PointIdentCore before O, base ClassSumsSpan and geometric LatFull remain explicit producers. MLat is available from the original LatFull input. This sharpens the rational-test work without claiming actual common geometry, Fits or eta agreement. Version0.46 proves finite generation of the original whole integral cusp lattice by actual integral-frame evaluation. It also proves the exact identity-cover, analytic unfolding, original-field coefficient descent, global clearing-scalar and both-level lattice-transport implications. Full spanning retains the supplied original IsNab connection and the five explicit principal-chart, level-system, raw-Gaussian, short-rule/dictionary and ramified-jet producers. The higher tuple covers, actual common class sums, PointIdentCore and ClassSumsSpan remain separate. CURRENT VERSION0.47 QUALIFICATION: Construct one original common geometric family and the higher CoversAt/all-tuple class sums and uniqueness with original phase rescalings, one PointIdentCore with common lambda,z, and base ClassSumsSpan. CNC now supplies the full original IsNab; PCL supplies the conditional trace-dual principal-lattice and quantitative normalization. Actual uniformization/HasPrinc and the five principal-chart, finite-level-system, Gaussian/reproduction, short-rule/simultaneous-dictionary and ramified-jet producer groups remain. Keep all original both-level, coefficient-field, Fits and eta agreements. Completed new original producers: ICL1-ICL21, CNC1-CNC16, PCL1-PCL12. Any earlier narrative naming these completed producers as outstanding is historical; this qualification is the current scope authority. CURRENT VERSION0.48 QUALIFICATION: Construct the original finite holomorphic FamilyUnif/N1c, normalized theta families and local lifting, then the common principal charts and finite-level theta systems. PUN now supplies the exact fixed-space models, all-family center lattices, A=1 exponent matching and individual primitive points with the full point-equality criterion; CNC supplies the full canonical connection and PCL the quantitative trace-dual normalization. Complete higher CoversAt and all-tuple class sums/uniqueness with the original phase rescalings, one PointIdentCore with common lambda,z, and base ClassSumsSpan. The Gaussian/reproduction, short-rule/simultaneous-dictionary and ramified-jet producers, and all original both-level, coefficient-field, Fits and eta agreements remain explicit. Completed current producers or conditional consumers: ICL1-ICL21, CNC1-CNC16, PCL1-PCL12, PUN1-PUN24, GCF1-GCF18 at their stated scopes. Earlier frontier descriptions are historical; this qualification is the current scope authority. Original whole-contract premises and quantifiers are retained. CURRENT VERSION0.49 QUALIFICATION: Construct the adjoining simultaneous principal framed points and test-ring dictionary. Full finite holomorphic FamilyUnif/N1c, normalized theta families and local lifts are now supplied by HUF. The next bounded target is allTest\_\allowbreak of\_\allowbreak secBaseChange: represent FrPt by the Lie-eigenframe bundles and theta-section frame line, use the full n=1 SecBaseChangeSt on every nonflat coefficient base change, and prove finite-generated subalgebra descent with the exact TRing/FromTest transport. Then supply the full FramesOK/CoverTri/FibTri and IsUnf/hasPrinc package, including adapted mod-eight symplectic basis, normalized theta multiplier/weight and arithmetic invariance. Continue higher CoversAt, all-tuple class sums and uniqueness with original phase rescalings, one PointIdentCore with common lambda,z, base ClassSumsSpan, Gaussian/reproduction, short-rule dictionaries and ramified jets, retaining both-level, coefficient-field, Fits and eta agreements. Common-family nonemptiness remains distinct from compatibility of an already supplied member. Completed current producers or conditional consumers: ICL1-ICL21, CNC1-CNC16, PCL1-PCL12, PUN1-PUN24, HUF1-HUF21 (full original FamilyUnif/N1c), GCF1-GCF18 and WRS1-WRS10 at their stated scopes. Earlier frontier descriptions are historical; this qualification is the current scope authority. Original whole-contract premises and quantifiers are retained.

\emph{Completed inputs (version 0.55).} SHV1–SHV21 (both original Serre conclusions, actual finite section model and full all-cover return), ASL1–ASL21 (actual section localization and full original arbitrary-base Perf producer), FCD1-FCD18 (full homogeneous-chart descent, graded presentation and Noetherian cube model), GSL1-GSL26 (full conditional section laws, universal basis and projective-model algebra), NFL1-NFL22, TDR1-TDR13, NFR1-NFR13, HUF1-HUF21, PUN1-PUN24, ICL1-ICL21, CNC1-CNC16, PCL1-PCL12, GCF/WRS at original supplied-member scope

\emph{Remaining producer scope.} SHV supplies both original Serre conjuncts: all-cover first-cohomology vanishing for every sufficiently high integer twist of an arbitrary line bundle on the original Noetherian very-ample scheme, and an actual homogeneous polynomial expression for every sufficiently high power-section. One finite homogeneous module and a finite syzygy list are fixed before the degree. Affine all-cover vanishing, two principal-cover exponents, refinement descent and the exact ordered sign return the chosen-cover calculation to the literal FirstCohomologyVanishes predicate. The original GSL\_\allowbreak SERRE input is now supplied; do not reopen it, the high-power polynomial theorem, or the ASL full arbitrary-base smooth-projective Perf producer. Remaining for GSL17 are PolBaseChangeSt, multiplication-pullback generated comparison, finite/flat/surjective and etale multiplication geometry, and invariant-function descent on every relevant open. Geometric integrality, principal rank one, torsion-indexed section decomposition, two-sided growth, and square/translation generators remain for the full P,V,R,T and GSL26/NFL joins. Original Moduli and HUF/PUN foundations remain. Later finite-level systems, Gaussian/reproduction, simultaneous dictionaries, ramified jets, all-tuple class sums, phase-sensitive uniqueness, common PointIdentCore and base ClassSumsSpan remain, retaining both levels, whole fields, Fits and eta agreement.



\paragraph{20. LOCAL\_\allowbreak NBUILT\_\allowbreak COMMON — C — Actual local specialization on the common selected family}

Construct, with the required closed-class property on the same Des(C,3), one eventual family of actual NBuilt certificates for the original Fm.D(x,i), retaining the certificate before forgetting to TSel. HLP now reconstructs full HeckeLift0, all-index commutation, HeckeStd0 and HeckeIntTw from the original MtInput and complex Hecke facts, in the full ambient GenAt range k<=nV. Every maximal ideal is an actual supported residual character class; literal MtSub is the O-span of projections of the original MtLat generators, is finite free, and localizes onto the exact NCov.part. Literal Psi has its exact range and off-space projection error, both original span assertions hold, and ActEx/ActHK are constructed on the same lattice under original semisimplicity. Original MjLat relation quotients, universal coefficient-map relation fields, selected tame generators, coefficient embeddings and all identity transports give the complete nonempty\_\allowbreak spec constructor on its full printed input list. A single compositum combines the four original coefficient requirements with big fixed beforehand. Whole ndSpec\_\allowbreak exists still needs actual common geometric/choice MtInput, complex HeckeStdPlain and the no-Jordan producer, original F3 bridges with valuation/GoodCoef/iq, original F4 common-overfield reading of the entire coefficient field, and exact arithmetic character/reciprocity/admissibility suppliers. Original AutSpec, Galois specialization and the remaining NBuilt fields on that same installation also remain. Do not replace the actual projected lattice or maps by abstract objects with the same dimensions.

\emph{Completed inputs (version 0.55).} SHV both Serre conclusions; ALC actual arithmetic charts; SVW same-setting volume join, ASL full actual perfect-complex model; RTC common-frame arithmetic inputs; AOS exact unchanged infinity normalization, FCD full algebraic model, ACP abstract arithmetic suppliers, GMR same-setting measure normalization, GSL/NFL full conditional section laws and uniformization, AFE/LFC full entry and conditional common frame, TDR/NFR and CAF/MFC, WRS selected built-class extraction at membership

\emph{Remaining producer scope.} Continue the distinct retained-defect, period-height and normalized spectral producers, and the original compatible second-layer selection on one common arithmetic/geometric family. GLOBAL\_\allowbreak LOCAL\_\allowbreak GAIN remains fully closed at its original v0.49 scope. Complete local comparison theorems do not themselves construct the whole common family. Common design existence, NBUILT compatibility, original true-column records, shared fields/embeddings, Fits and eta agreement, simultaneous rows and eventual-family thresholds remain explicit. A stronger historical WPt/TSel construction may serve other producers but is not needed to reopen the proved selected-class local-gain estimate. SHV supplies both original Serre conjuncts: all-cover first-cohomology vanishing for every sufficiently high integer twist of an arbitrary line bundle on the original Noetherian very-ample scheme, and an actual homogeneous polynomial expression for every sufficiently high power-section. One finite homogeneous module and a finite syzygy list are fixed before the degree. Affine all-cover vanishing, two principal-cover exponents, refinement descent and the exact ordered sign return the chosen-cover calculation to the literal FirstCohomologyVanishes predicate. The original GSL\_\allowbreak SERRE input is now supplied; do not reopen it, the high-power polynomial theorem, or the ASL full arbitrary-base smooth-projective Perf producer. Remaining for GSL17 are PolBaseChangeSt, multiplication-pullback generated comparison, finite/flat/surjective and etale multiplication geometry, and invariant-function descent on every relevant open. Geometric integrality, principal rank one, torsion-indexed section decomposition, two-sided growth, and square/translation generators remain for the full P,V,R,T and GSL26/NFL joins. Original Moduli and HUF/PUN foundations remain. Later finite-level systems, Gaussian/reproduction, simultaneous dictionaries, ramified jets, all-tuple class sums, phase-sensitive uniqueness, common PointIdentCore and base ClassSumsSpan remain, retaining both levels, whole fields, Fits and eta agreement.



\paragraph{21. SECOND\_\allowbreak LAYER\_\allowbreak COMPATIBLE\_\allowbreak SELECTION — C — Original second-layer selector and class duties}

For the same A, supply Duty on A and a Closed common original TSel family whose selected records are the original true-column rec7 records (or have the entire dependent tuple transported), with their real witness and original coefficient-field root data. Existence/compatibility remains arithmetic; merely abstractly isomorphic fields or independently selected columns do not suffice. CSR21–22 now construct the same-R Bundle/BTw/BWit/BunU outputs from the exact selected realization/row-side packet and NEb, with the original EvP/ParEv quantifier order. Actual common selection, the five stated realization packets and NEb remain upstream; none is replaced by an independently chosen compatible record. Version 0.38 closes the selected analytic extension and reconstructs the first updated DivDef3 consumer; it leaves the original common selection, geometric readings, rational tests, base rows, host/parts/NEb and exact retained DI premises upstream, with unchanged EvP/ParEv thresholds. Version 0.39 supplies the original given-data reader constructors and exact RowPairInt/first DivDef3 on original premises. It leaves one actual common geometric/comparison/test member, eta agreement, compatible selection, rational tests, base rows, host/parts/NEb and original EvP/ParEv thresholds upstream. Version 0.40 also supplies original comparison/RowPk/fitting-stage Moved algorithms and exact common data transport. The actual initial geometric member, fitting tests, eta agreement and compatible second-layer selector with original eventual thresholds remain upstream. Version 0.41 supplies the finite phase-correct rational theta route: exact coefficient extension, split representatives, square-law/PhiIs, Pins from CoversAt, class-sum/point-identity transfer to OneMapOverL, actual-good-projector RatG and both-level LiftTestB. Its original higher covers, all-tuple class-sum uniqueness including phase rescaling, one PointIdentCore before O, base ClassSumsSpan and geometric LatFull remain explicit producers. MLat is available from the original LatFull input. This sharpens the rational-test work without claiming actual common geometry, Fits or eta agreement. CURRENT VERSION0.49 QUALIFICATION: Continue the distinct global denominator, period/height, normalized spectral and second-layer compatible-selection producers, together with existence of the original common arithmetic/geometric family. GLOBAL\_\allowbreak LOCAL\_\allowbreak GAIN and its actual GroupCompatibilityFamily supplier are now complete on the unchanged selected built class; no additional WPt/TSel realization is needed for that estimate. Do not reopen local first-growth, mean, auxiliary admission, selected D3 tops/Adm2F, the single GIs2 stage or the exact comparison/kernel transport. A full RecT or TSel may remain useful for other tasks, but is not claimed by this local-gain closure. Completed current producers or conditional consumers: GCF1-GCF18 at their stated scopes, WRS1-WRS10 original selected-class compatibility/local gain. Earlier frontier descriptions are historical; this qualification is the current scope authority. Original whole-contract premises and quantifiers are retained.

\emph{Completed inputs (version 0.55).} IFD full invariant matrix descent; SHV/ALC/SVW full named nested producers, WIS actual inclusion and invariant completeness; AOS exact scalar; ASL/RTC full named nested producers, WLD full generic lattice constructors, GMR actual-setting normalization, FCD/ACP exact nested suppliers, WRS1-WRS10 original selected-class compatibility/local gain, DNE full denominator consumer, DFD/SSC canonical-data and signature comparisons, GSL/NFL and AFE/LFC conditional constructors, Inherited exact same-record realization/transport consumers

\emph{Remaining producer scope.} Continue the distinct retained-defect, period-height and normalized spectral producers, and the original compatible second-layer selection on one common arithmetic/geometric family. GLOBAL\_\allowbreak LOCAL\_\allowbreak GAIN remains fully closed at its original v0.49 scope. Complete local comparison theorems do not themselves construct the whole common family. Common design existence, NBUILT compatibility, original true-column records, shared fields/embeddings, Fits and eta agreement, simultaneous rows and eventual-family thresholds remain explicit. A stronger historical WPt/TSel construction may serve other producers but is not needed to reopen the proved selected-class local-gain estimate.



\paragraph{22. BASE\_\allowbreak THETA\_\allowbreak READINGS — C — Actual base reading on the common bundle}

For every family in the original class, eventually on the SAME selected Bundle and for each i<m, supply the actual R,F,Read,ReadL, weak models/Weil data and phase-one pins. Identify their LiftTie.theta, mtRead, mjRead and ThDef3 basis with that Bundle.P.A3, P.Zs, P.S0 and its stored maps. Exact field, embeddings, metrics, measures, row family and weak choices are retained. This remains one part of the coupled BTH12 producer. PKT now supplies the complete original comparison packet, RowPk and fitting-stage Moved with Copy/Coh on original given geometry, Inputs/Pins/Holds/ExAt/InGeo and Fits. Its ChT/LfT/Reps, actual lattice carriers, original Haar and norm constants, period maps, representatives and all-place RrT are constructed together; its one numerical record precedes every eta/V/S/N/G. RDR gives the original level-zero side, moved-Read and ReadPk/pin constructors on those objects, with one Side0 before every reader and all twelve Read ties. Full ReadL retains the exact o.eta(j,v)=q.eta(j,v) guard for every 1<=j<=m and place v; arbitrary Over does not imply it. Inherited CSR fixes one ReadC before all TwArgs, transports both metrics and all thirteen reading maps, and gives actual period/evaluation with kappa=(c*nP)\textasciicircum{}(-1/2). RPI supplies exact row-pair integrability, so first updated DivDef3 follows from the retained original Pins/Holds, nonzero chain, five real-group predicates and group tie. Initial theta extension remains complete. Remaining work is actual initial common geometry and its original predicates, Fits/primitive rational tests, eta agreement and outer family/eventual/same-Bundle duties, reFace compatibility, MultEq support, original chain/group data, actual base rows and any stronger DivRows3 theta square. No new compatible reading, family threshold or Haar normalization is substituted. Version 0.41 supplies the finite phase-correct rational theta route: exact coefficient extension, split representatives, square-law/PhiIs, Pins from CoversAt, class-sum/point-identity transfer to OneMapOverL, actual-good-projector RatG and both-level LiftTestB. Its original higher covers, all-tuple class-sum uniqueness including phase rescaling, one PointIdentCore before O, base ClassSumsSpan and geometric LatFull remain explicit producers. MLat is available from the original LatFull input. This sharpens the rational-test work without claiming actual common geometry, Fits or eta agreement.

\emph{Completed inputs (version 0.55).} SHV1–SHV21 (both original Serre conclusions, actual finite section model and full all-cover return), ASL1–ASL21 (actual section localization and full original arbitrary-base Perf producer), FCD1-FCD18 (full homogeneous-chart descent, graded presentation and Noetherian cube model), GSL1-GSL26, NFL1-NFL22, TDR1-TDR13, NFR1-NFR13, HUF1-HUF21, PUN1-PUN24, CNC1-CNC16, Inherited actual-reading consumers at original inputs

\emph{Remaining producer scope.} SHV supplies both original Serre conjuncts: all-cover first-cohomology vanishing for every sufficiently high integer twist of an arbitrary line bundle on the original Noetherian very-ample scheme, and an actual homogeneous polynomial expression for every sufficiently high power-section. One finite homogeneous module and a finite syzygy list are fixed before the degree. Affine all-cover vanishing, two principal-cover exponents, refinement descent and the exact ordered sign return the chosen-cover calculation to the literal FirstCohomologyVanishes predicate. The original GSL\_\allowbreak SERRE input is now supplied; do not reopen it, the high-power polynomial theorem, or the ASL full arbitrary-base smooth-projective Perf producer. Remaining for GSL17 are PolBaseChangeSt, multiplication-pullback generated comparison, finite/flat/surjective and etale multiplication geometry, and invariant-function descent on every relevant open. Geometric integrality, principal rank one, torsion-indexed section decomposition, two-sided growth, and square/translation generators remain for the full P,V,R,T and GSL26/NFL joins. Original Moduli and HUF/PUN foundations remain. Later finite-level systems, Gaussian/reproduction, simultaneous dictionaries, ramified jets, all-tuple class sums, phase-sensitive uniqueness, common PointIdentCore and base ClassSumsSpan remain, retaining both levels, whole fields, Fits and eta agreement. Complete the original F3 local coefficient bridges and F4 common-overfield readings, with valuation/norm/GoodCoef/iq and the entire coefficient field, and actual complex Hecke standard/no-Jordan plus arithmetic character data on the same form datum. Construct the actual base rows, host/frame lengths, nonzero chains and group ties, MultEq support, eigensystem/carrier/parts, all-embedding logLam, HD/d3b and NEb. Keep the stronger second DivRows3 scope distinct from the completed first DivDef3 and row-pair integrability. The full arithmetic entry and conditional common-frame construction do not by themselves install HOST\_\allowbreak NUMBER\_\allowbreak S, HOST\_\allowbreak FRAME\_\allowbreak LENGTH, TASK\_\allowbreak HOST\_\allowbreak VAN\_\allowbreak E\_\allowbreak INSTALL, TASK\_\allowbreak HOST\_\allowbreak XPP\_\allowbreak INSTALL or TASK\_\allowbreak COUNT\_\allowbreak V\_\allowbreak INSTALL.



\paragraph{23. BASE\_\allowbreak THETA\_\allowbreak RATIONAL\_\allowbreak TESTS — C — Simultaneous original base rational tests}

On the SAME original Bundle/reading/theta map, for every i<m supply one finite separating complex-linear test family whose values lie in the fixed embedded coefficient field on theta of EVERY original rational source vector and on EVERY stored target-coordinate basis vector. BTH8/BTH11 use this same LiftTestB for rational adjoint image and matrix coordinates. RTG now reconstructs prescribed-field coefficient extension, original split class representatives/RepsOK, the self-dual split-lattice square law and finite residue PhiIs; every Pins field follows from the actual CoversAt. With the base coefficient record literally fixed, exact class-sum existence/uniqueness including phase-rescaled O, PointIdentCore with common lambda,z before every O, and base ClassSumsSpan imply all-tuple finite lift span and OneMapOverL. The phase is exactly cj(twAt embLL A3.A.v0)(R.phs i). RatG uses the actual good projector eOp(y), and full LiftTestB follows at BOTH original levels. These finite implication proofs are complete. Remaining original producers are the prescribed higher covers, those actual class sums/uniqueness, PointIdentCore, base ClassSumsSpan and the base/common geometric datum; exact MLat is supplied only from original LatFull, whose geometric production remains separate. Finite dimension or an unrelated rational theta model cannot replace these same-object identities. Original outer eventual-family/Bundle quantifiers remain. Version0.46 proves finite generation of the original whole integral cusp lattice by actual integral-frame evaluation. It also proves the exact identity-cover, analytic unfolding, original-field coefficient descent, global clearing-scalar and both-level lattice-transport implications. Full spanning retains the supplied original IsNab connection and the five explicit principal-chart, level-system, raw-Gaussian, short-rule/dictionary and ramified-jet producers. The higher tuple covers, actual common class sums, PointIdentCore and ClassSumsSpan remain separate. CURRENT VERSION0.47 QUALIFICATION: Construct one original common geometric family and the higher CoversAt/all-tuple class sums and uniqueness with original phase rescalings, one PointIdentCore with common lambda,z, and base ClassSumsSpan. CNC now supplies the full original IsNab; PCL supplies the conditional trace-dual principal-lattice and quantitative normalization. Actual uniformization/HasPrinc and the five principal-chart, finite-level-system, Gaussian/reproduction, short-rule/simultaneous-dictionary and ramified-jet producer groups remain. Keep all original both-level, coefficient-field, Fits and eta agreements. Completed new original producers: ICL1-ICL21, CNC1-CNC16, PCL1-PCL12. Any earlier narrative naming these completed producers as outstanding is historical; this qualification is the current scope authority. CURRENT VERSION0.48 QUALIFICATION: Construct the original finite holomorphic FamilyUnif/N1c, normalized theta families and local lifting, then the common principal charts and finite-level theta systems. PUN now supplies the exact fixed-space models, all-family center lattices, A=1 exponent matching and individual primitive points with the full point-equality criterion; CNC supplies the full canonical connection and PCL the quantitative trace-dual normalization. Complete higher CoversAt and all-tuple class sums/uniqueness with the original phase rescalings, one PointIdentCore with common lambda,z, and base ClassSumsSpan. The Gaussian/reproduction, short-rule/simultaneous-dictionary and ramified-jet producers, and all original both-level, coefficient-field, Fits and eta agreements remain explicit. Completed current producers or conditional consumers: ICL1-ICL21, CNC1-CNC16, PCL1-PCL12, PUN1-PUN24. Earlier frontier descriptions are historical; this qualification is the current scope authority. Original whole-contract premises and quantifiers are retained. CURRENT VERSION0.49 QUALIFICATION: Construct the adjoining simultaneous principal framed points and test-ring dictionary. Full finite holomorphic FamilyUnif/N1c, normalized theta families and local lifts are now supplied by HUF. The next bounded target is allTest\_\allowbreak of\_\allowbreak secBaseChange: represent FrPt by the Lie-eigenframe bundles and theta-section frame line, use the full n=1 SecBaseChangeSt on every nonflat coefficient base change, and prove finite-generated subalgebra descent with the exact TRing/FromTest transport. Then supply the full FramesOK/CoverTri/FibTri and IsUnf/hasPrinc package, including adapted mod-eight symplectic basis, normalized theta multiplier/weight and arithmetic invariance. Continue higher CoversAt, all-tuple class sums and uniqueness with original phase rescalings, one PointIdentCore with common lambda,z, base ClassSumsSpan, Gaussian/reproduction, short-rule dictionaries and ramified jets, retaining both-level, coefficient-field, Fits and eta agreements. Completed current producers or conditional consumers: ICL1-ICL21, CNC1-CNC16, PCL1-PCL12, PUN1-PUN24, HUF1-HUF21 (full original FamilyUnif/N1c and theta family). Earlier frontier descriptions are historical; this qualification is the current scope authority. Original whole-contract premises and quantifiers are retained.

\emph{Completed inputs (version 0.55).} SHV1–SHV21 (both original Serre conclusions, actual finite section model and full all-cover return), ASL1–ASL21 (actual section localization and full original arbitrary-base Perf producer), FCD1-FCD18 (full homogeneous-chart descent, graded presentation and Noetherian cube model), GSL1-GSL26, NFL1-NFL22, TDR1-TDR13, NFR1-NFR13, HUF1-HUF21, PUN1-PUN24, ICL1-ICL21, CNC1-CNC16, PCL1-PCL12, Inherited both-level RatG/LiftTestB consumers

\emph{Remaining producer scope.} SHV supplies both original Serre conjuncts: all-cover first-cohomology vanishing for every sufficiently high integer twist of an arbitrary line bundle on the original Noetherian very-ample scheme, and an actual homogeneous polynomial expression for every sufficiently high power-section. One finite homogeneous module and a finite syzygy list are fixed before the degree. Affine all-cover vanishing, two principal-cover exponents, refinement descent and the exact ordered sign return the chosen-cover calculation to the literal FirstCohomologyVanishes predicate. The original GSL\_\allowbreak SERRE input is now supplied; do not reopen it, the high-power polynomial theorem, or the ASL full arbitrary-base smooth-projective Perf producer. Remaining for GSL17 are PolBaseChangeSt, multiplication-pullback generated comparison, finite/flat/surjective and etale multiplication geometry, and invariant-function descent on every relevant open. Geometric integrality, principal rank one, torsion-indexed section decomposition, two-sided growth, and square/translation generators remain for the full P,V,R,T and GSL26/NFL joins. Original Moduli and HUF/PUN foundations remain. Later finite-level systems, Gaussian/reproduction, simultaneous dictionaries, ramified jets, all-tuple class sums, phase-sensitive uniqueness, common PointIdentCore and base ClassSumsSpan remain, retaining both levels, whole fields, Fits and eta agreement.



\paragraph{24. SLR\_\allowbreak BASE\_\allowbreak ROWS\_\allowbreak PRODUCER — C — Actual base row identities on that same bundle}

For the SAME selected Bundle, common original fields and readings used by the other three BTH12 inputs, supply actual RowsHold0\_\allowbreak 3 uniformly eventually for every family. Its original primitive row realization, spectrum/type and basis identifications must be retained. BTH supplies a theta consumer only; it does not construct the base row identity. No separate witness may be chosen merely to satisfy this row requirement.



\clearpage
\section{Serre vanishing on every cover and polynomial high-power sections}
\label{sec:serre-h1-v055}

\status{This chapter proves the full original
\texttt{exists\_\allowbreak h1Zero\_\allowbreak tw\_\allowbreak of\_\allowbreak very\allowbreak Ample\allowbreak By}:
over a Noetherian ring, every actual line bundle has vanishing first
cohomology after every sufficiently large integer twist by the supplied
very-ample bundle.  Vanishing means the original coordinate-cocycle
statement for every open cover.  No smoothness assumption is added.
The same argument also proves the original homogeneous-polynomial
description of every sufficiently high power-section, and hence both
conjuncts of \texttt{serre\_\allowbreak h1Zero\_\allowbreak and\_\allowbreak sections\_\allowbreak polynomial}.
The separate multiplication, polarized base-change, geometric rank and
square-law suppliers remain explicit.}

There are two different finite constructions.  First choose one finite
homogeneous submodule of the actual graded section module.  Its
localizations already give the actual sections on every nonempty
intersection in the prescribed cover, in every integer degree.  Next,
over a Noetherian base, take finitely many successive kernels of graded
free surjections.  Their free-module bounds have one common maximum.
These choices occur before the twist degree, before a cocycle and before
an arbitrary cover used to test vanishing.

The passage from the prescribed affine cover to every other cover is
proved separately.  An affine open has the required all-cover vanishing
by explicit denominator clearing.  Local primitives on the prescribed
affine opens differ by a cocycle on that cover.  Solving this latter
cocycle makes the local primitives agree, so they descend to the original
test cover.  Exactness on one cover is therefore connected to the literal
source predicate by an actual gluing argument.

\paragraph{Notation and inherited objects.}
Let $R$ be a commutative ring, $Y$ a scheme and
$f:Y\to\operatorname{Spec}R$.  For a line bundle $M$ and an open $V$,
write $\operatorname{Sec}(M,V)$ for its actual compatible coordinate
families.  Its $R$-action uses the original map
$c_f(V):R\to\mathcal O(V)$.  If $H$ is another line bundle, put
$M(t)=M\otimes H^t$ for every $t\in\mathbb Z$, with the actual
transition convention of ASL2.  Negative powers are dual powers.

Fix the original data $\operatorname{VeryAmpleBy}(f,H,s)$:
$s_0,\ldots,s_m$ are sections of $H$; their nonzero loci
$D_i=D(s_i)$ cover $Y$; every $D_i$ is affine; and the coordinate
ratios $r_{ij}=s_j/s_i$ give a surjection
\[
 \pi_i:R[X_0,\ldots,X_m]\longrightarrow\mathcal O(D_i),
 \qquad X_j\longmapsto r_{ij}.
                                                        \tag{SHV.1}
\]
The ratios and this evaluation use the original structure map $f$.
ASL9 supplies quasi-compactness, quasi-separatedness and affine nonempty
intersections of this cover.  ASL8, ASL12, ASL14--ASL16 supply the actual
graded module
\[
 G=\bigoplus_{t\in\mathbb Z}\operatorname{Sec}(M(t),Y),
 \qquad X_i a=s_i a,
                                                        \tag{SHV.2}
\]
its localization-to-section maps, its dehomogenizations and the finite
projective affine section modules.  FCD10 and FCD13 supply the proved
finite saturation and quotient arguments.  None of these inherited
results requires the smoothness of $Y$.

For the algebraic part let $\iota$ be any finite linearly ordered set,
$\nu=|\iota|$, $S=R[X_i:i\in\iota]$, and
$Q=\bigoplus_{t\in\mathbb Z}Q_t$ a graded $S$-module with variables
of degree one.  For $I\subseteq\iota$ put $X_I=\prod_{i\in I}X_i$
and
\[
 Q_{I,t}=(Q[X_I^{-1}])_t
 =\left\{\frac{a}{X_I^k}:\
       k\ge0,\ a\in Q_{t+k|I|}\right\}.
                                                        \tag{SHV.3}
\]
This is the source's literal \texttt{hom\allowbreak Loc}; $Q_{\varnothing,t}=Q_t$.
Its augmented localization complex is
\[
 A^p(Q,t)=\bigoplus_{\substack{I\subseteq\iota\\|I|=p}}Q_{I,t},
 \qquad
 (d x)_J=\sum_{j\in J}(-1)^{|\{i\in J:i<j\}|}
                   x_{J\setminus\{j\}}\big|_J .
                                                        \tag{SHV.4}
\]
There are finitely many subsets, so the displayed direct sums are also
products.  The complex is zero outside degrees $0,\ldots,\nu$.
Insertion of two different indices in opposite orders gives opposite
signs, so $d^2=0$.  The ordinary complex is
$C^q(Q,t)=A^{q+1}(Q,t)$, $q\ge0$, with the same differential.
Thus ordinary degree-one exactness is augmented degree-two exactness.

\subsection{One degree bound for the whole localization complex}

% BEGIN SHV1
\paragraph{SHV1: graded localization preserves the actual short exact sequence.}
\label{claim:SHV1}
If $0\to Q'\xrightarrow{u}Q\xrightarrow{v}Q''\to0$ is a short
exact sequence of degree-preserving $S$-linear maps, then for every
$I$ and every integer $t$,
\[
 0\longrightarrow Q'_{I,t}\xrightarrow{u_I}Q_{I,t}
   \xrightarrow{v_I}Q''_{I,t}\longrightarrow0
                                                        \tag{SHV.5}
\]
is exact.  The maps commute with every localization restriction, and
therefore give a short exact sequence of the complexes in (SHV.4).
The vertex assertion holds for an arbitrary variable set and a finite
$I$; finiteness of $\iota$ is only needed for the finite sums in (SHV.4).

\begin{proof}
A homogeneous element of degree $a$ in $Q''$ has a homogeneous lift
of the same degree in $Q$: take any lift and project it to degree $a$.
The other homogeneous components map to zero, since the grading is a
direct sum.  This proves surjectivity in (SHV.5) by lifting the
numerator of a fraction.

Suppose $b/X_I^k$ maps to zero.  The equality of this fraction with
zero means that $X_I^\ell v(b)=0$ for some $\ell\ge0$.
Thus $X_I^\ell b=u(a)$ for a homogeneous $a$ of degree
$t+(k+\ell)|I|$; the homogeneous lift follows by the same projection
argument.  The fraction $a/X_I^{k+\ell}$ maps to $b/X_I^k$.
If a fraction from $Q'$ maps to zero, injectivity of $u$ applied
after multiplication by the killing power of $X_I$ proves that the
original fraction was zero.  These arguments also apply to
$I=\varnothing$, where $X_I=1$.

For $I\subseteq J$ the restriction is
$a/X_I^k\mapsto X_{J\setminus I}^ka/X_J^k$.
An $S$-linear map commutes with that formula.  Taking the finite
sums over all subsets of each cardinality proves termwise exactness
and compatibility with the signed differentials.  No tensoring by
an assumed flat coefficient algebra has been used.
\end{proof}
% END SHV1

% BEGIN SHV2
\paragraph{SHV2: the free polynomial complex and its positive-degree contraction.}
\label{claim:SHV2}
For every commutative $R$, the cohomology of $A^\bullet(S,t)$
is a finite free $R$-module in each degree.  It is zero in every
degree when $t\ge1$.

\begin{proof}
We use the Laurent-monomial calculation underlying FCD15 and spell
out its degree bound.  In the vertex $I$ a basis is
\[
 X^\alpha,\qquad
 \alpha\in\mathbb Z^\iota,\quad
 \sum_i\alpha_i=t,\quad
 N(\alpha):=\{i:\alpha_i<0\}\subseteq I .
                                                        \tag{SHV.6}
\]
To obtain such a monomial as a fraction, choose $k$ large enough that
$\alpha_i+k\ge0$ for $i\in I$.  Its numerator has those exponents
on $I$ and exponents $\alpha_i$ outside $I$, and has degree
$t+k|I|$.  Conversely every homogeneous numerator is a finite sum
of these monomials.  Their linear independence follows by clearing
one common denominator and using the usual polynomial monomial
basis.  Multiplication by $X_I$ is injective on a polynomial ring
over any $R$, because it only shifts its monomial coefficients.
This argument does not assume that $R$ is a domain.

The restrictions carry a monomial to the identical Laurent monomial.
Consequently the whole complex is a direct sum, over $\alpha$, of
complexes with basis $e_I$ for $I\supseteq N(\alpha)$ in degree $|I|$.
If $N(\alpha)\ne\iota$, choose $j_0\notin N(\alpha)$ and set
\[
 h(e_I)=
 \begin{cases}
 (-1)^{|\{i\in I:i<j_0\}|}e_{I\setminus\{j_0\}},&j_0\in I,\\
 0,&j_0\notin I.
 \end{cases}                                           \tag{SHV.7}
\]
The smaller subset still contains $N(\alpha)$.  In $dh+hd$ the
term inserting and deleting $j_0$ has coefficient one.  For
another index $j$, the two possible orders of insertion and deletion
have signs differing by one in their parity: one of $j,j_0$ lies
before the other in the order.  These two terms cancel.  Hence
$dh+hd=1$.  This includes the augmented degree-zero term.

Only vectors with $N(\alpha)=\iota$ can contribute cohomology;
each gives a single copy of $R$ in degree $\nu$ and no other term.
For $\nu>0$, write $\alpha_i=-1-\beta_i$.  The possibilities are
the nonnegative integer solutions of
$\sum_i\beta_i=-t-\nu$, a finite set for each fixed $t$.
If $\nu=0$, there is one empty exponent vector when $t=0$ and
none otherwise.  Thus the surviving direct sum is finite free.
Direct sums cause no completion issue: every cochain has finite
monomial support, and the contraction preserves that support.

For $t\ge1$, an exponent vector with every entry negative cannot
have sum $t$.  When $\iota$ is empty its sum is zero, also
impossible.  Thus all summands contract and the complex is exact
in every degree.  This proves both assertions, including the
empty variable set.
\end{proof}
% END SHV2

% BEGIN SHV3
\paragraph{SHV3: an explicit common bound for all shifts of a finite free module.}
\label{claim:SHV3}
Let $F=\bigoplus_{\ell=1}^r S e_\ell$ with $\deg e_\ell=d_\ell$
in $\mathbb Z$.  Every $H^p(A^\bullet(F,t))$ is finite over $R$.
The single bound
\[
                    b(F)=1+\sum_{\ell=1}^r|d_\ell|
                                                        \tag{SHV.8}
\]
makes $A^\bullet(F,t)$ exact in every degree whenever $t\ge b(F)$.

\begin{proof}
Localizing and taking a homogeneous piece commute with this finite
direct sum.  Its $\ell$th summand is the complex
$A^\bullet(S,t-d_\ell)$, since division by $X_I^k$ leaves the
total degree of $X^\alpha e_\ell$ equal to
$\sum_i\alpha_i+d_\ell$.  The restrictions and signs are identical.
SHV2 gives finiteness for each summand.

For every $\ell$,
$d_\ell\le|d_\ell|\le\sum_j|d_j|$.  Therefore
$t\ge b(F)$ implies $t-d_\ell\ge1$, and SHV2 contracts all the
summands.  This is one bound for every cohomological degree, not a
bound chosen after that degree.  If $r=0$, the module and complex
are zero; the displayed bound is still valid.
\end{proof}
% END SHV3

% BEGIN SHV4
\paragraph{SHV4: finitely many graded kernels can be chosen before the degree.}
\label{claim:SHV4}
Suppose $R$ is Noetherian and $Q$ is finite over $S$.
For every $a\ge1$ there are finite graded free modules
$F_0,\ldots,F_{a-1}$ and finite graded modules
$Q_0=Q,Q_1,\ldots,Q_a$, with short exact sequences
\[
 0\longrightarrow Q_{j+1}\longrightarrow F_j
       \longrightarrow Q_j\longrightarrow0
       \qquad(0\le j<a).                               \tag{SHV.9}
\]
All these objects can be fixed before any twist degree or
cohomological degree is specified.

\begin{proof}
Replace a finite generating family of $Q$ by the finitely many
homogeneous components of all its members.  They still generate,
because their span contains the original generators; they belong to
$Q$, so their span cannot be larger than $Q$.  Declare a free basis
element to have the degree of the corresponding generator.  Mapping
that basis element to the generator gives a degree-preserving
surjection $F_0\to Q$.

Its kernel is graded.  Indeed, if the image of a sum of homogeneous
components is zero, each image is zero, since the target grading
is a direct sum.  Hilbert's basis theorem makes $S$ Noetherian,
and a submodule of the finite free module $F_0$ is finite.
Its homogeneous components again give a finite homogeneous
generating family.  Repeat this construction $a$ times.
This is a finite recursion depending only on the original graded
module and the chosen finite families.  It does not assert that the
last kernel is zero or projective, and it does not assume finite
global dimension for $R$.
\end{proof}
% END SHV4

% BEGIN SHV5
\paragraph{SHV5: the exactness transfer used in dimension shifting.}
\label{claim:SHV5}
Let $0\to B^\bullet\xrightarrow{u}E^\bullet
\xrightarrow{v}D^\bullet\to0$ be a termwise short exact sequence
of cochain complexes.  If $E^\bullet$ is exact in degree $p$
and $B^\bullet$ is exact in degree $p+1$, then $D^\bullet$
is exact in degree $p$.

\begin{proof}
Let $z\in D^p$ be a cocycle.  Lift it to $e\in E^p$.
Since $v(de)=d z=0$, write $de=u(b)$ with $b\in B^{p+1}$.
The equality $u(db)=d^2e=0$ and injectivity of $u$ show that
$b$ is a cocycle.  Exactness at $p+1$ gives $b=db_0$ for
$b_0\in B^p$.  Thus $e-u(b_0)$ is a cocycle in $E^p$.
It equals $de_0$ by exactness there, and applying $v$ gives
$z=d(v(e_0))$.  This is the required coboundary.

For nonnegative complexes at $p=0$, put the degree $-1$ terms
equal to zero.  Exactness of $E$ at zero says directly that
$e-u(b_0)=0$, so $z=0$.  Thus the same argument includes the
bottom degree.  Every boundary is a cocycle because $d^2=0$,
completing the exactness assertion.
\end{proof}
% END SHV5

% BEGIN SHV6
\paragraph{SHV6: the full uniform large-degree exactness theorem.}
\label{claim:SHV6}
For every Noetherian $R$, every finite $\iota$ and every finite
graded $S$-module $Q$, there is $t_0\in\mathbb Z$ such that
\[
 t\ge t_0\quad\Longrightarrow\quad
          A^\bullet(Q,t)\text{ is exact in every degree}.
                                                        \tag{SHV.10}
\]
This is the full original
\texttt{cube\allowbreak Complex\_\allowbreak gr\allowbreak Cube\_\allowbreak exact\allowbreak At\_\allowbreak of\_\allowbreak finite}.

\begin{proof}
Use SHV4 with $a=\nu+1$, obtaining (SHV.9) through $Q_{\nu+1}$.
Set
$t_0=\max_{0\le j\le\nu} b(F_j)$.
For every $t\ge t_0$, SHV3 makes all the complexes
$A^\bullet(F_j,t)$ exact in every degree, and SHV1 turns
(SHV.9) into short exact sequences of complexes.

Fix any $p\ge0$.  The complex $A^\bullet(Q_{\nu+1},t)$ is
zero in degree $p+\nu+1$, because this degree exceeds $\nu$.
It is therefore exact there.  Apply SHV5 to the last short exact
sequence to obtain exactness of $A^\bullet(Q_\nu,t)$ in
degree $p+\nu$.  Applying SHV5 successively gives exactness for
$Q_{\nu-1}$ in degree $p+\nu-1$, and finally for $Q$ in degree
$p$.  Each application uses an already exact free complex.

The same finite list $F_0,\ldots,F_\nu$ and the same maximum
$t_0$ work for all $p$.  No infinite family of degree-dependent
thresholds was replaced by a maximum.  Noetherianity enters only
through the finiteness of the successive graded kernels in SHV4.
This proof also handles $\nu=0$ and arbitrary negative generator
degrees.
\end{proof}
% END SHV6

\subsection{One finite module for the actual sections in every degree}

% BEGIN SHV7
\paragraph{SHV7: finite chart saturation gives all nonempty localized vertices.}
\label{claim:SHV7}
Let $G$ be any graded $S$-module over an arbitrary commutative $R$.
Suppose $G/(X_i-1)G$ is finite over $S$ for every $i\in\iota$.
There is one finite homogeneous submodule $K\subseteq G$ such that
for every integer $t$ and every nonempty $I\subseteq\iota$ the
inclusion induces a bijection
\[
                    K_{I,t}\longrightarrow G_{I,t}.
                                                        \tag{SHV.11}
\]
These bijections are the vertices of the original inclusion map
of localization cubes.  More generally, every homogeneous submodule
satisfying (SHV.12) has these bijections, without a finiteness
assumption on either module.

\begin{proof}
FCD10 gives a finite homogeneous family $g_\ell\in G$ whose
$S$-span $K$ satisfies
\[
       \forall i,\ \forall y\in G,\ \exists k\ge0:
                         X_i^k y\in K.                 \tag{SHV.12}
\]
The submodule is graded because each homogeneous component of a
polynomial multiple of a $g_\ell$ is again such a multiple.
It is finite by its displayed generating family.  This construction
makes no choice involving $t$.

Fix $i\in I$.  If $y/X_I^a$ is a homogeneous fraction, choose
$k$ as in (SHV.12).  Then
$X_I^k y=(\prod_{j\in I\setminus\{i\}}X_j^k)X_i^k y$ belongs
to $K$.  Its homogeneous degree is $t+(a+k)|I|$, so
$(X_I^k y)/X_I^{a+k}$ is a preimage in $K_{I,t}$.
For injectivity, a fraction in $K$ that is zero in $G[X_I^{-1}]$
has a numerator killed by a power of $X_I$ in $G$.  Since $K$
is an actual submodule, that equality also holds in $K$.

The maps are induced by this same inclusion; the restriction
formula following (SHV.5) proves their compatibility.  No
isomorphism at the empty vertex is asserted here: $K_t$ can be
strictly smaller than $G_t$ in small degrees.
\end{proof}
% END SHV7

% BEGIN SHV8
\paragraph{SHV8: the full finite graded module for the actual section cube.}
\label{claim:SHV8}
For the original $f,H,s,M$ satisfying
$H.\operatorname{IsLB}$, $M.\operatorname{IsLB}$ and
$\operatorname{VeryAmpleBy}(f,H,s)$ over any commutative $R$,
there is one finite graded $S$-module $K$ and, for every integer
$t$, a cube map
\[
 \Phi_t:\bigl(I\longmapsto K_{I,t}\bigr)
       \longrightarrow
       \bigl(I\longmapsto\operatorname{Sec}(M(t),D_I)\bigr),
 \qquad D_I=\bigcap_{i\in I}D_i,
                                                        \tag{SHV.13}
\]
bijective at every nonempty vertex.  This is the full original
\texttt{exists\_\allowbreak finite\_\allowbreak graded\_\allowbreak module\_\allowbreak sec\allowbreak Cube};
$K$ is chosen before all $t$.

\begin{proof}
Use the actual graded module $G$ in (SHV.2).
ASL14 gives the original $\pi_i$-semilinear dehomogenization
$G\to P_i$, where
$P_i=\operatorname{Sec}(M(0),D_i)$, together with
\[
                  G/(X_i-1)G\simeq P_i.                 \tag{SHV.14}
\]
This quotient equality is FCD13 applied to the actual map.
ASL16 makes $P_i$ finite projective over
$B_i=\mathcal O(D_i)$.  Since $\pi_i:S\to B_i$ is surjective,
any finite $B_i$-generating family also generates $P_i$ over $S$.
Thus SHV7 applies and chooses one finite homogeneous
$K\subseteq G$.

Compose the localization cube of this inclusion with the
actual cube comparison of ASL12.  On a fraction its formula is
\[
       \Phi_{t,I}\left(\frac{a}{X_I^k}\right)
              =\left(\prod_{i\in I}s_i\right)^{-k}a|_{D_I}.
                                                        \tag{SHV.15}
\]
The factors being inverted are units on $D_I$.  ASL12 proves
this formula compatible with every original restriction and
bijective for the cube of $G$; SHV7 proves bijectivity of the
first map at every nonempty vertex.  For the empty vertex the
map is the actual inclusion $K_t\subseteq G_t$.

The scalar action, the transitions of $M(t)$ and the sections
$s_i$ in (SHV.15) have not changed.  The source allows an
abstract existential module $N$; the concrete choice $N=K$
supplies all its module, grading and decomposition structures.
No Noetherian or smoothness hypothesis was needed for this
choice or its simultaneous all-integer-degree comparison.
\end{proof}
% END SHV8

% BEGIN SHV9
\paragraph{SHV9: uniform positive-degree exactness on the prescribed affine cover.}
\label{claim:SHV9}
Now assume $R$ is Noetherian.  There is one integer $t_0$ such
that for every $t\ge t_0$ the ordinary section complex of
$M(t)$ on $D_0,\ldots,D_m$ is exact in every positive degree.
In particular its degree-one exactness is the original identity
\[
 \operatorname{im}\bigl(\operatorname{Cube.d}(\operatorname{secCube})\,1\bigr)
 =\ker\bigl(\operatorname{Cube.d}(\operatorname{secCube})\,2\bigr).
                                                        \tag{SHV.16}
\]

\begin{proof}
Fix the one module $K$ and its maps from SHV8.  Apply SHV6 to
$K$ and choose its integer bound $t_0$.  At every $t\ge t_0$,
the augmented complex $A^\bullet(K,t)$ is exact in every
degree, and hence its ordinary complex is exact in degrees
$q\ge1$.  Its terms and differentials in those degrees use
only nonempty subsets of the chart index set.

The maps (SHV.13) are bijective on all these vertices and
commute with every restriction.  Taking their finite products
gives an isomorphism of the ordinary complexes with the
actual section complex.  Transport exactness through this
isomorphism.  Ordinary degree one corresponds to subsets of
cardinality two, so its incoming and outgoing maps are exactly
the two maps printed in (SHV.16).

The assertion concerns the given sections on the given cover.
The threshold was fixed from $K$ before choosing $t$ or a
cocycle.  Upgrading it to every cover requires the next
arguments, rather than a change in the meaning of cohomology.
\end{proof}
% END SHV9

\subsection{From the chosen affine cover to every original test cover}

For a line bundle $L$ on a scheme $X$, a section cocycle on opens
$(O_a)$ is a family
$\phi_{ab}\in\operatorname{Sec}(L,O_a\cap O_b)$ satisfying
\[
                   \phi_{ab}+\phi_{bc}=\phi_{ac}
                   \quad\hbox{on }O_a\cap O_b\cap O_c.
                                                        \tag{SHV.17}
\]
Its restrictions satisfy the same identity on every smaller open,
which is exactly the quantifier in the source's
\texttt{Cech\allowbreak One.\allowbreak Is\allowbreak Cocycle}.  A primitive is a family
$\psi_a\in\operatorname{Sec}(L,O_a)$ with
$\phi_{ab}=\psi_a-\psi_b$ on intersections.
Write $\mathcal V_L(V)$ for the assertion that every section cocycle
on every open cover of $V$, with all members contained in $V$,
has such a primitive.  Cover index sets may be arbitrary.

% BEGIN SHV10
\paragraph{SHV10: a primitive on a refinement descends to the original cover.}
\label{claim:SHV10}
Let $\phi$ be a section cocycle on $(O_a)$.  Suppose
$P_k\subseteq O_{\tau(k)}$, every $O_a$ is covered by the
$P_k$, and the restricted cocycle on $(P_k)$ has a primitive
$u_k$.  Then $\phi$ has a primitive on the original $(O_a)$.

\begin{proof}
On $O_a\cap P_k$ define
\[
                   v_{a,k}=\phi_{a,\tau(k)}+u_k.
                                                        \tag{SHV.18}
\]
All terms in this and the following identities are restricted
to the displayed common open.  On
$O_a\cap P_k\cap P_l$, the cocycle identity says
$\phi_{a,\tau(l)}-\phi_{a,\tau(k)}
 =\phi_{\tau(k),\tau(l)}$, whereas the refined primitive
says $\phi_{\tau(k),\tau(l)}=u_k-u_l$.
Hence $v_{a,k}=v_{a,l}$.  The sheaf gluing property gives
one actual section $\psi_a$ on $O_a$ with these restrictions.

On $O_a\cap O_b\cap P_k$,
\[
 \psi_a-\psi_b
   =\phi_{a,\tau(k)}-\phi_{b,\tau(k)}
   =\phi_{ab}.
\]
These opens cover $O_a\cap O_b$, so equality of sections gives
the required identity there.  No finite-cover hypothesis,
separation property of $X$, or local trivialization was needed.
In particular the argument returns a primitive on the original
cover, rather than merely on some selected refinement.
\end{proof}
% END SHV10

% BEGIN SHV11
\paragraph{SHV11: section cocycles are precisely the original coordinate predicate.}
\label{claim:SHV11}
For every actual line bundle $L$,
\[
 \mathcal V_L(X)
       \quad\Longleftrightarrow\quad
             \operatorname{FirstCohomologyVanishes}(L).
                                                        \tag{SHV.19}
\]
The right side retains every original cover, chart assignment,
transition unit and coordinate function.

\begin{proof}
Write $U_p$ for the original bundle charts and $g_{pq}$ for
their units, with coordinate convention $z_p=g_{pq}z_q$.
An original test of the right side consists of opens
$V_a\subseteq U_{c(a)}$ covering $X$ and scalar functions
$z_{ab}$ on $V_a\cap V_b$ with
\[
 z_{ab}+g_{c(a)c(b)}z_{bc}=z_{ac}
                    \quad\hbox{on triple intersections}.
                                                        \tag{SHV.20}
\]
Regard $z_{ab}$ as the $c(a)$-coordinate of a section
$\phi_{ab}$ of $L$ on $V_a\cap V_b$.  The chosen bundle
chart gives a bijection between functions and sections there.
In the $c(a)$-coordinate, the section
$\phi_{ab}+\phi_{bc}$ has the left side of (SHV.20).
Thus these sections satisfy (SHV.17).
If $\mathcal V_L(X)$ holds, let $\psi_a$ be a section primitive
and let $w_a$ be its $c(a)$-coordinate.  Taking that same
coordinate on each overlap gives
\[
             z_{ab}=w_a-g_{c(a)c(b)}w_b,
                                                        \tag{SHV.21}
\]
the literal output required by the original predicate.

Conversely, let a section cocycle be given on any cover $(O_a)$.
Refine by $P_{(a,p)}=O_a\cap U_p$.  These opens cover $X$
and have the displayed original chart assignment $(a,p)\mapsto p$.
Take the $p$-coordinate of $\phi_{ab}$ on
$P_{(a,p)}\cap P_{(b,q)}$.  The section cocycle identity becomes
the original coordinate identity with the transition $g_{pq}$.
Applying the coordinate predicate gives scalar primitives on the
refinement.  The inverse chart identifications turn them into
section primitives with exactly the refined coboundary formula.
SHV10 descends those sections to the original $(O_a)$.
This proves the converse, including infinite and repeated-index
covers.
\end{proof}
% END SHV11

% BEGIN SHV12
\paragraph{SHV12: two finite powers correct a principal-cover cocycle.}
\label{claim:SHV12}
Let $V$ be affine, $B=\mathcal O(V)$, and let
$D_1,\ldots,D_r$, $D_i=D_V(f_i)$, be a finite principal cover.
For any section cocycle $c_{ij}$ of an actual line bundle $L$
on this cover, there are integers $N\ge1$, $K\ge0$ and sections
$\rho_{ik}\in\operatorname{Sec}(L,D_i)$ such that on
$D_i\cap D_j$,
\[
 f_k^K(\rho_{ik}-\rho_{jk})=f_k^{N+K}c_{ij}
                    \qquad(1\le i,j,k\le r).
                                                        \tag{SHV.22}
\]
All occurrences of $f_k$ use its original restriction from $B$.

\begin{proof}
On the affine open $D_i$ the intersection $D_i\cap D_k$
is the principal open defined by $f_k|_{D_i}$.
ASL15 extends $c_{ik}$ after a power of this function.
There are finitely many ordered pairs, so choose one $N\ge1$
at least all their exponents.  Multiplying an already chosen
extension by the extra power gives $\rho_{ik}$ with
\[
                  \rho_{ik}|_{D_i\cap D_k}=f_k^N c_{ik}.
                                                        \tag{SHV.23}
\]
For fixed $i,j,k$, the section
$\rho_{ik}-\rho_{jk}-f_k^N c_{ij}$ on $D_i\cap D_j$
restricts to zero on the intersection with $D_k$.
Indeed (SHV.23) makes that restriction
$f_k^N(c_{ik}-c_{jk}-c_{ij})$, which vanishes by the cocycle
identity.

The open $D_i\cap D_j=D_V(f_i f_j)$ is affine.  Apply the
power-killing part of ASL15 there to the section just displayed
and the function $f_k$.  A second maximum over the finitely many
triples gives one exponent $K$ killing all these discrepancies.
Expanding the resulting identity is precisely (SHV.22).
Increasing an exponent is legitimate because multiplication by
an additional power preserves a zero section.

The first power extends sections; the second makes an identity
hold on the entire pairwise overlap.  Equality only after
restriction to $D_k$ was never cancelled across a possible
zero divisor.  The argument applies over arbitrary $B$, without
a Noetherian hypothesis.
\end{proof}
% END SHV12

% BEGIN SHV13
\paragraph{SHV13: full all-cover first-cohomology vanishing on an affine open.}
\label{claim:SHV13}
For every scheme $X$, every actual line bundle $L$ and every
affine open $V\subseteq X$, the predicate $\mathcal V_L(V)$
holds.  This is the full original
\texttt{cech\allowbreak HOne\allowbreak Vanishes\allowbreak On\allowbreak Open\_\allowbreak of\_\allowbreak is\allowbreak Affine\allowbreak Open}.

\begin{proof}
First solve a section cocycle on a finite principal cover
$D_i=D_V(f_i)$.  The functions $f_i$ generate the unit ideal
in $B=\mathcal O(V)$: a proper ideal containing all of them
would be contained in a maximal ideal, giving a point outside
the cover.  Put $q=N+K\ge1$, with the exponents supplied
by SHV12.  Their powers $f_i^q$ also generate the unit ideal.
For an explicit reason, write $1=\sum_i b_i f_i$ and raise
this identity to $r(q-1)+1$.  In every resulting monomial
some $f_i$ occurs at least $q$ times, so grouping terms yields
\[
                         1=\sum_i a_i f_i^q
                         \quad(a_i\in B).
                                                        \tag{SHV.24}
\]
Define
$u_i=\sum_k a_k f_k^K\rho_{ik}$ on $D_i$.
All coefficients are restricted from $B$.  Equation (SHV.22)
then gives, on $D_i\cap D_j$,
\[
 u_i-u_j
   =\sum_k a_k f_k^{N+K}c_{ij}
   =c_{ij}.
\]
Thus every cocycle on this finite principal cover is a coboundary.

Now start with an arbitrary cover $(O_a)$ of $V$ and an
arbitrary section cocycle.  Principal opens form a basis of the
affine topology, so every point of $V$ has a principal
neighborhood contained in some $O_a$.  Quasi-compactness of $V$
selects finitely many such neighborhoods $D_i$, with
$D_i\subseteq O_{\tau(i)}$.  Restrict the original cocycle to
this finite principal cover and apply the construction above.
SHV10 descends its primitive to the original cover.

If $V$ is empty, all section groups are zero and the zero
family already solves every cocycle.  Thus the empty finite
cover causes no exception.  At no point was the original
cover assumed finite, affine, or subordinate to the bundle
charts; that last issue was handled by actual sections.
\end{proof}
% END SHV13

% BEGIN SHV14
\paragraph{SHV14: local all-cover vanishing plus one cover's cocycles.}
\label{claim:SHV14}
Let $(D_i)$ cover an open $V$, with each $D_i\subseteq V$.
Suppose $\mathcal V_L(D_i)$ holds for every $i$, and suppose
every section cocycle on the particular cover $(D_i)$ has
a primitive there.  Then $\mathcal V_L(V)$ holds.  The index
sets here need not be finite.

\begin{proof}
Take a section cocycle $\phi_{ab}$ on any cover $(O_a)$ of $V$.
On each $D_i$, local vanishing supplies sections
$\sigma_{i,a}$ on $D_i\cap O_a$ such that
\[
                 \phi_{ab}=\sigma_{i,a}-\sigma_{i,b}
                 \quad\hbox{on }D_i\cap O_a\cap O_b.
                                                        \tag{SHV.25}
\]
On $D_i\cap D_j\cap O_a$ form
$\sigma_{i,a}-\sigma_{j,a}$.
For two indices $a,b$ the difference of these expressions on
their common open is
$(\sigma_{i,a}-\sigma_{i,b})
 -(\sigma_{j,a}-\sigma_{j,b})=\phi_{ab}-\phi_{ab}=0$.
They therefore glue to an actual section
$\theta_{ij}$ on $D_i\cap D_j$.

On every $D_i\cap D_j\cap D_k\cap O_a$ one has
$\theta_{ij}+\theta_{jk}=\theta_{ik}$ by cancellation of
the three $\sigma$-terms.  As these opens cover the triple
intersection, $\theta$ is a section cocycle on $(D_i)$.
By the hypothesis choose $\xi_i$ with
$\theta_{ij}=\xi_i-\xi_j$.

For fixed $a$, the sections
$\sigma_{i,a}-\xi_i$ agree on
$D_i\cap D_j\cap O_a$, by the definition of $\theta$.
Glue them on the covering opens $D_i\cap O_a$ to a section
$\psi_a$ on $O_a$.  On $D_i\cap O_a\cap O_b$,
$\psi_a-\psi_b=\sigma_{i,a}-\sigma_{i,b}=\phi_{ab}$.
Sheaf equality gives this identity on all of $O_a\cap O_b$.
These $\psi_a$ solve the original test cocycle.

The proof uses only actual restrictions and the sheaf axiom;
it does not require the pairwise intersections $D_i\cap D_j$
to be affine or quasi-compact.
\end{proof}
% END SHV14

% BEGIN SHV15
\paragraph{SHV15: the exact ordered differential and its coboundary sign.}
\label{claim:SHV15}
On a linearly ordered cover $(D_i)$, degree-one exactness of
the ordinary section complex is equivalent to the assertion
that every section cocycle on that cover has a primitive.
For a finite index set this is exactly the relationship between
the two \texttt{Cube.\allowbreak d} maps in (SHV.16) and the source's
\texttt{Is\allowbreak Coboundary}.

\begin{proof}
An ordered degree-one cochain is a family $x_{ij}$ for $i<j$,
and the first two differentials are
\[
 (\delta u)_{ij}=u_j-u_i,\qquad
 (\delta x)_{ijk}=x_{jk}-x_{ik}+x_{ij}\quad(i<j<k).
                                                        \tag{SHV.26}
\]
Restrictions to the appropriate common open are implicit.
These are the signs in (SHV.4), not a newly chosen convention.

A section cocycle $\phi$ has $\phi_{ii}=0$, since its
cocycle equation with three equal indices gives
$\phi_{ii}+\phi_{ii}=\phi_{ii}$.  It also has
$\phi_{ji}=-\phi_{ij}$, by the equation with first and third
indices equal.  Taking $x_{ij}=\phi_{ij}$ for $i<j$ gives
$\delta x=0$.  Exactness supplies $u$ with $\delta u=x$.
Set $\psi_i=-u_i$.  Then
$\phi_{ij}=\psi_i-\psi_j$ for $i<j$ by (SHV.26);
the diagonal and antisymmetry identities extend it to every
ordered pair.  The minus sign is needed because the source's
primitive uses first minus second, while $\delta u$ uses
second minus first.

Conversely, let $\delta x=0$.  Define $\phi_{ij}=x_{ij}$
for $i<j$, define $\phi_{ij}=-x_{ji}$ for $j<i$, and put
$\phi_{ii}=0$.  The expression
$\phi_{ij}+\phi_{jk}-\phi_{ik}$ on a triple intersection
changes sign under either adjacent transposition of its
indices, by antisymmetry, and is zero when two indices agree.
For increasing indices it is zero by $\delta x=0$.
Sorting the three indices therefore proves (SHV.17) for
every triple.  A cocycle primitive $\psi$ gives
$x=\delta(-\psi)$, proving exactness.  This argument is also
valid for an infinite ordered index set, where the cochains
are products and each displayed differential still uses only
finitely many terms at each component.
\end{proof}
% END SHV15

% BEGIN SHV16
\paragraph{SHV16: exactness on an affine cover implies the full coordinate predicate.}
\label{claim:SHV16}
Let $L$ be an actual line bundle on any scheme $X$.
If a linearly ordered affine open cover $(D_i)$ has its
ordinary section complex exact in degree one, then
$\operatorname{FirstCohomologyVanishes}(L)$ holds.
In particular, for every finite ordered cover this proves
the full original
\texttt{first\allowbreak Cohomology\allowbreak Vanishes\_\allowbreak of\_\allowbreak exact\_\allowbreak sec\allowbreak Cube}
with its exact \texttt{Cube.\allowbreak d} hypothesis.

\begin{proof}
SHV13 gives $\mathcal V_L(D_i)$ for each affine member.
SHV15 translates the given exactness into the existence of
primitives for every section cocycle on the particular
cover $(D_i)$.  SHV14 now gives $\mathcal V_L(X)$.
Finally SHV11 gives the literal coordinate predicate.

Thus for each arbitrary original chart-subordinate cover
and each of its scalar cocycles, the construction returns
coordinate functions on that same cover satisfying (SHV.21).
No assumption on the cardinality of this test cover has
entered, and no generic sheaf-cohomology comparison is left
as an unproved premise.
\end{proof}
% END SHV16

% BEGIN SHV17
\paragraph{SHV17: full original Serre first-cohomology vanishing.}
\label{claim:SHV17}
Let $R$ be any Noetherian commutative ring,
$f:Y\to\operatorname{Spec}R$ any morphism, and let
$H,M$ be the actual line bundles in an original
$\operatorname{VeryAmpleBy}(f,H,s)$ presentation.
Then there is an integer $t_0$ such that
\[
 \forall t\in\mathbb Z,\quad t\ge t_0
   \ \Longrightarrow\
       \operatorname{FirstCohomologyVanishes}(M\otimes H^t).
                                                        \tag{SHV.27}
\]
This proves \texttt{exists\_\allowbreak h1Zero\_\allowbreak tw\_\allowbreak of\_\allowbreak very\allowbreak Ample\allowbreak By}
at its full original scope.

\begin{proof}
Choose the single finite graded module of SHV8 and the single
bound in SHV9.  At every $t\ge t_0$, the actual section complex
of $M(t)$ on the prescribed $D_i=D(s_i)$ is exact in degree one.
The opens $D_i$ are affine and cover $Y$, and $M(t)$ is a line
bundle for every integer $t$ because its transition functions
are the original units of $M$ times the $t$th powers of those
of $H$.  SHV16 applies to these actual data and gives (SHV.27).

The choices of the finite graded submodule, its finite list of
graded kernels and the bound precede $t$.  For each such $t$
the conclusion quantifies every cover, every assignment of an
original $M(t)$-chart to a cover member, and every coordinate
cocycle on it.  The solution is on that original cover and
uses the original transition units.  Smoothness, reducedness,
integrality, a field hypothesis and a global section
base-change theorem have not been assumed.
\end{proof}
% END SHV17

\subsection{The accompanying high-power polynomial statement}

% BEGIN SHV18
\paragraph{SHV18: a saturated finite span contains every sufficiently high piece.}
\label{claim:SHV18}
Suppose $R$ is Noetherian and $G$ is a graded $S$-module,
which need not be finite.  Let $g_1,\ldots,g_r$ be homogeneous
and let $K=\sum_\ell Sg_\ell$.  Assume
\[
 \begin{split}
 &\forall i,\ \forall y\in G,\ \exists k\ge0:\ X_i^k y\in K,\\
 &\forall y\in G,\ 
       \bigl(\forall i,\ \exists k\ge0:\ X_i^k y=0\bigr)
           \Longrightarrow y=0.
 \end{split}                                           \tag{SHV.28}
\]
Then there is $t_0\in\mathbb Z$ such that $G_t\subseteq K$
for every $t\ge t_0$.  This is the full original
\texttt{mem\_\allowbreak span\_\allowbreak of\_\allowbreak saturating\_\allowbreak of\_\allowbreak le}.

\begin{proof}
The module $K$ is finite and graded.  The first condition in
(SHV.28) is precisely the saturation used in the proof of
(SHV.11), so that proof gives bijections
$K_{I,t}\to G_{I,t}$ for all nonempty $I$.
There is no need here to assume finite dehomogenized quotients:
the finite family and its saturation have already been supplied.
The bijections commute with every restriction.

Choose a bound for $K$ by SHV6.  For $t$ above it and
$y\in G_t$, the images of $y$ in the singleton localizations
form a degree-zero cocycle of the ordinary complex of $G$:
the restrictions of the same $y$ agree on every overlap.
Lift this cocycle through the isomorphism of ordinary
localization complexes to a cocycle $z$ for $K$.
Exactness of $A^\bullet(K,t)$ in augmented degree one gives
$a\in K_t$ whose singleton images are $z$.

Hence $y-a$ maps to zero in $G[X_i^{-1}]$ for every $i$.
The fraction criterion gives some $k_i$ with
$X_i^{k_i}(y-a)=0$ for each $i$.  The second condition
in (SHV.28) forces $y=a$, proving $y\in K$.
The same bound works for every $y$ in every sufficiently
large degree.  When the variable set is empty, that second
condition already forces $G=0$, so the conclusion still holds.
\end{proof}
% END SHV18

% BEGIN SHV19
\paragraph{SHV19: the actual graded sections have no common chart-power torsion.}
\label{claim:SHV19}
Let $s_i$ be sections of an actual line bundle $H$ whose
nonzero loci cover $Y$, and let $M$ be any actual line bundle.
For the graded section module $G$ in (SHV.2),
\[
       \bigl(\forall i,\ \exists k\ge0:\ X_i^k y=0\bigr)
                    \quad\Longrightarrow\quad y=0 .
                                                        \tag{SHV.29}
\]
No affine-chart or Noetherian hypothesis is needed.

\begin{proof}
Write $y$ as its finite sum of homogeneous sections $y_t$.
Since multiplication by $X_i^k$ shifts each homogeneous
degree by $k$, the equality $X_i^k y=0$ implies
$s_i^k y_t=0$ separately for every $t$.  On $D(s_i)$
the coordinate of $s_i$ is a unit, by ASL2.  Thus
$y_t|_{D(s_i)}=0$.  The opens $D(s_i)$ cover $Y$,
so the sheaf equality of actual sections gives $y_t=0$
for every component.  Their finite sum is zero.

In this argument one cancels $s_i$ only on its nonzero
locus.  There is no assertion that multiplication by $s_i$
is injective on all global sections.
\end{proof}
% END SHV19

% BEGIN SHV20
\paragraph{SHV20: a degree-zero generator in a chart gives power saturation.}
\label{claim:SHV20}
For this claim the polynomial variable set may be arbitrary;
fix one variable $X_i$.
Let $\pi:S\to B$ be a surjective ring map with $\pi(X_i)=1$,
and let $\epsilon:G\to P$ be $\pi$-semilinear, where $G$ is
graded and $P$ is a $B$-module.  Suppose every homogeneous
element of $\ker\epsilon$ is killed by some power of $X_i$.
Let $e\in G_0$ and suppose every element of $P$ is
$b\,\epsilon(e)$ for some $b\in B$.  For every $S$-submodule
$K\subseteq G$ containing $e$ and every $y\in G$, some
$X_i^k y$ belongs to $K$.

\begin{proof}
First let $y\in G_t$ be homogeneous, allowing $t<0$.
Choose $b\in B$ with $\epsilon(y)=b\epsilon(e)$ and choose
a polynomial $F$ with $\pi(F)=b$.
Choose an integer $D\ge0$ at least $t$ and the total degree
of every monomial of $F$.  This is a finite choice even
when the variable set is infinite.  Homogenize with $X_i$:
\[
 F_D=\sum_\alpha c_\alpha X_i^{D-|\alpha|}X^\alpha .
                                                        \tag{SHV.30}
\]
It is homogeneous of degree $D$ and $\pi(F_D)=\pi(F)$
because $\pi(X_i)=1$.
Both $X_i^{D-t}y$ and $F_D e$ have degree $D$ and have
the same image under $\epsilon$.
By the homogeneous-kernel hypothesis, for some $a\ge0$,
\[
 X_i^a(X_i^{D-t}y-F_D e)=0.
\]
Thus $X_i^{a+D-t}y=X_i^aF_D e$ belongs to $K$.
All exponents used as powers are nonnegative.

For a general $y$, apply this argument to its finitely many
homogeneous components.  Increase their exponents to a common
maximum, using that $K$ is an $S$-submodule, and add.
This proves the assertion without assuming $\epsilon$
surjective, $G$ finite or the base ring Noetherian.
It is the full original
\texttt{exists\_\allowbreak X\_\allowbreak pow\_\allowbreak smul\_\allowbreak mem\_\allowbreak span\_\allowbreak of\_\allowbreak dehomogenisation},
including its arbitrary containing submodule.
\end{proof}
% END SHV20

% BEGIN SHV21
\paragraph{SHV21: every high-power section is an actual homogeneous polynomial.}
\label{claim:SHV21}
For every original $\operatorname{VeryAmpleBy}(f,H,s)$ over
a Noetherian $R$, there is a natural number $T$ such that
for every natural number $t\ge T$ and every actual section
$\sigma$ of $H^t$, there is $F\in R[X_0,\ldots,X_m]$,
homogeneous of degree $t$, with
\[
 \sigma_k=
  \operatorname{eval}\bigl(c_f(U_k),(s_j)_k\bigr)(F)
                 \qquad\text{for every original }H\text{-chart }U_k.
                                                        \tag{SHV.31}
\]
This proves the full original \texttt{exists\_\allowbreak form\_\allowbreak eq\_\allowbreak sec\_\allowbreak pw}.
Together with SHV17 it proves both original conjuncts of
\texttt{serre\_\allowbreak h1Zero\_\allowbreak and\_\allowbreak sections\_\allowbreak polynomial}.

\begin{proof}
Use the actual trivial line bundle for $M$ in (SHV.2).
Let $e\in G_0$ be the section whose coordinate is $1$.
The canonical zero-twist identification identifies
$\operatorname{Sec}(M(0),D_i)$ with $B_i=\mathcal O(D_i)$,
and the degree-zero restriction formula of ASL14 sends
$e$ to $1$ on each chart.  Every element of that chart
module is therefore a multiple of $\epsilon_i(e)$.
The original evaluation $\pi_i$ is surjective and has
$\pi_i(X_i)=1$, and ASL14 supplies its homogeneous
power-kernel property.

Apply SHV20 to each chart with the single submodule $K=S e$.
It gives the first condition in (SHV.28) for every element
of $G$, while SHV19 gives the second.  The homogeneous
family consists of this one degree-zero element $e$.
SHV18 now supplies one integer $t_0$ with $G_t\subseteq S e$
for all $t\ge t_0$.  Choose a natural number $T$ with
$T\ge\max(0,t_0)$.

An actual section $\sigma$ of $H^t$ defines an element
$y\in G_t$ by retaining its coordinates in the factor
$\mathcal O_Y\otimes H^t$.  This uses the trivial bundle's
single chart and coordinate $1$, so it introduces no choice
of trivialization or of a new section.
For $t\ge T$, write $y=F'e$.
The map $S\to G$, $F'\mapsto F'e$, preserves degree.
Taking its degree-$t$ component shows
$y=F e$ for the degree-$t$ homogeneous part $F$ of $F'$.
This remains valid for $F=0$.

Finally the definition of the polynomial action gives its
coordinate formula.  Multiplication by $X_j$ multiplies by
the original coordinate $(s_j)_k$, and a coefficient $r\in R$
acts by $c_f(U_k)(r)$.  Iterating for a monomial and adding
its finitely many coefficients shows that the $k$-coordinate
of $F e$ is exactly the right side of (SHV.31), because
$e$ has coordinate one.  The coordinate of $y$ is $\sigma_k$.
Thus (SHV.31) holds on every original chart, with one
homogeneous polynomial and the original scalar map.

The bound $T$ was fixed before $t$ and before $\sigma$.
Unlike SHV17's bound, which can depend on the arbitrary
bundle $M$, this bound uses only the fixed trivial-bundle
graded construction for $f,H,s$.
SHV17 and this assertion have exactly the two quantifier
orders in the original conjunction, so pairing them gives
that conjunction without any additional geometric premise.
\end{proof}
% END SHV21

\subsection{Return to the remaining section-law arguments}

SHV17 supplies the full input named \texttt{GSL\_\allowbreak SERRE} in
the earlier contracts.  That input was explicitly the first,
all-cover $H^1$ conjunct; SHV21 additionally supplies the
high-power polynomial conjunct of the original Serre theorem.
Neither conclusion assumes the smooth-projective \texttt{Perf}
record.  ASL21 supplies that record independently at its stated
smoothness hypothesis; this chapter uses the earlier actual
localization and affine-section arguments before smoothness enters.

In GSL17 the Serre input can therefore be removed from the retained
premises.  Its positive-power vanishing conclusion still requires the
original polarized base-change framework, the multiplication-pullback
generated line-bundle comparison, finite/flat/surjective and etale
multiplication geometry, and descent of every invariant function on
every relevant open.  The prime and residue-class argument in GSL17
then applies exactly as proved.  Those remaining producers have not
been supplied by a theorem about high twists.

Likewise the principal rank-one theorem, geometric integrality,
torsion-indexed section decomposition, the required two-sided section
growth, and the square/translation generators remain distinct inputs
to the later global section laws.  A polynomial description in high
degree alone does not identify those dimensions or produce those
geometric comparisons.  No whole selected RH or Dirichlet-GRH
obligation is declared closed by this chapter.

\paragraph{Pinned source map.}
All source references are to
\texttt{Coco\allowbreak Sgt/\allowbreak riemann-hypothesis-lean}, commit
\texttt{7cbf7a0c32fca1badf05931879a37c3f8a2b03b4}.
Write $C$ for
\texttt{src/\allowbreak Development/\allowbreak Abelian\allowbreak Variety\allowbreak Theory/\allowbreak Proofs/\allowbreak Coherent\allowbreak Cohomology\allowbreak Proofs/\allowbreak },
$A$ for the sibling
\texttt{Proofs/\allowbreak Abelian\allowbreak Scheme\allowbreak Law\allowbreak Proofs/\allowbreak }, and $D$ for the sibling
\texttt{Definitions/\allowbreak Quasicoherent\allowbreak Cohomology/\allowbreak }.
The accompanying provenance gives exact source ranges and hashes.
The principal original declarations are as follows.

\begin{itemize}
\item SHV1--SHV6: $D$ \texttt{Exists\allowbreak Away\allowbreak Frac\allowbreak Of\allowbreak Mem\allowbreak Hom\allowbreak Loc.\allowbreak lean},\par
$D$ \texttt{Is\allowbreak Graded\allowbreak Map\allowbreak Linear\allowbreak Combination.\allowbreak lean}, and
$C$ \texttt{Cube\allowbreak Complex\allowbreak Graded\allowbreak Module\allowbreak Exact\allowbreak At\allowbreak Of\allowbreak Large\allowbreak Degree.\allowbreak lean}.
The free calculation is in
$C$ \texttt{Exists\allowbreak Monomial\allowbreak Basis\allowbreak Cech\allowbreak Cube\allowbreak Polynomial\allowbreak Grading.\allowbreak lean},
\texttt{Cube\allowbreak Complex\allowbreak Exact\allowbreak At\allowbreak Of\allowbreak Monomial\allowbreak Basis.\allowbreak lean},
\texttt{Cube\allowbreak Complex\allowbreak Poly\allowbreak Grading\allowbreak Finite\allowbreak Exact.\allowbreak lean} and
\texttt{Cube\allowbreak Complex\allowbreak Gr\allowbreak Free\allowbreak Finite\allowbreak And\allowbreak Exact\allowbreak At.\allowbreak lean}.
\item SHV7--SHV9: $A$
\texttt{Exists\allowbreak Finite\allowbreak Graded\allowbreak Module\allowbreak Line\allowbreak Bundle\allowbreak Sections.\allowbreak lean},
especially \texttt{exists\_\allowbreak finite\_\allowbreak gr\allowbreak Cube\_\allowbreak hom\_\allowbreak bijective} and
\texttt{exists\_\allowbreak finite\_\allowbreak graded\_\allowbreak module\_\allowbreak sec\allowbreak Cube}.
\item SHV10--SHV16: $A$
\texttt{Cech\allowbreak HOne\allowbreak Vanishes\allowbreak On\allowbreak Open\allowbreak Of\allowbreak Is\allowbreak Affine\allowbreak Open.\allowbreak lean},\par
from \texttt{is\allowbreak Coboundary\_\allowbreak of\_\allowbreak refine} through
\texttt{h1\allowbreak Zero\_\allowbreak of\_\allowbreak exact\allowbreak At\_\allowbreak one}; and $C$
\texttt{First\allowbreak Cohomology\allowbreak Vanishes\allowbreak Of\allowbreak Exact\allowbreak Sec\allowbreak Cube.\allowbreak lean}.
The literal coordinate predicate is in
\texttt{Definitions/\allowbreak Abelian\allowbreak Scheme\allowbreak Laws/\allowbreak First\allowbreak Cohomology\allowbreak Vanishes\allowbreak Spec\allowbreak Of\allowbreak Is\allowbreak Local\allowbreak Ring.\allowbreak lean}.
\item SHV17: $C$ \texttt{Serre\allowbreak Vanishing\allowbreak H1\allowbreak Twist.\allowbreak lean},
\texttt{exists\_\allowbreak h1\allowbreak Zero\_\allowbreak tw\_\allowbreak of\_\allowbreak very\allowbreak Ample\allowbreak By}.
\item SHV18--SHV21: $C$ \texttt{Mem\allowbreak Span\allowbreak Of\allowbreak Saturating\allowbreak Of\allowbreak Le.\allowbreak lean}
and \texttt{Exists\allowbreak Homogeneous\allowbreak Form\allowbreak Eq\allowbreak High\allowbreak Power\allowbreak Section.\allowbreak lean};
the final two-conjunct interface is $A$
\texttt{Serre\allowbreak Vanishing\allowbreak Very\allowbreak Ample.\allowbreak lean}.
\end{itemize}


\section{The actual arithmetic local charts and their residual representations}
\label{sec:arithmetic-local-charts-v055}

\paragraph{What the local construction must preserve.}
The component argument uses the entry algebra of the supplied universal
representation, not an arbitrary ring with a similar spectrum.  At a
closed point it passes to a specified coefficient extension and a
complete local ring.  We prove that the original universal property
makes this ring the actual completion, prove its flatness, and identify
the representation and all three entry ideals after that map.  We also
construct these local rings at every original permitted closed point.
The residual argument treats ranks two and three separately where
their matrix patterns differ.

The upstream results RTC7 and RTC9 are available at their full original
scope.  Thus the actual weight-zero coefficient ring is complete
Noetherian local, the actual graded entry algebra is Noetherian and
finite type over it, and all entries and ordinary characters have the
required adic continuity.  None of those results assumes a later
\(\mathrm{FinA}\), finite-point, propagation or automorphic conclusion.

\paragraph{Notation.}
In the arithmetic statements let \(O\) be the original complete
discrete valuation ring, with finite residue field \(k\) of
characteristic \(\ell\), and let \(F\) satisfy the original
\(F.\mathrm{Ok}(C_0)\).  Fix \(n=2\) or \(3\) and a supplied
\(S:\mathrm{Space}(O,F,n)\).  Write \(\mathcal R=S.R\),
\(B=S.B\), \(A=S.A\), and \(u=S.u\).  Their inclusions are
the original inclusions in \(\mathcal R\).  Put
\[
 \mathfrak m_B=B\cap\mathfrak m_{\mathcal R},\qquad
 x_0=\mathfrak m_B A,\qquad
 \pi(\mathfrak p)=\mathfrak p\cap B.
                                                        \tag{ALC.1}
\]
The weights are \(w(2)=(1,-1)\) and \(w(3)=(1,0,-1)\).
The ideals \(\mathfrak t_+,\mathfrak t_-\) are generated by the
entries below and above the diagonal.  The ideal \(\mathfrak e\)
is generated by entries whose weight difference is odd.  Thus
\[
 \mathrm{TI}(\mathfrak p)
 \Longleftrightarrow
 \mathfrak t_+\not\subseteq\mathfrak p,
 \quad\mathfrak t_-\not\subseteq\mathfrak p,
 \quad(n=3\Rightarrow\mathfrak e\not\subseteq\mathfrak p).
                                                        \tag{ALC.2}
\]
An original \(L:S.\mathrm{Loc}(x)\) contains a finite coefficient
extension \(O\to O'\), with
\(\mathfrak m_OO'=\mathfrak m_{O'}\), a coefficient algebra
\(T=L.R\) over \(O'\), and the specified \(O\)-algebra map
\(\iota:A\to T\).  It requires
\(\iota^{-1}\mathfrak m_T=x\) and the residue-compatible
universal property used explicitly below.  Its representation is
\(\rho_L=u\otimes\iota\); for an ideal \(P\subset T\) write
\(L.\mathrm{im}(P)=\iota^{-1}P\) and
\(L.\pi(P)=\pi(\iota^{-1}P)\).

The first completion statements have the wider original algebraic
scope: \(O\) is merely a commutative local ring.  A coefficient
algebra over it means exactly \(\mathrm{IsCO}\): a complete
Noetherian local \(O\)-algebra, with a local structure map and
the specified residue identification with \(\operatorname{Res}(O)\).
No finiteness of that residue field is needed in these statements.

\subsection{The universal property determines every finite quotient}

% BEGIN ALC1
\paragraph{ALC1. Locality of coefficient maps and the nilpotent coefficient quotients.}
\label{claim:ALC1}
Every \(O\)-algebra map between coefficient algebras over a
commutative local ring \(O\) is local, preserves the specified
residue maps, and is continuous for all powers of the maximal
ideals.  If \(D\) is a Noetherian \(O\)-algebra and
\(\varphi:D\to k=\operatorname{Res}(O)\) is an
\(O\)-algebra map, put \(y=\ker\varphi\).  Then
\(D/y^N\) is a coefficient algebra for every \(N>0\), with
nilpotent maximal ideal \(y/y^N\) and residue map induced by
\(\varphi\).

\begin{proof}
Let \(h:R\to C\) be such a map and \(r\in\mathfrak m_R\).
Choose \(a\in O\) whose image in \(C\) has the same residue
as \(h(r)\).  If \(a\) were a unit, then \(r-a\) would be
a unit in \(R\), whereas \(h(r-a)\in\mathfrak m_C\).
This is impossible.  Thus \(a\in\mathfrak m_O\), and the
local structure map gives \(h(r)\in\mathfrak m_C\).
For arbitrary \(r\), choose \(a\) with \(r-a\in\mathfrak m_R\)
and apply this result.  It proves
\(\operatorname{res}_C(h(r))=\operatorname{res}_R(r)\).
Also \(h(\mathfrak m_R^N)\subseteq\mathfrak m_C^N\) for
every \(N\), which is precisely the required continuity.

The map \(\varphi\) is surjective because its restriction to
\(O\) is the residue map.  Hence \(y\) is maximal.  In
\(D/y^N\), the ideal \(y/y^N\) is nilpotent and its quotient
is \(k\).  An element outside this ideal has an inverse modulo
the ideal; the finite geometric series corrects that inverse to
an actual inverse.  Consequently this is the unique maximal ideal.
Its \(N\)th power is zero, so its adic completion is the ring
itself and its adic filtration is separated.  Noetherianity follows
from that of \(D\).  The map from \(O\) induces the original
residue map and sends nonunits into the maximal ideal.  These are
all the \(\mathrm{IsCO}\) clauses.  In particular this applies
to \(R/\mathfrak m_R^N\) whenever \(R\) is a coefficient
algebra.  It also applies with
\(\varphi=\operatorname{res}_R\circ\iota\) for a supplied
map \(\iota:D\to R\).
\end{proof}
% END ALC1

% BEGIN ALC2
\paragraph{ALC2. The two quotient maps are actual inverses.}
\label{claim:ALC2}
Let \(D\) be a Noetherian \(O\)-algebra, let \(R\) be a
coefficient algebra, and let \(\iota:D\to R\) be an
\(O\)-algebra map.  Assume the original universal property:
for every coefficient algebra \(C\) and every \(f:D\to C\)
with \(\operatorname{res}_C f=\operatorname{res}_R\iota\),
there is a unique \(O\)-algebra map \(g:R\to C\) satisfying
\(g\iota=f\).  Put \(y=\iota^{-1}\mathfrak m_R\).
For every \(N>0\), the map induced by \(\iota\) is an
isomorphism
\[
 \alpha_N:D/y^N\xrightarrow{\ \sim\ }R/\mathfrak m_R^N.
                                                        \tag{ALC.3}
\]
These isomorphisms commute with every reduction to a smaller
positive exponent.

\begin{proof}
ALC1 makes \(D/y^N\) a coefficient algebra with the required
residue map.  Apply the universal property to its quotient map
to obtain \(g_N:R\to D/y^N\), with
\(g_N\iota=q_D\).  By ALC1 the map \(g_N\) is local.
The target maximal ideal has \(N\)th power zero, so
\(g_N(\mathfrak m_R^N)=0\).  It therefore factors as
\(\beta_N:R/\mathfrak m_R^N\to D/y^N\).
Conversely \(\iota(y^N)\subseteq\mathfrak m_R^N\) defines
\(\alpha_N\).

The equality \(\beta_N\alpha_N=1\) follows by evaluating on
the surjective image of \(D\): both composites send \(d\)
to its class modulo \(y^N\).  To prove the other equality,
use the universal property with the coefficient algebra
\(R/\mathfrak m_R^N\).  The two maps
\(q_R\) and \(\alpha_Ng_N\) from \(R\) agree after
\(\iota\), and their common prescribed map has the required
residue.  Uniqueness gives \(\alpha_Ng_N=q_R\), hence
\(\alpha_N\beta_N=1\).  Reduction compatibility for
\(\alpha_N\) follows directly from their definition on \(D\);
their inverses consequently have the same compatibility.
\end{proof}
% END ALC2

% BEGIN ALC3
\paragraph{ALC3. Full completion from the original universal property.}
\label{claim:ALC3}
Under exactly the hypotheses of ALC2, \(\iota\) satisfies the
original \(\mathrm{IsCplAt}(y,R)\) record:
\[
 \iota^{-1}(\mathfrak m_R^N)=y^N,
 \qquad
 \forall r\in R\ \exists d\in D:
           r-\iota(d)\in\mathfrak m_R^N
                     \quad(N\ge0),                  \tag{ALC.4}
\]
and \(R\) is Noetherian and complete for its maximal ideal.
This proves the full original \(\mathrm{cpl\_\allowbreak of\_\allowbreak univ}\).

\begin{proof}
For \(N>0\), injectivity of (ALC.3) is exactly the first
identity in (ALC.4), and surjectivity is exactly the second.
For \(N=0\), the first identity says that the inverse image
of the unit ideal is the unit ideal; the second holds with
\(d=0\).  Noetherianity and completeness are already fields
of the supplied coefficient-algebra structure.  These are
literally the four fields in \(\mathrm{IsCplAt}\).

In particular, the compatible isomorphisms (ALC.3), followed
by completeness and separatedness of \(R\), identify it with
\(\varprojlim_N D/y^N\).  This identification sends the
original element \(d\) to \(\iota(d)\); it does not replace
the original map with an unmarked isomorphism of rings.
\end{proof}
% END ALC3

% BEGIN ALC4
\paragraph{ALC4. The full original arithmetic local-completion theorem.}
\label{claim:ALC4}
Return to the arithmetic data and let \(L:S.\mathrm{Loc}(x)\)
be any supplied local object.  For every \(O'\)-algebra map
\[
 \iota':D=O'\otimes_OA\longrightarrow T,
 \qquad \iota'(1\otimes a)=\iota(a)\quad(a\in A),
                                                        \tag{ALC.5}
\]
the original \(\mathrm{IsCplAt}\) record holds at
\(y=(\iota')^{-1}\mathfrak m_T\).  This is the full
original \(\mathrm{loc\_\allowbreak cpl}\), for both ranks and every
original \(x,L,\iota'\).

\begin{proof}
By RTC7, \(A\) is Noetherian.  The specified \(O'\) is a
finite \(O\)-module, so \(D\) is a finite \(A\)-module:
tensor a finite set of \(O\)-module generators by \(A\).
It is therefore a Noetherian ring.

Let \(C\) be any coefficient algebra over \(O'\), and let
\(f:D\to C\) be an \(O'\)-algebra map whose residue agrees
with \(\iota'\).  Give \(C\) its composite \(O\)-algebra
structure and the resulting scalar tower.  Restrict \(f\)
along \(a\mapsto1\otimes a\).  By (ALC.5), its residue
agrees with \(\iota:A\to T\).  The actual \(L.\mathrm{univ}\)
therefore supplies a unique \(O'\)-algebra map \(g:T\to C\)
agreeing on \(A\).

For every pure tensor \(b\otimes a\), \(O'\)-linearity
and the equality on \(A\) give
\(g(\iota'(b\otimes a))=b f(1\otimes a)=f(b\otimes a)\).
Pure tensors additively generate \(D\), so \(g\iota'=f\).
Any other such map agrees on the original \(A\), hence is
equal by \(L.\mathrm{univ}\).  We have verified every
antecedent of ALC3 with coefficient ring \(O'\), the given
\(T\), and this very \(\iota'\).  Its conclusion is the
required record.
\end{proof}
% END ALC4

\subsection{Completion is flat, and the primes remain attached to the original map}

% BEGIN ALC5
\paragraph{ALC5. Constructing the Noetherian completion and its exact powers.}
\label{claim:ALC5}
Let \(E\) be a Noetherian commutative ring, let
\(I=(a_1,\ldots,a_r)\), and define
\(\widehat E=\varprojlim_N E/I^N\).  Then
\[
 \ker(\widehat E\to E/I^N)=I^N\widehat E,
 \qquad \widehat E/I^N\widehat E=E/I^N.          \tag{ALC.6}
\]
The ring \(\widehat E\) is Noetherian, complete and separated
for \(I\widehat E\), and \(E\) is dense in it.  If \(I\)
is maximal, \(\widehat E\) is local with maximal ideal
\(I\widehat E\), and the source's particular
\(\mathrm{Cpl}(I)=\widehat{E_I}\) has the full
\(\mathrm{IsCplAt}(I,\mathrm{Cpl}(I))\) record.  The last
assertion is the original \(\mathrm{cpl\_\allowbreak is}\).

\begin{proof}
Formal substitution defines a ring map
\[
 E[[X_1,\ldots,X_r]]\longrightarrow\widehat E,
                \qquad X_i\longmapsto a_i.           \tag{ALC.7}
\]
It is defined by finite truncations: terms of total degree at
least \(N\) vanish modulo \(I^N\).  Thus multiplication
and addition are legitimate in every finite quotient.

For \(z\in\widehat E\), choose its constant coefficient
in \(E\) to match \(z\) modulo \(I\).  Suppose terms of
degrees below \(N\) have been chosen.  Their discrepancy
with \(z\), reduced modulo \(I^{N+1}\), belongs to
\(I^N/I^{N+1}\).  The degree-\(N\) monomials in the
\(a_i\) generate that module over \(E\), so choose their
coefficients to correct the discrepancy.  Continuing gives a
formal series mapping to \(z\).  This proves surjectivity of
(ALC.7).  If \(z\) is zero modulo \(I^N\), start with all
terms of degree below \(N\) equal to zero.  The resulting
series lies in \((X_1,\ldots,X_r)^N\), and hence
\(z\in I^N\widehat E\).  The reverse inclusion is immediate.
This proves (ALC.6), including the case \(N=0\).

The finite-variable formal power-series ring over a Noetherian
ring is Noetherian; its surjective image (ALC.7) is consequently
Noetherian.  The exact kernel identities identify the adic
quotients of \(\widehat E\) with the inverse system defining
it.  They prove completeness and separatedness, and their
surjectivity from \(E\) proves density.

When \(I\) is maximal, the quotient by \(I\widehat E\) is
the field \(E/I\).  An element invertible in that quotient
has a lift of its inverse whose product differs from one by
an element of \(I\widehat E\).  The convergent geometric
series corrects the lift.  Therefore every element outside
\(I\widehat E\) is a unit, proving the local assertion.

For the source's particular completion, observe that
\(E/I^N\) is already local when \(I\) is maximal and
\(N>0\): its maximal ideal is nilpotent and its quotient
is a field.  Every element of \(E\setminus I\) therefore
becomes a unit in it.  Localization consequently gives
\[
        E/I^N\simeq E_I/(IE_I)^N.
\]
These maps respect all reductions and the original maps
from \(E\).  Passing to the inverse limits identifies
\(\widehat E\) with \(\widehat{E_I}\) as an
\(E\)-algebra.  The exact quotient identities and density
above give all four \(\mathrm{IsCplAt}\) fields for this
specific ring and map.
\end{proof}
% END ALC5

% BEGIN ALC6
\paragraph{ALC6. The induced filtration on a finite submodule.}
\label{claim:ALC6}
For a finite module \(M\) over a Noetherian ring \(E\), a
submodule \(N\subset M\), and an ideal \(I\subset E\),
there is \(c\ge0\) such that
\[
 N\cap I^qM=I^{q-c}(N\cap I^cM)\quad(q\ge c).
                                                        \tag{ALC.8}
\]
In particular the filtrations \(I^qN\) and \(N\cap I^qM\)
are cofinal.

\begin{proof}
The Rees algebra \(\bigoplus_{q\ge0}I^qt^q\) is generated
over \(E\) by finitely many elements \(a_it\), hence is
Noetherian.  The Rees module \(\bigoplus_{q\ge0}I^qMt^q\)
is finite over it, generated by a finite set of generators of
\(M\) in degree zero.  Its graded submodule
\(\bigoplus_{q\ge0}(N\cap I^qM)t^q\) therefore has a
finite homogeneous generating family.  Choose \(c\) at
least their degrees.

A generator of degree \(d\le c\) contributes in degree
\(q\ge c\) an element of
\(I^{q-d}(N\cap I^dM)\).  This lies in
\(I^{q-c}(N\cap I^cM)\).  The opposite inclusion follows
from multiplication inside \(M\), giving (ALC.8).
Now \(I^qN\subset N\cap I^qM\), while
\(N\cap I^{q+c}M\subset I^qN\).  These are the two
cofinality inclusions.  No assumption that \(I\) lies in a
Jacobson radical was used.
\end{proof}
% END ALC6

% BEGIN ALC7
\paragraph{ALC7. Exactness of finite-module completion and its tensor description.}
\label{claim:ALC7}
Completion for an ideal \(I\) of a Noetherian ring \(E\)
is exact on short exact sequences of finite modules.  For
every finite \(E\)-module \(M\), its canonical map gives
\[
               \widehat E\otimes_E M\simeq\widehat M.
                                                        \tag{ALC.9}
\]

\begin{proof}
Consider \(0\to N\to M\to Q\to0\).  ALC6 identifies
the completion of \(N\) in its own filtration with that in
the filtration induced from \(M\), so its map into
\(\widehat M\) is injective.

For surjectivity onto \(\widehat Q\), choose successive
representatives \(q_j\in Q\) with
\(q_{j+1}-q_j\in I^jQ\).  Having lifted \(q_j\) to
\(m_j\), lift the difference to an element of \(I^jM\)
and add it to \(m_j\).  This produces a coherent sequence
\(m_j\) lifting the prescribed sequence.  Thus
\(\widehat M\to\widehat Q\) is surjective.

If \(z\in\widehat M\) maps to zero, represent it by
\(m_j\in M\) modulo \(I^jM\).  Its image in \(Q\)
lies in \(I^jQ\); subtract a lift from \(I^jM\) to obtain
\(n_j\in N\).  The differences of these representatives
belong to \(N\cap I^jM\).  Cofinality in ALC6 makes them
a coherent element of \(\widehat N\), with image \(z\).
This proves exactness also at the middle term.

Choose a finite presentation \(E^s\to E^t\to M\to0\),
which exists by Noetherianity.  Applying the just proved
exactness, and using that completion commutes with finite
free modules, identifies \(\widehat M\) with the cokernel
of the same matrix over \(\widehat E\).  Right exactness
of tensor product identifies \(\widehat E\otimes_E M\)
with that cokernel as well.  These identifications agree on
\(1\otimes m\), and hence give the canonical (ALC.9).
\end{proof}
% END ALC7

% BEGIN ALC8
\paragraph{ALC8. Flatness of completion for the actual algebra structure.}
\label{claim:ALC8}
The completion \(\widehat E\) of every Noetherian ring at
every ideal is flat over \(E\).  More generally, if a
Noetherian ring \(D\), an ideal \(y\), and a local
\(D\)-algebra \(R\) satisfy the original
\(\mathrm{IsCplAt}(y,R)\), then this very \(D\)-algebra
\(R\) is flat.

\begin{proof}
For every finitely generated ideal \(J\subset E\), complete
the injection \(J\hookrightarrow E\).  ALC7 and its
canonical tensor description give an injection
\[
              J\otimes_E\widehat E\longrightarrow\widehat E.
                                                        \tag{ALC.10}
\]
The ideal criterion for flatness states that an \(E\)-module
is flat precisely when these maps are injective for all
finitely generated ideals.  It is the finite-relation form
of tensor exactness: relations between a finite set of
coefficients can be resolved in their generated ideal, and
then successively in a finite free presentation.  Applying
that criterion to (ALC.10) proves the first assertion.

For the second assertion, the exact power-contraction and
density fields of \(\mathrm{IsCplAt}\) give canonical
isomorphisms \(D/y^N\simeq R/\mathfrak m_R^N\) for
every positive \(N\).  They commute with reduction.  The
complete separated ring \(R\) is therefore isomorphic to
\(\varprojlim_ND/y^N\) as a \(D\)-algebra, with its
original map from \(D\).  The first assertion applies to
that completion.  Transport of modules through this
\(D\)-algebra isomorphism proves flatness for the specified
structure on \(R\).
\end{proof}
% END ALC8

% BEGIN ALC9
\paragraph{ALC9. Full prime and minimal-prime return for a completed local chart.}
\label{claim:ALC9}
Let \(D\) be Noetherian, \(y\) maximal, and \(R\) a
local \(D\)-algebra satisfying \(\mathrm{IsCplAt}(y,R)\).
Besides flatness, the following all hold:
\begin{enumerate}
\item every prime \(p\subseteq y\) of \(D\) is the
contraction of a prime of \(R\);
\item every prime of \(R\) contracts into \(y\);
\item every minimal prime of \(R\) contracts to a minimal
prime of \(D\);
\item for every ideal \(J\subset D\), every prime minimal
over \(JR\) contracts to a prime minimal over \(J\).
\end{enumerate}
Together these prove the full original theorem \(\mathrm{c3a}\).

\begin{proof}
ALC8 proves flatness, and the power-contraction identity with
exponent one gives \(\mathfrak m_R\cap D=y\).
We spell out the prime argument so that it applies to the
actual map and to every ideal \(J\).

For a prime \(Q\subset R\), put \(p=Q\cap D\).
The induced map \(D_p\to R_Q\) is flat and local.  It is
faithfully flat: a proper ideal of \(D_p\) extends into the
maximal ideal of \(R_Q\), and hence remains proper.  Any
nonzero module has a nonzero cyclic submodule; flatness
preserves its injection and the preceding proper-ideal
assertion keeps that cyclic module nonzero after tensoring.
If \(p_0\subseteq p\), the fiber of \(R_Q\) over the
prime \(p_0D_p\) is consequently a nonzero ring.  Choose
a prime of that fiber.  Its contraction to \(R\) gives
\(Q_0\subseteq Q\) with \(Q_0\cap D=p_0\).
This proves the needed going-down assertion directly.

Apply it first to \(Q=\mathfrak m_R\) to obtain the first
claim.  The second follows from \(Q\subseteq\mathfrak m_R\).
For the fourth, let \(Q\) be minimal over \(JR\), and
let \(p_0\) be a prime containing \(J\) and contained in
\(Q\cap D\).  Going down produces \(Q_0\subseteq Q\)
over \(p_0\).  Since \(J\subseteq p_0\), this prime
contains \(JR\).  Minimality forces \(Q_0=Q\), hence
\(p_0=Q\cap D\).  This is exactly minimality over \(J\).
The third claim is its case \(J=0\).  No choice of a
different completion or of a different prime above \(y\)
has entered the argument.
\end{proof}
% END ALC9

% BEGIN ALC10
\paragraph{ALC10. Finite-residue contractions and the original coefficient extension.}
\label{claim:ALC10}
If a ring \(D\) maps to a local ring with finite residue
field, the inverse image of its maximal ideal is maximal.
If \(O,O'\) are discrete valuation domains and
\(\mathfrak m_OO'=\mathfrak m_{O'}\), then the specified
map \(O\to O'\) is injective and \(O'\) is flat over \(O\).

\begin{proof}
For the first assertion the image of \(D\) in the finite
residue field is a finite integral domain.  Multiplication
by any nonzero element of that image is injective, hence
surjective, so the image is a field.  Its kernel is the
stated contraction, proving maximality.

For the second choose a uniformizer \(\varpi\) of \(O\).
Its image generates \(\mathfrak m_{O'}\), so is nonzero:
otherwise the discrete valuation ring \(O'\) would be a
field.  A nonzero element of \(O\) is a unit times a power
of \(\varpi\).  Its image is nonzero in the domain \(O'\),
proving injectivity.  Thus \(O'\) is torsion-free over \(O\).
Every nonzero ideal of a DVR is principal, and its tensor
map into \(O'\) is multiplication by a nonzero element.
It is injective by torsion-freeness.  The same ideal
criterion used in ALC8 proves flatness.
\end{proof}
% END ALC10

% BEGIN ALC11
\paragraph{ALC11. Full original flatness and maximality of the arithmetic chart.}
\label{claim:ALC11}
For every original \(x\) and every supplied
\(L:S.\mathrm{Loc}(x)\), the map \(\iota:A\to T\) is
flat, \(x\) is maximal, and
\[
                    \mathfrak m_B\subseteq\pi(x).
                                                        \tag{ALC.11}
\]
In fact equality holds in (ALC.11).  This supplies every
field of the full original \(\mathrm{loc\_\allowbreak flat}\) for
both allowed ranks.

\begin{proof}
Use the canonical \(O'\)-algebra map
\(\iota':O'\otimes_OA\to T\), sending
\(b\otimes a\) to \(b\iota(a)\).  ALC4 identifies it
as the original completion map, and ALC8 proves it flat.
ALC10 and the supplied equality
\(\mathfrak m_OO'=\mathfrak m_{O'}\) show that \(O'\)
is flat over \(O\).  Tensoring therefore makes
\(A\to O'\otimes_OA\), \(a\mapsto1\otimes a\),
flat.  The composite is exactly \(\iota\), so it is flat.

The finite residue field of \(T\) is the specified residue
field of \(O'\).  Applying the first part of ALC10 to
\(\iota\) shows that its maximal-ideal contraction \(x\)
is maximal.  Apply the same assertion to the composite
\(B\to A\to T\).  Its contraction \(\pi(x)\) is a
maximal ideal of the actual local ring \(B\).  RTC7 says
that this ring's maximal ideal is exactly \(\mathfrak m_B\).
Thus \(\pi(x)=\mathfrak m_B\), proving (ALC.11).
\end{proof}
% END ALC11

\subsection{The residual types in ranks two and three}

% BEGIN ALC12
\paragraph{ALC12. Opposite triangular directions cannot both survive on the special fiber.}
\label{claim:ALC12}
Let \(K\) be any field and \(f:A\to K\) any ring map,
with no prescribed coefficient-field structure on \(K\).
If \(\mathfrak m_B\subseteq\pi(\ker f)\), then
\[
             \mathfrak t_+\subseteq\ker f
             \quad\hbox{or}\quad
             \mathfrak t_-\subseteq\ker f.
                                                        \tag{ALC.12}
\]
This holds for both original ranks, including rank three.

\begin{proof}
Suppose neither inclusion holds.  Since these are generated
entry ideals, choose a below-diagonal entry \(a=u(g)_{ij}\)
and an above-diagonal entry \(b=u(h)_{rs}\) with
\(f(a)\ne0\) and \(f(b)\ne0\).  In both allowed ranks
the weights strictly decrease with the index.  Thus the
weight of \(a\) is \(-p\), and that of \(b\) is \(q\),
for positive integers \(p,q\).

The product \(a^qb^p\) has weight zero.  By the definition
of the actual coefficient algebra it belongs to \(B\).
Both entries are off the diagonal, so their images in the
residue field of the original \(\mathcal R\) are zero.
Consequently \(a^qb^p\in\mathfrak m_B\).  The assumed
contraction condition gives
\[
                       f(a)^q f(b)^p=0.
\]
This is impossible in a field when both factors are nonzero.
The argument uses an actual weight-zero entry monomial; it
does not use the rank-two triangular point realization, nor
does it require a finite or algebraic residue-field target.
\end{proof}
% END ALC12

% BEGIN ALC13
\paragraph{ALC13. The marked difference gives every coordinate projector.}
\label{claim:ALC13}
The actual \(B\)-linear span of \(u(G_0)\) contains each
coordinate projector \(E_{ii}\).  More generally the full
original \(\mathrm{idem}\) statement holds: for any original
coefficient algebra \(C\), any commutative \(C\)-algebra
\(K\), and any normalized \(\rho\) with
\(\mathrm{Integral}(O,C,\rho)\), the \(C\)-span of
\(\rho(G_0)\) contains all \(E_{ii}\).

\begin{proof}
Here is the elementary matrix calculation behind both assertions.
The linear span of a multiplicative matrix representation is
an algebra: products of two generators are again generators,
and its identity is the matrix at the group identity.
Suppose \(2\) and \(\zeta\) are units in the scalar ring
and that this algebra contains
\[
 \rho(g_0)-\rho(g_0^*)
             =\zeta\operatorname{diag}(w(n)_i).
                                                        \tag{ALC.13}
\]
Multiplying by \(\zeta^{-1}\) puts
\(Y=\operatorname{diag}(w(n)_i)\) in the algebra.  For
rank two its projectors are
\[
                  E_{00}=(1+Y)/2,
                  \qquad E_{11}=(1-Y)/2.
                                                        \tag{ALC.14}
\]
For rank three they are
\[
 E_{00}=(Y^2+Y)/2,
 \qquad E_{11}=1-Y^2,
 \qquad E_{22}=(Y^2-Y)/2.
                                                        \tag{ALC.15}
\]
Each formula is checked on the diagonal entries \(1,-1\)
or \(1,0,-1\); off-diagonal entries vanish.  In particular
no division by \(3\) occurs.

For the original \(\mathrm{idem}\) statement, the
\(\mathrm{Integral.ze}\) field supplies \(\zeta\in C\)
whose residue is \(\Psi(g_0)-\Psi(g_0)^{-1}\), together
with the diagonal entries of (ALC.13).  The off-diagonal
entries vanish by the normalized \(\mathrm{diag}\) condition.
The residue of \(\zeta\) is nonzero because
\(\Psi(g_0)^2\ne1\), the original marked-element field
of \(F.\mathrm{Ok}\).  Thus \(\zeta\) is a unit in
the local ring \(C\).  Since \(\ell\ne2\), the residue
of \(2\) is also nonzero, so \(2\) is a unit.  Applying
(ALC.14) or (ALC.15) proves the full claimed span statement
over every original \(C,K\).  The additional trace and
ordinary-character fields of \(\mathrm{Integral}\) need
no strengthening.

For the actual representation over \(A\), a full
\(\mathrm{Integral}\) record need not be imported.  Set
\(\zeta=u(g_0)_{00}-u(g_0)_{n-1,n-1}\in B\).
Its residue is the same nonzero difference.  RTC7 makes
\(B\) a local coefficient ring, so \(2,\zeta\) are units
there.  Polarization gives
\(u(g_0^*)_{ij}=u(g_0)_{j^{\rm rev},i^{\rm rev}}\).
Together with the original diagonal normalization, this says
that the marked difference has diagonal \((\zeta,-\zeta)\)
or \((\zeta,0,-\zeta)\).  This is (ALC.13), and the same
formulas prove the actual \(B\)-span assertion directly.
\end{proof}
% END ALC13

% BEGIN ALC14
\paragraph{ALC14. Coordinate-stable subspaces and the rank-three connectivity cases.}
\label{claim:ALC14}
Let \(a:G\to\operatorname{Mat}_n(K)\) be a representation
over a field.  Suppose every \(a(G)\)-stable subspace is
stable under each \(E_{ii}\).  Put an arrow \(j\to i\)
when \(a(g)_{ij}\ne0\) for some \(g\).  If this directed
graph is strongly connected, the representation is irreducible.
For rank two, nonzero entries in both triangular directions
give this connectivity.  For rank three, it follows from
both triangular directions, at least one odd-weight entry,
and the reflection symmetry
\[
 \bigl(\exists g:\ a(g)_{ij}\ne0\bigr)
 \Longleftrightarrow
 \bigl(\exists g:\ a(g)_{j^{\rm rev},i^{\rm rev}}\ne0\bigr).
                                                        \tag{ALC.16}
\]

\begin{proof}
Let \(W\ne0\) be stable.  A nonzero vector in \(W\) has
a nonzero coordinate \(v_j\); applying \(E_{jj}\) and
dividing by \(v_j\) puts \(e_j\) in \(W\).  If \(e_j\)
belongs to \(W\) and \(a(g)_{ij}\ne0\), applying
\(E_{ii}a(g)\) and dividing by that entry puts \(e_i\)
in \(W\).  Strong connectivity therefore puts every basis
vector in \(W\), proving \(W=K^n\).

In rank two the two nonzero triangular directions are exactly
the arrows \(0\to1\) and \(1\to0\).
In rank three, reflection pairs the adjacent lower arrows
\(0\to1\), \(1\to2\), and pairs the adjacent upper
arrows \(2\to1\), \(1\to0\).  Call the existence of
the lower pair \(L\), and of the upper pair \(U\).
The odd-weight condition is exactly \(L\) or \(U\).
If both hold, both directions along the adjacent chain exist.
If \(L\) holds but \(U\) does not, the nonzero upper
triangular direction must be the long arrow \(2\to0\).
Together these form the directed cycle
\(0\to1\to2\to0\).  If \(U\) holds but \(L\) does
not, the lower direction must be \(0\to2\), giving the
opposite directed cycle.  Each case is strongly connected.
These cases exhaust the original hypotheses.
\end{proof}
% END ALC14

% BEGIN ALC15
\paragraph{ALC15. Full original classification of the field-valued types.}
\label{claim:ALC15}
For every field \(K\) and every ring map \(f:A\to K\),
\[
 \operatorname{MatIrr}\bigl(u\otimes f\bigr)
       \Longleftrightarrow \mathrm{TI}(\ker f).
                                                        \tag{ALC.17}
\]
Moreover, if \(\mathfrak m_B\subseteq\pi(\ker f)\),
one of the two triangular entry ideals lies in \(\ker f\).
These are both clauses of the full original
\(\mathrm{types}\) theorem, for ranks two and three and
every original field-valued point.

\begin{proof}
The second assertion is ALC12.  For (ALC.17), an entry ideal
is contained in \(\ker f\) precisely when every one of
its generating entries vanishes after applying \(f\).
If the lower entries all vanish, the first coordinate line
is stable; if the upper entries all vanish, the last
coordinate line is stable.  These are nonzero proper lines
because \(n\ge2\).  In rank three, if the odd-weight
entries all vanish, the middle coordinate line is stable.
Indeed the two entries in its column outside the middle
position have odd weight.  Irreducibility therefore implies
all conditions in (ALC.2).

Conversely suppose (ALC.2) holds.  By ALC13 each coordinate
projector belongs to the actual \(B\)-span of \(u(G_0)\).
After applying \(f\), a subspace stable under the represented
group is stable under those projectors.  Polarization gives
the reflection symmetry (ALC.16): the entry at the reflected
position is the original entry evaluated at \(g^*\).
Applying the same identity twice gives the converse vanishing
implication as well.  The noncontainments in (ALC.2) say
exactly that the graph has the arrows required by ALC14.
That claim proves irreducibility.

No topology, cardinality, or prescribed \(k\)-algebra
structure was imposed on \(K\).  The units used to form
the projectors were already units in \(B\), and their
images remain units under the arbitrary ring map \(f\).
Thus the argument retains the original broad point scope.
It also applies without change when \(\ell=3\).
\end{proof}
% END ALC15

% BEGIN ALC16
\paragraph{ALC16. The continuity, residue and entry-ideal dictionaries.}
\label{claim:ALC16}
For a supplied original local object \(L\), the representation
\(\rho_L\) is adically continuous for \(\mathfrak m_T\).
For every field map \(h:k\to\operatorname{Res}(T)\)
satisfying the original \(\mathrm{IsResMap}(O,T,h)\), its
reduction has the exact \(\mathrm{IsBar}(h)\) diagonal and
ordinary characters.  For every ideal \(P\subset T\),
\[
 \begin{aligned}
 \rho_L.\mathfrak t_+\subseteq P
       &\Longleftrightarrow \mathfrak t_+\subseteq\iota^{-1}P,\\
 \rho_L.\mathfrak t_-\subseteq P
       &\Longleftrightarrow \mathfrak t_-\subseteq\iota^{-1}P,\\
 \rho_L.\mathfrak e\subseteq P
       &\Longleftrightarrow \mathfrak e\subseteq\iota^{-1}P.
 \end{aligned}                                             \tag{ALC.18}
\]
No primality hypothesis on \(P\) is needed.

\begin{proof}
ALC11 gives \(\mathfrak m_B\subseteq\pi(x)\), and the
original \(L.\mathrm{lies}\) gives
\(x=\iota^{-1}\mathfrak m_T\).  Therefore
\(\iota(x_0^N)\subseteq\mathfrak m_T^N\) for every
\(N\).  Compose each locally constant quotient map in
RTC9 with the induced quotient map to
\(T/\mathfrak m_T^N\).  This proves continuity of every
matrix entry and every ordinary-character value, with all
original local group arguments and all powers \(N\).

For the exact residue assertion, let \(b\in B\), and
choose \(a\in O\) lifting its residue in the original
\(\mathcal R\).  Then \(b-a\in\mathfrak m_B\), so
\(\iota(b-a)\in\mathfrak m_T\).  The prescribed
compatibility of \(h\) with \(O\) gives
\[
 \operatorname{res}_T(\iota(b))
       =h\bigl(\operatorname{res}_{\mathcal R}(b)\bigr).
                                                        \tag{ALC.19}
\]
Every diagonal entry and every ordinary-character value is
an actual element of \(B\).  Their residues in the supplied
space are, respectively, \(\Psi(g)^{w(n)_i}\) and
\(\Psi(\operatorname{loc}_v s)^{w(n)_i}\).
Applying (ALC.19) is exactly the two-part
\(\mathrm{IsBar}(h)\) assertion.  No equality of
semisimplifications is substituted for these literal entries.

Finally, an ideal generated by the \(\iota\)-images of
a family lies in \(P\) if and only if every original
generator lies in \(\iota^{-1}P\).  Apply this statement
to each of the three defining entry predicates.  It gives
(ALC.18) for every ideal, including the zero and unit ideals.
\end{proof}
% END ALC16

% BEGIN ALC17
\paragraph{ALC17. Full original local representation theorem.}
\label{claim:ALC17}
For every original \(L:S.\mathrm{Loc}(x)\), all clauses
of \(\mathrm{loc\_\allowbreak rho}\) hold.  In addition to the
continuity, every specified residue map, and all ideal
comparisons in ALC16, its final assertion is
\[
 \mathfrak t_+\not\subseteq x
 \quad\Longrightarrow\quad
 \rho_L.\mathfrak t_+=T,
 \qquad \rho_L.\mathfrak t_-\subseteq\mathfrak m_T.
                                                        \tag{ALC.20}
\]
The statement holds for both allowed ranks on the same
original \(F,S,L\).

\begin{proof}
All preceding fields are ALC16.  If
\(\rho_L.\mathfrak t_+\) were a proper ideal of the local
ring \(T\), it would lie in \(\mathfrak m_T\).
The first equivalence in (ALC.18) and
\(\iota^{-1}\mathfrak m_T=x\) would give
\(\mathfrak t_+\subseteq x\), a contradiction.
Thus \(\rho_L.\mathfrak t_+=T\).

Consider the actual field-valued map
\(A\xrightarrow{\iota}T\to\operatorname{Res}(T)\).
Its kernel is \(x\), and ALC11 gives the required
\(\mathfrak m_B\) contraction.  The special-fiber clause
of the full types theorem, ALC15, says that one triangular
entry ideal is contained in \(x\).  The hypothesis excludes
the lower one, so \(\mathfrak t_-\subseteq x\).
The second equivalence in (ALC.18) gives the remaining
assertion in (ALC.20).  In rank three this uses the proved
general opposite-weight product argument, not a rank-two
point-realization shortcut.
\end{proof}
% END ALC17

\subsection{Existence at every original closed point}

% BEGIN ALC18
\paragraph{ALC18. The original closed point has a finite residue field.}
\label{claim:ALC18}
If \(x\subset A\) is maximal and
\(\mathfrak m_B\subseteq\pi(x)\), then \(K=A/x\)
is a finite field with its prescribed \(O\)-algebra
structure, and \(\mathfrak m_O\) acts by zero on it.

\begin{proof}
RTC7 makes \(A\) a finite-type algebra over the local
coefficient ring \(B\), whose residue field is the original
finite \(k\).  The contraction hypothesis makes \(K\)
a finite-type \(k\)-algebra, and maximality makes it a
field.  We recall the finite-type field argument to make
the residue extension explicit.

For a field generated as a \(k\)-algebra by finitely many
elements, choose a maximal algebraically independent subset
\(t_1,\ldots,t_d\) of these generators.  The remaining
generators are algebraic over \(k(t_1,\ldots,t_d)\).
Clearing the finitely many denominators of their monic
equations makes them integral over
\(k[t_1,\ldots,t_d,h^{-1}]\) for one nonzero polynomial
\(h\).  The whole generated field is a finite integral
algebra over this ring, since it already contains \(h^{-1}\).
A subring over which a field is integral is itself a field:
the monic equation for the inverse of a nonzero element
expresses that inverse in the subring.

If \(d>0\), choose an irreducible polynomial in \(k[t_1]\)
which does not divide \(h\).  Such a polynomial exists
because a one-variable polynomial ring over a field has
infinitely many irreducibles, by the usual product-plus-one
argument.  It remains prime after adjoining the other
variables and is not made invertible by inverting \(h\).
Thus \(k[t_1,\ldots,t_d,h^{-1}]\) is not a field, a
contradiction.  Hence \(d=0\), so the original field is
a finite algebraic extension of \(k\), and is finite as
a set because \(k\) is finite.

Apply this to \(K=A/x\).  The local map \(O\to B\)
sends \(\mathfrak m_O\) into \(\mathfrak m_B\), whose
image in \(K\) is zero.  The resulting \(k\)-embedding
is the one induced by the original \(O\)-structure; no
replacement of that embedding has been made.
\end{proof}
% END ALC18

% BEGIN ALC19
\paragraph{ALC19. Constructing the required finite unramified coefficient extension.}
\label{claim:ALC19}
Let \(O\) be any original coefficient DVR and let \(K\)
be a finite field with an \(O\)-algebra structure killing
\(\mathfrak m_O\).  There is an explicitly constructed
finite free \(O\)-algebra \(O'\), itself a coefficient
DVR of residue characteristic \(\ell\), such that
\[
 \mathfrak m_OO'=\mathfrak m_{O'},\qquad
              K\xrightarrow{\ \sim\ }\operatorname{Res}(O')
                  \quad\hbox{as \(O\)-algebras}.
                                                        \tag{ALC.21}
\]
This supplies the full original \(\mathrm{unr\_\allowbreak ex}\)
output, with an isomorphism in place of its required embedding.

\begin{proof}
The given map factors through an embedding of the finite
field \(k=\operatorname{Res}(O)\).  Choose a generator
\(\gamma\) of the cyclic group \(K^\times\); it also
generates \(K\) as a field over \(k\).  For completeness,
cyclicity follows from the elementary finite-subgroup argument:
in a finite abelian subgroup of a field's units, an element
has order equal to the group exponent, by combining its
prime-primary maximal orders.  The polynomial \(X^e-1\)
has at most \(e\) roots, so this exponent is at least the
group order and therefore equal to it.

Let \(g\in k[X]\) be the monic minimal polynomial of
\(\gamma\).  Lift its coefficients to \(O\), retaining
the leading coefficient one, to obtain a monic polynomial
\(G\) of the same positive degree.  Define
\[
                    O'=O[X]/(G).
                                                        \tag{ALC.22}
\]
Division by the monic polynomial gives the finite free
\(O\)-basis \(1,X,\ldots,X^{\deg G-1}\).  Let
\(\varpi\) be a uniformizer of \(O\).  This basis shows
that \(O'\) is complete and separated for \(\varpi O'\),
and that multiplication by \(\varpi\) is injective.  Its
quotient is exactly
\[
          O'/\varpi O'=k[X]/(g)\simeq K.
                                                        \tag{ALC.23}
\]
An element outside \(\varpi O'\) has an inverse modulo
that ideal.  The convergent geometric series corrects it
to an actual inverse.  Thus \(O'\) is local with maximal
ideal \(\varpi O'\).

Every nonzero element \(b\in O'\) has a finite maximal
power \(\varpi^r\) dividing it: in the displayed free
basis, take the minimum valuation of its nonzero
coefficients.  Thus \(b=\varpi^r u\) with \(u\) a
unit.  Since \(\varpi\) is a non-zero-divisor, products
of nonzero elements are nonzero.  The ring is therefore
a domain.  In any nonzero ideal, choose an element of
smallest such exponent \(r\); the ideal is then exactly
\(\varpi^rO'\).  This proves the discrete valuation
property and Noetherianity directly.  The uniformizer is
nonzero and not a unit, so the ring is not a field.

Its maximal-ideal completeness is the already proved
\(\varpi\)-adic completeness.  Its residue is finite and
has characteristic \(\ell\), proving all coefficient
clauses.  The image of \(\mathfrak m_O=(\varpi)\) is
exactly the maximal ideal.  Finally the inverse of the
specific isomorphism in (ALC.23), sending the polynomial
class to \(\gamma\), is the map in (ALC.21).  The
coefficient lift and this map respect the prescribed
\(O\)-action throughout.
\end{proof}
% END ALC19

% BEGIN ALC20
\paragraph{ALC20. The residue-compatible coefficient completion and its universal property.}
\label{claim:ALC20}
Let \(O\) be any commutative local ring, \(D\) a
Noetherian \(O\)-algebra, and
\(\varphi:D\to\operatorname{Res}(O)\) an
\(O\)-algebra map.  There are one coefficient algebra
\(R\) and one map \(j:D\to R\) such that
\[
          \operatorname{res}_R(j(d))=\varphi(d),
                                                        \tag{ALC.24}
\]
and every residue-compatible map from \(D\) to a coefficient
algebra factors uniquely through \(j\).  The same conclusion
holds for any already supplied \(\mathrm{IsCplAt}(\ker\varphi,R)\)
with its given \(O\)-structure and scalar tower.  These
are the full original \(\mathrm{cpl\_\allowbreak ex}\) and
\(\mathrm{cpl\_\allowbreak univ}\) statements.

Independently of a chosen coefficient ring \(O\), let
\(D\) be any Noetherian ring, \(y\) any maximal ideal,
and \(R\) any local \(D\)-algebra satisfying
\(\mathrm{IsCplAt}(y,R)\).  For every commutative \(T\),
every ideal \(J\) for which \(T\) is complete and
separated, and every \(f:D\to T\) with \(f(y)\subseteq J\),
there is a unique extension sending \(\mathfrak m_R\)
into \(J\); moreover
\(R\simeq_D\mathrm{Cpl}(y)\) and \(yR=\mathfrak m_R\).
These are all outputs of the full original
\(\mathrm{cpl\_\allowbreak uniq}\).

\begin{proof}
The map \(\varphi\) is surjective on the image of \(O\),
so \(y=\ker\varphi\) is maximal.  Construct
\(R=\varprojlim_N D/y^N\) and the canonical \(j\).
ALC5 proves that \(R\) is Noetherian local and complete,
with maximal ideal \(yR\) and residue \(D/y\).
The map induced by \(\varphi\) identifies that residue
with the original \(k\).  The composite \(O\to R\)
is local and surjective on residues: a nonunit has residue
zero, and every residue is lifted from \(O\).  Thus it
is the required coefficient algebra and (ALC.24) holds.

If \(f:D\to C\) has the stated residue, then
\(f(y^N)\subseteq\mathfrak m_C^N\) for every \(N\).
The compatible maps
\(D/y^N\to C/\mathfrak m_C^N\), followed by completeness
of \(C\), define an \(O\)-algebra map \(g:R\to C\).
It extends \(f\).  Any other such \(O\)-algebra map is
local and continuous by ALC1.  Since \(j(D)\) is dense,
the two maps agree modulo every \(\mathfrak m_C^N\),
and separatedness gives equality.  The objects \(R,j\)
were fixed before \(C,f\), as the original contract requires.

For an already supplied \(R\) with
\(\mathrm{IsCplAt}(y,R)\), its exact contraction and
density give \(R/\mathfrak m_R\simeq D/y\), respecting
the given scalar tower.  This makes the structure map
\(O\to R\) local with the required residue field, so
\(R\) is a coefficient algebra.  Its compatible finite
quotients identify it, with the original map from \(D\),
with the same inverse limit.  The preceding construction
and uniqueness therefore apply to that ring and that map,
proving the wider \(\mathrm{cpl\_\allowbreak univ}\) assertion too.

For the coefficient-independent assertion, the same exact
quotient and density argument identifies the given \(R\)
with \(\varprojlim D/y^N\).  No local-ring or coefficient
structure on \(T\) is needed.  The maps
\(D/y^N\to T/J^N\) define the extension by \(J\)-adic
completeness.  The identity \(yR=\mathfrak m_R\) follows
from this inverse-limit identification and (ALC.6); it
shows that the extension sends \(\mathfrak m_R\) into
\(J\).  Any two extensions with that property agree
modulo \(J^N\), by approximating each \(r\in R\) with
an element of \(D\) modulo \(\mathfrak m_R^N\).
Separatedness gives uniqueness.  Finally ALC5 identifies
this inverse limit with the actual \(\mathrm{Cpl}(y)\),
as a \(D\)-algebra.  This proves the claimed full
\(\mathrm{cpl\_\allowbreak uniq}\) scope.
\end{proof}
% END ALC20

% BEGIN ALC21
\paragraph{ALC21. Full original existence of the arithmetic local object.}
\label{claim:ALC21}
For every original maximal ideal \(x\subset A\) with
\(\mathfrak m_B\subseteq\pi(x)\), an original
\(L:S.\mathrm{Loc}(x)\) exists.  This is the full
\(\mathrm{loc\_\allowbreak ex}\) theorem at both allowed ranks.

\begin{proof}
By ALC18, \(K=A/x\) is a finite field with the required
\(O\)-structure.  ALC19 constructs a finite coefficient
DVR extension \(O'\), with
\(\mathfrak m_OO'=\mathfrak m_{O'}\), and an
\(O\)-algebra embedding \(e:K\to k'\), where
\(k'=\operatorname{Res}(O')\).  Keep this choice fixed.
Put \(D=O'\otimes_O A\), a Noetherian ring because it
is finite over the Noetherian ring \(A\).  Define
\[
 \varphi:D\longrightarrow k',\qquad
       b\otimes a\longmapsto\operatorname{res}_{O'}(b)
                                  e(a\bmod x).
                                                        \tag{ALC.25}
\]
The prescribed scalar compatibility makes this a well-defined
\(O'\)-algebra map.  ALC20 constructs its coefficient
completion \(j:D\to T\).  Give \(T\) the composite
\(O\)-structure and put
\[
                     \iota(a)=j(1\otimes a).
                                                        \tag{ALC.26}
\]
Then
\(\operatorname{res}_T(\iota(a))=e(a\bmod x)\).
Injectivity of \(e\) gives
\(\iota^{-1}\mathfrak m_T=x\), the literal
\(\mathrm{lies}\) field.

It remains to prove the exact \(\mathrm{univ}\) field,
including its scalar structures.  Let \(C\) be any
coefficient algebra over \(O'\), with the supplied
\(O\)-structure and scalar tower, and let \(f:A\to C\)
be an \(O\)-algebra map whose residue agrees with
\(\iota\).  The tensor map
\[
               f':D\to C,\qquad b\otimes a\mapsto b f(a)
\]
is \(O'\)-linear.  Its residue is (ALC.25), first on
pure tensors by the assumed equality for \(a\), then
on their sums.  The universal property in ALC20 supplies
a unique \(O'\)-algebra map \(g:T\to C\) with
\(gj=f'\).  Restriction to \(1\otimes A\) gives
\(g\iota=f\).

If another \(O'\)-algebra map has this equality on \(A\),
then its composite with \(j\) agrees with \(f'\) on
every pure tensor, and hence on \(D\).  Uniqueness in
ALC20 forces it to equal \(g\).  This proves exactly the
original local universal property.  The finite coefficient
extension, its unramified maximal-ideal identity, all
coefficient-algebra fields of \(T\), both scalar structures
and their tower, \(\iota\), \(\mathrm{lies}\), and
\(\mathrm{univ}\) have now all been constructed on the
same object.
\end{proof}
% END ALC21

% BEGIN ALC22
\paragraph{ALC22. The actual local component charts are now available to propagation.}
\label{claim:ALC22}
For every original closed point \(x\) over
\(\mathfrak m_B\), the local object constructed in ALC21
has all the completion, flatness and representation properties
proved above.  The actual spectral map
\[
       f_L:\operatorname{Spec}T\longrightarrow\operatorname{Spec}A,
                     \qquad P\longmapsto\iota^{-1}P
                                                        \tag{ALC.27}
\]
is the full component chart required by the original
\(\mathrm{ChartOfFlat}\) consumer: it is continuous,
sends the closed point to \(x\), sends minimal primes to
generic points of irreducible components, and realizes every
irreducible component through \(x\) in this way.  For a
core point, its local lower-entry ideal is the unit ideal
and its upper-entry ideal lies in \(\mathfrak m_T\).

\begin{proof}
The construction supplies an actual \(L\), so ALC4,
ALC11 and ALC17 apply without another existence premise.
In particular \(A\to T\) is flat and
\(\iota^{-1}\mathfrak m_T=x\).  Apply the already
proved ACP16 with the identity open chart of
\(\operatorname{Spec}A\).  Its assumptions are exactly
these properties, and its map is (ALC.27), with no change
of the underlying algebra structure.  This gives every
component-chart clause.  Alternatively the prime and
minimal-prime argument in ALC9 is the same going-down
mechanism after the coefficient tensor extension.

For a point of the original core, the extra condition is
\(\mathfrak t_+\not\subseteq x\).  ALC17 gives the
two asserted local triangular ideals.  The ideal comparisons
of ALC16 identify the local \(\mathrm{TI}\), exceptional
entry and good loci with the inverse images of their
original global conditions.  In particular, for an original
ideal \(M\subset B\),
\[
          M\subseteq L.\pi(P)
             \Longleftrightarrow
          M\subseteq\pi(\iota^{-1}P),
                                                        \tag{ALC.28}
\]
and the right side is the original global good condition.
Thus the charts fit the actual local-to-global comparison
used in \(\mathrm{ComponentsGoodRankTwo}\), and its
rank-three counterpart, with the same rings and ideals.

This does not supply the remaining numerical and arithmetic
fields of their propagation records.  The sharp component
chains, dimension bounds, nice-prime hypotheses, global
density and seed incidence must still be supplied.  The
conclusion here is the completed original local producer,
not the whole propagation theorem.
\end{proof}
% END ALC22

\paragraph{Scope of the completed producer.}
The full original \(\mathrm{cpl\_\allowbreak of\_\allowbreak univ}\),
\(\mathrm{loc\_\allowbreak cpl}\), \(\mathrm{c3a}\),
\(\mathrm{cpl\_\allowbreak is}\), \(\mathrm{cpl\_\allowbreak uniq}\),
\(\mathrm{loc\_\allowbreak flat}\), \(\mathrm{idem}\),
\(\mathrm{types}\), \(\mathrm{loc\_\allowbreak rho}\),
\(\mathrm{unr\_\allowbreak ex}\), \(\mathrm{cpl\_\allowbreak univ}\),
\(\mathrm{cpl\_\allowbreak ex}\), and \(\mathrm{loc\_\allowbreak ex}\)
are supplied at the scopes stated above.  The elementary
completion lemmas allow arbitrary Noetherian base rings;
the coefficient-completion universal lemmas allow arbitrary
commutative local coefficient rings.  The arithmetic
theorems retain the original DVR, \(F.\mathrm{Ok}\),
rank-two-or-three, and supplied \(\mathrm{Space}\) premises.
No finite-point conclusion has been used upstream.

Rank-two and rank-three spaces and their coefficient ideals
remain separate.  Their local coefficient extensions need
not be identified.  All constructions keep the original
frame, group maps, ordinary characters and marked element;
the original arithmetic \(P\) premise remains where the
AFE entry uses it.  The allowed Host case \(\ell=3\) is
preserved.  Existence of the supplied universal spaces,
sharp dimension and component bounds, closed-point and
seed incidence, the rank-three corner and exceptional
ideal, full propagation, \(\mathrm{OffFin}\), automorphic
seeds and patching/classicality remain separate.  The
stronger Alt outputs still retain their original
\(\mathrm{WitA}\) and \(\mathrm{WitAQ}\) producers.
No whole selected RH/GRH obligation is closed by this chapter.

\subsection{Exact source correspondence}

All primary files are at the accepted commit
\texttt{7cbf7a0c32fca1badf05931879a37c3f8a2b03b4}.
Write \(F\) for
\texttt{src/\allowbreak Development/\allowbreak Automorphic\allowbreak Arithmetic/\allowbreak }
\texttt{Automorphic\allowbreak Varieties/\allowbreak Frobenius\allowbreak Smooth\allowbreak Vectors/\allowbreak }
and \(H\) for the sibling
\texttt{Commutative\allowbreak Algebra\allowbreak Height/\allowbreak } directory.
Full snapshots, exact ranges and hashes are in this component's
source manifest.

\begin{itemize}
\item ALC1--ALC4 match \(F\)\texttt{Loc\allowbreak Cpl.\allowbreak lean},
lines 7--188, especially the complete universal theorem at
100--161 and actual arithmetic application at 162--188.
The original \(\mathrm{IsCO}\) record is in the sibling
\texttt{Commutative\allowbreak Algebra/\allowbreak Res\allowbreak Map\allowbreak Bijective.\allowbreak lean}, lines 9--13;
\(\mathrm{IsCplAt}\) is in \(H\)\texttt{Is\allowbreak Local\allowbreak Hom.\allowbreak lean}.
\item ALC5--ALC9 reconstruct the completion and flatness
mathematics used by \(H\)\texttt{Cpl\allowbreak Is.\allowbreak lean},
\texttt{Cpl\allowbreak Uniq.\allowbreak lean}, and
\texttt{Highest\allowbreak Weight\allowbreak Space\allowbreak Finite.\allowbreak lean}; the last contains
the full \(\mathrm{c3a}\) statement and its prime clauses.
\item ALC10--ALC11 give all of \(F\)\texttt{Loc\allowbreak Flat.\allowbreak lean},
lines 9--69.  ALC12--ALC15 give the full \(\mathrm{types}\)
statement and its supporting matrix arguments in
\(F\)\texttt{Irr\allowbreak Of\allowbreak Conn.\allowbreak lean}, with the original projector
theorem in \texttt{Add\allowbreak Mul\allowbreak Self\allowbreak Eq\allowbreak Of\allowbreak Orthogonal\allowbreak Idempotents.\allowbreak lean}.
\item ALC16--ALC17 give all outputs of
\(F\)\texttt{Loc\allowbreak Rho.\allowbreak lean}, lines 8--66, including every
ideal, every specified residue map, and both ranks in its
last implication.
\item ALC18--ALC21 reconstruct the source suppliers and
full existence theorem in \(F\)\texttt{Res\allowbreak Eq\allowbreak Of.\allowbreak lean},
lines 11--178.  The monic-polynomial construction in ALC19
proves the needed unramified coefficient extension directly.
\item ALC22 applies the unchanged ACP16 to the actual map
used in the sibling
\texttt{Curve\allowbreak DVRSpecialization/\allowbreak Components\allowbreak Good\allowbreak Rank\allowbreak Two.\allowbreak lean},
lines 107--128 and 149--158.  The remaining dimension,
density and seed inputs in that source are retained explicitly.
\end{itemize}


\section{Stabilizer volumes from the original smooth Siegel--Weil identity}

The purpose of this chapter is to compare two \emph{arithmetic quotient}
volumes.  AOS computed the archimedean orbital scalar without changing
the supplied measures.  Here the additional input is the source's
smooth Siegel--Weil identity.  We identify its continuous local densities
with the already constructed Gram-fiber measures, unfold one rational
orbit, and cancel one positive integral.  The stabilizer quotient is
proved to have finite mass during this argument; its finite mass is
not an extra premise.

The endpoint is the full original \texttt{Vol\allowbreak Stab\allowbreak From\allowbreak SWS}, including
every character, Haar measure, Gram-fiber measure and fundamental-domain
witness in \texttt{Vol\allowbreak Stab\allowbreak At}.  The two retained mathematical premises
are exactly \texttt{SWa1} and
\texttt{Stabilizer\allowbreak First\allowbreak Chern}\allowbreak\texttt{Compatibility}.  Their content is recorded
below.  In particular neither is renamed as the desired volume equality.
The final application supplies the literal \texttt{In\allowbreak Volumes\allowbreak At}
implication, and then its actual \texttt{In\allowbreak Volumes} conclusion on the
GMR19 rescaled setting.  The latter uses the unipotent normalization
already present in GMR19; there is no further rescaling.

\subsection{The common objects and the retained premises}

Let $E/F$ be the original CM extension of number fields, with involution
$\sigma$, and let $J$ be a nondegenerate $\sigma$-Hermitian matrix of
size $M$.  Write
\[
 V=E^M,\qquad G=U(J)(\mathbb A_F),\qquad \Gamma=U(J)(F).
 \tag{SVW.1}
\]
The notation $U(J)(F)$ means the rational unitary matrices with entries
in $E$; their embedding into $G$ is the source's
\texttt{HFin\allowbreak In}.  For an integer $a$ and an ordered rational tuple
$X_0=(x_1,\ldots,x_a)$, assume
\[
 1\leq a,\qquad 2a<M,\qquad T=X_0^*JX_0,\qquad \det T\ne0.
 \tag{SVW.2}
\]
Keep any supplied proof of the Hermitian symmetry of $T$.  Set
\[
 W=(\mathbb A_F\otimes_F E)^{M\times a},\quad
 Y_T=\{X:X^*JX=T\},\quad H=\operatorname{Stab}_G(X_0),\quad
 \Delta=\Gamma\cap H.
 \tag{SVW.3}
\]
Here $X_0$ is embedded diagonally, and $H$ is its pointwise stabilizer.
The source writes tuples as functions rather than column matrices;
these two descriptions have the same ordered columns.

An admissible test is a pure product
\[
 \Phi(X)=\Phi_\infty(X_\infty)\prod_v\Phi_v(X_v),
 \tag{SVW.4}
\]
where $\Phi_\infty$ is continuous with compact support, each finite-place
$\Phi_v$ is locally constant with compact support, and outside a finite
set $\Phi_v$ is the indicator of the tuple lattice obtained from one
specified $F$-basis of $E$.  Formula (SVW.4), including the finite-adelic
function between its two factors, is exactly \texttt{Is\allowbreak Adm}.  The
predicate \texttt{Is\allowbreak Adm\allowbreak S} additionally requires the infinite factor to
be smooth in \emph{every} continuous additive real coordinate chart.
We prove below that the products are well-defined, continuous and
compactly supported in the adelic topology.

For $\Phi$ as above put
\[
 \Theta_{T,\Phi}(g)=
  \sum_{X\in V^a,\ X^*JX=T}\Phi(g^{-1}X),\qquad
 \operatorname{Mean}_{\mu,\mathcal F}\Theta
   =\mu(\mathcal F)_{\mathbb R}^{-1}
                      \int_{\mathcal F}\Theta(g)\,d\mu(g).
 \tag{SVW.5}
\]
The subscript $\mathbb R$ denotes the source's real value of an
extended nonnegative measure; it is zero at infinity.  Source complex
integrals are Bochner integrals with value zero when the function is
not integrable.  We state explicitly the one place where the
nonzero-integral hypothesis rules out that fallback convention.

Here are the retained inputs in mathematical language.

\begin{itemize}
\item \texttt{SWa1} asserts existence of an ambient Haar measure and,
for \emph{every} ambient Haar measure $\mu$, existence of a fundamental
domain $\mathcal F$ for $\Gamma\backslash G$ and finiteness of
$\mu(\mathcal F)$ for \emph{every} such domain.  Its definition
contains no assertion about $H$ or $\Delta\backslash H$.
\item \texttt{Stabilizer\allowbreak First\allowbreak Chern}\allowbreak\texttt{Compatibility} applies to every
nondegenerate rational Hermitian $T$ of size $a$ with $2a<M$,
every specified rational basis, and every pure test with smooth compact
infinite factor.  For every ambient Haar measure, its every
fundamental domain, every allowed global character $\psi$, and
\emph{every} pair of self-dual tuple/Hermitian measures at infinity
and at all finite places, it supplies continuous local densities
$g_\infty,g_v$ on the invertible Hermitian locus such that
\[
 \operatorname{Mean}_{\mu,\mathcal F}\Theta_{T,\Phi}
       =g_\infty(T)\prod_v g_v(T),
 \tag{SVW.6}
\]
and the product is convergent.  Each density satisfies the original
\texttt{Is\allowbreak Cont\allowbreak Density\allowbreak On}: integration of $\Phi_*$ over the inverse
image of any compact Hermitian set in that locus equals integration
of $g_*$ over that set.  The original statement also asks rapid decay
of $\Phi_\infty$; compact support supplies it in SVW5.  The nested
existence-and-universality clauses in this premise are retained.
\end{itemize}

The original \texttt{Is\allowbreak Om\allowbreak T} for a measure $\nu$ on $W$ says that it is
supported on $Y_T$ and has the following property.  Self-dual tuple
and Hermitian measures and full continuous fiber families exist in
the specified nested order; for \emph{every} such collection and every
admissible test, the test is integrable, the local integral product
converges, and
\[
 \int_W\Phi\,d\nu
   =\left(\int\Phi_\infty\,d\omega_{\infty,T}\right)
      \prod_v\int\Phi_v\,d\omega_{v,T}.
 \tag{SVW.7}
\]
The original \texttt{Is\allowbreak Compat\allowbreak T} for supplied $\mu,\mu_H,\nu$ is
support on the actual orbit and the identity
\[
 \int_G f(g)\,d\mu(g)
   =\int_W\int_H f(s(y)h)\,d\mu_H(h)\,d\nu(y)
 \tag{SVW.8}
\]
for every nonnegative measurable $f$ and every section $s$ of the
orbit map.  No measurability assumption on $s$ is added.  The source
uses extended nonnegative integrals in (SVW.8).

We reuse GMR1--5 for the character, self-dual measures and pointwise
continuous fibers, GMR9--12 for the adelic Gram measure and its
quotient normalization, and GMR14 for local unimodularity.  These are
human arguments already proved, not invocations of formal certificates.
The proof below makes the additional universal-measure and arithmetic
quotient steps explicit.

\subsection{The same fiber measure in the density identity}

% BEGIN SVW1
\paragraph{SVW1. The adelic topology and the rational discrete subgroups.}
\label{claim:SVW1}
The tuple space $W$, the group $G$ and the closed stabilizer $H$ are
locally compact, Hausdorff and second countable.  The rational tuple
space $V^a$ is a closed discrete subset of $W$.  The groups $\Gamma$
and $\Delta$ are countable discrete subgroups of their respective
ambient groups.

\begin{proof}
Choose one $F$-basis of $E$.  Tensor coordinates identify $W$ with
a finite Cartesian power of $\mathbb A_F$ and its finite part with
the restricted product of the corresponding finite local vector
spaces, relative to their coordinate integer lattices.  This is the
actual tensor topology, as in AMD1.  There are countably many finite
places.  Each local field has a countable dense subfield and is a
metrizable locally compact space; its integer ring is compact open.
The restricted product is the union, over finite sets of places, of
open products having only finitely many unrestricted local factors.
Each such product is locally compact and second countable.  A
countable union of their countable bases is a basis for the restricted
product.  The finite-dimensional infinite factor has the same
properties.  This proves the assertions about $W$ and matrix spaces.

We recall why the diagonal $F$ is discrete in $\mathbb A_F$.  Its
intersection with $F_\infty\times\prod_v\mathcal O_v$ is
$\mathcal O_F$ in the infinite coordinates.  In a rational basis,
the Minkowski embedding of $\mathcal O_F$ is a lattice: an integral
basis gives an invertible real coordinate matrix, whose determinant
is nonzero by separability.  A sufficiently small infinite
neighborhood of zero therefore meets that diagonal intersection only
at zero.  This gives an adelic neighborhood isolating zero in $F$.
A discrete subgroup of a Hausdorff topological group is closed: if
there were an accumulation point, two distinct subgroup elements
arbitrarily close to it would have quotient arbitrarily close to
the isolated identity.  Thus the diagonal is closed.  Finite tensor
coordinates give the assertion for $V^a$.

The equation $g^*Jg=J$ defines a closed subset of matrix space.
It forces $\det g$ to be a unit and gives the inverse
$J^{-1}g^*J$, continuously.  Hence this closed subset has exactly
the topology and group operations of the source's $G$.  The further
equations $gX_0=X_0$ make $H$ closed.  Finally the rational matrices
are countable, and their entries lie in the just established
discrete diagonal tensor space.  Restriction to $G$ and then $H$
proves the asserted countability and discreteness.  All Haar
measures on these groups are consequently Radon and sigma-finite.
\end{proof}
% END SVW1

% BEGIN SVW2
\paragraph{SVW2. Admissible cutoffs and the full original finite-theta statement.}
\label{claim:SVW2}
Every original admissible test is continuous and compactly supported
on $W$.  Every compact $K\subset W$ has a nonnegative admissible
real test $\rho$, with smooth infinite factor, satisfying
$\rho\geq1$ on $K$.  For every admissible $\Phi$ and every $g\in G$,
only finitely many rational tuples have $\Phi(g^{-1}X)\ne0$.
In particular this proves the full original \texttt{Fin\allowbreak Theta}, and
the associated theta functions are measurable.

\begin{proof}
Fix the basis in the admissibility data and enlarge a finite set
$S$ to contain all exceptional local factors.  A finite-adelic tuple
is in the coordinate lattice outside finitely many places.  Thus its
local factors are one outside a finite set, so the source's finite
product is the ordinary finite product of its nontrivial values.
A nonzero product requires membership in the compact support of
each factor at $S$ and in the coordinate lattice outside $S$.
These conditions define a compact set: the unrestricted factors
are finitely many compact supports and the remaining factors are
compact lattices.  The same set is contained in an open
restricted-product chart.  There the function is a finite product
of locally constant functions; outside that chart it vanishes.
Since the chart is clopen, extension by zero is continuous.
Together with the compact infinite support this proves the first
assertion.

For the cutoff, cover the compact finite projection of $K$ by the
open restricted-product charts and take a finite subcover.  One
finite set $S$ contains all their exceptional places.  At each
$v\in S$, finitely many compact open local boxes cover the projected
compact set; their union is compact open.  Use its indicator, and
use the lattice indicator outside $S$.  Choose a smooth nonnegative
compact infinite cutoff at least one on the infinite projection.
Their product is the required $\rho$.  Smoothness in any other
allowed additive chart follows because a continuous additive map
between finite-dimensional real vector spaces is real linear:
first use rational linearity and then continuity.  Thus the change
of chart is smooth and preserves smoothness of the cutoff.

For fixed $g$, the possible rational tuples lie in
$g\operatorname{supp}\Phi$, a compact set.  A closed discrete subset
meets a compact set in a finite set: the intersection is compact
and discrete, and its cover by singleton neighborhoods has a finite
subcover.  Apply SVW1.  This gives precisely \texttt{Fin\allowbreak Theta}, with
no restriction on the rational tuples other than those in its
original statement.  The rational set is countable, and each term
$g\mapsto\Phi(g^{-1}X)$ is continuous.  The theta function is the
pointwise limit of its finite partial sums, hence measurable.
For a nonnegative real test it is pointwise nonnegative and finite.
\end{proof}
% END SVW2

% BEGIN SVW3
\paragraph{SVW3. Every supplied \texttt{Is\allowbreak Om\allowbreak T} is the same Radon measure.}
\label{claim:SVW3}
Fix any allowed character $\psi$ and the data (SVW.1)--(SVW.2).
Every measure $\nu$ satisfying the full original \texttt{Is\allowbreak Om\allowbreak T}
equals the measure $\nu_T$ of GMR9.  It is nonzero, Radon and positive
on every nonempty relatively open subset of $Y_T$.  This assertion
does not select new local measures or a new normalization.

\begin{proof}
The existence clauses inside \texttt{Is\allowbreak Om\allowbreak T} allow us to choose one
collection of self-dual measures and full fiber families, in their
prescribed order.  GMR2 and GMR5 identify them with the corresponding
measures in GMR9.  Equation (SVW.7) therefore gives
\[
                \int\Phi\,d\nu=\int\Phi\,d\nu_T
          \quad\hbox{for every admissible }\Phi.
 \tag{SVW.9}
\]
It also asserts actual integrability, rather than only an equality
under a zero fallback convention.

For a compact $K$, choose the cutoff $\rho$ of SVW2.  Integrability
and $\rho\geq1_K$ imply
$\nu(K)\leq\int\rho\,d\nu<\infty$.  Thus $\nu$ is locally finite.
On a locally compact second countable Hausdorff space, a locally
finite Borel measure is Radon.  One may see the relevant regularity
by using a countable basis of relatively compact open sets, finite
regularity on their compact closures, and monotone exhaustion.
In particular both measures in (SVW.9) are sigma-finite Radon
measures; no such hypothesis was silently imposed on the supplied
$\nu$.

The admissible tests determine them.  Indeed, a compact open finite
adelic box is a product of compact open local boxes with the
standard lattice at almost all places.  Its indicator is a pure
finite factor allowed in (SVW.4).  Multiplying it by any continuous
compact infinite test gives an admissible function.  Equality
(SVW.9), followed by the usual continuous cutoff approximation in
the finite-dimensional infinite space, proves equality on all
relatively compact open infinite sets times these finite boxes.
Those rectangles form a basis stable under finite intersections
after finite refinement.  Both spaces have countable such bases.
Disjointification and monotone convergence therefore give equality
on open sets, and Radon regularity gives equality on all Borel sets.
Hence $\nu=\nu_T$.  The nonzero and full-support assertions are now
exactly the properties proved in GMR9, with support understood on
the actual Gram fiber.  The argument used a common collection of
local witnesses only to compare the two measures; it proves equality
for every originally supplied $\nu$.
\end{proof}
% END SVW3

% BEGIN SVW4
\paragraph{SVW4. Continuous density equals the prescribed fiber integral at the specified parameter.}
\label{claim:SVW4}
At infinity and at every finite place, let the tuple and Hermitian
Haar measures and the full fiber family be \emph{any} original
witnesses.  Let $\Phi_*$ be continuous with compact support at
infinity, or an original Schwartz--Bruhat test at a finite place.
If $g_*$ satisfies the original \texttt{Is\allowbreak Cont\allowbreak Density\allowbreak On} on the
invertible Hermitian locus, then
\[
                   g_*(T)=\int\Phi_*\,d\omega_{*,T}.
 \tag{SVW.10}
\]
These are the full original \texttt{Dens\allowbreak Fib\allowbreak Inf} and
\texttt{Dens\allowbreak Fib\allowbreak V}, including their arbitrary density and witness
quantifiers.

\begin{proof}
Write $q(X)=X^*JX$ and let $\mathcal B$ be the invertible Hermitian
locus.  It is open, and the rational hypothesis $\det T\ne0$
places the image of $T$ inside it at every completion.  A tuple
with invertible Gram matrix is independent: its column map has
left inverse $T^{-1}X^*J$.  Thus $q^{-1}(\mathcal B)$ is contained
in the independent locus $U$ on which the full fibers are defined.

Choose a continuous compact cutoff $\chi$ on the Hermitian space
with support in $\mathcal B$, equal to one on an open neighborhood
$B_0$ of $T$.  Such a cutoff exists in the real space, and a compact
open box gives one in the nonarchimedean space.  The function
$f(X)=\Phi_*(X)\chi(q(X))$ is continuous with compact support
contained in $U$.  The defining continuity clause of the fiber
family, applied separately to real and imaginary parts, shows that
$h(S)=\int f\,d\omega_{*,S}$ is continuous.  For $S\in B_0$ its
support clause gives $h(S)=\int\Phi_*\,d\omega_{*,S}$.

Let $C\subset B_0$ be compact.  The original disintegration identity
for nonnegative functions gives, by splitting real and imaginary
parts and using their integrability,
\[
 \int_{q^{-1}(C)}\Phi_*\,dX
     =\int_C h(S)\,dT(S).
 \tag{SVW.11}
\]
All these integrals are finite: on the tuple side there is compact
support for an ambient Haar measure, and on the parameter side the
continuous fiber integral is bounded on $C$.  The supplied
continuous-density identity makes the same left side equal to
$\int_C g_*(S)\,dT(S)$.  Consequently $h$ and $g_*$ have identical
integrals on every compact subset of $B_0$.

They agree pointwise.  Otherwise, at a point where their difference
is nonzero, multiply that complex difference by a unit scalar so
its real part is positive there.  Continuity gives a neighborhood
where that real part is bounded below by a positive number.  A
compact neighborhood inside it has positive finite Haar measure,
contradicting equality of the two integrals.  In particular they
agree at $T$, proving (SVW.10).  Nothing in the argument depends on
which fiber family was supplied, and it does not replace fixed-$T$
equality by almost-everywhere equality.  The finite-place test is
continuous because it is locally constant.  Finite products of
real or split local factors cause no change to the argument.
\end{proof}
% END SVW4

% BEGIN SVW5
\paragraph{SVW5. The full original theta mean on the same \texttt{Is\allowbreak Om\allowbreak T} measure.}
\label{claim:SVW5}
Under the two retained premises, every original allowed character,
ambient Haar measure $\mu$, ambient fundamental domain $\mathcal F$,
supplied $\nu$ satisfying \texttt{Is\allowbreak Om\allowbreak T}, and admissible smooth test
$\Phi$ satisfy
\[
           \operatorname{Mean}_{\mu,\mathcal F}\Theta_{T,\Phi}
                          =\int\Phi\,d\nu.
 \tag{SVW.12}
\]
This is the full original \texttt{SWeq\allowbreak Om}.

\begin{proof}
Retain the basis and pure-product decomposition belonging to
$\Phi$.  A continuous compact infinite test has rapid decay in the
source's exact size function: for each $N$, the function
$|\Phi_\infty(X)|(1+\operatorname{size}(X))^N$ is continuous with
compact support, hence bounded.  This is precisely the required
inequality, with a constant allowed to depend on $N$ and the given
basis.  Thus the original compatibility premise applies to this
same test.

Choose self-dual measures and fiber families from the nested
existence clauses of the supplied \texttt{Is\allowbreak Om\allowbreak T}.  Apply the
universal parts of the compatibility premise at the already fixed
$\mu,\mathcal F,\psi$ and at these very self-dual measures.  It
returns densities $g_\infty,g_v$ and (SVW.6).  By SVW4 each of their
values at $T$ is the corresponding fiber integral for these same
witnesses.  Substitution in (SVW.6), place by place, gives the right
side of (SVW.7).  That equation identifies it with $\int\Phi\,d\nu$.
Both product expressions have their original convergence assertions,
so this substitution uses equality of the factors, not a rearranged
conditionally convergent series.

The choice of local witnesses was possible for every supplied
$\nu$, and each conclusion of the retained premise is universal
in those witnesses.  The character, ambient measure and fundamental
domain have never changed.  This proves all quantifiers of the
original statement.
\end{proof}
% END SVW5

\subsection{Arithmetic domains and rational orbit unfolding}

% BEGIN SVW6
\paragraph{SVW6. The rational Gram fiber is exactly one rational orbit.}
\label{claim:SVW6}
Every rational tuple $X$ with $X^*JX=T$ equals $\gamma X_0$ for a
rational unitary $\gamma\in\Gamma$.  The rational Gram fiber is
therefore identified with the right-coset set $\Gamma/\Delta$,
by $\gamma\Delta\mapsto\gamma X_0$.  This is the full original
\texttt{Witt} and the exact \texttt{gram\_\allowbreak orbit\_\allowbreak equiv} needed for
the theta sum.

\begin{proof}
Invertibility of $T$ makes both tuples independent and their spans
nondegenerate.  The map carrying the ordered columns of $X_0$ to
those of $X$ is consequently an isometry of nondegenerate
subspaces.  The extension argument already proved in GMR10 is
algebraic over any quadratic field extension, so it applies over
$E/F$ itself.  To recall the essential construction, diagonalize
the nondegenerate restricted form into nonisotropic lines and match
them successively in the orthogonal complement of the previously
matched lines.  Two vectors $u,v$ of norm $t\ne0$ are matched by
\[
 R(z)=z-(u-v)d^{-1}h(u-v,z),\qquad d=h(u-v,u),
 \tag{SVW.13}
\]
when $d\ne0$.  Expanding the Hermitian pairing and using
$h(u-v,u-v)=d+\overline d$ proves that $R$ is unitary and
$R(u)=v$.  If $d=0$, use $u,-v$, whose denominator is $2t$, and
then negate in the current complement.  All coefficients remain
in $E$ and all earlier matched vectors are fixed.  Composition
extends the ordered subspace isometry to the ambient space.

Thus the tuple map from $\Gamma$ is onto.  Two group elements have
the same image precisely when their quotient belongs to the
pointwise stabilizer, namely $\Delta$.  Diagonal rational embedding
is injective by SVW1, so the rational tuple and its adelic image
give the same orbit equivalence.  This also verifies that the
theta sum counts each rational tuple once, without an extra
stabilizer multiplicity.
\end{proof}
% END SVW6

% BEGIN SVW7
\paragraph{SVW7. A strict Borel fundamental domain without a mass assumption.}
\label{claim:SVW7}
For a second countable topological group $K$ and a discrete subgroup
$\Lambda$, there is a Borel set $D\subset K$ such that each left
$\Lambda$-orbit meets $D$ exactly once.  For the groups in (SVW.3)
this gives the full original \texttt{Ex\allowbreak FDStab}, for every supplied
Haar measure on $H$.  It does not assert finite mass.

\begin{proof}
Choose an identity neighborhood $V$ meeting $\Lambda$ only at the
identity, and then an open identity neighborhood $U$ with
$UU^{-1}\subset V$.  Second countability, and hence the Lindelof
property, gives a countable cover $K=\bigcup_{j\geq1}Ux_j$.
Two distinct points of a single $Ux_j$ cannot be left
$\Lambda$-equivalent: their quotient belongs to $UU^{-1}\cap
\Lambda=\{1\}$.

Define the open saturated sets
$S_j=\bigcup_{\lambda\in\Lambda}\lambda Ux_j$ and put
\[
 D_j=Ux_j\setminus\bigcup_{i<j}S_i,
                    \qquad D=\bigcup_{j\geq1}D_j.
 \tag{SVW.14}
\]
Each $D_j$ is Borel; the saturated union is open even without a
countability assertion about $\Lambda$.  Any orbit meets some
$Ux_j$; take the first such index.  There it contributes exactly
one point to $D_j$, and at all later indices it has been removed.
This proves both existence and uniqueness of the representative.
Thus $D$ is a strict Borel fundamental domain and is an allowed
fundamental domain for every measure.

Apply this with $K=H$ and $\Lambda=\Delta$, using SVW1.
The literal source subgroup \texttt{rat\allowbreak In} is this $\Delta$ inside
$H$: its two membership conditions are membership in $\Gamma$
and in $H$.  Therefore the result is the whole \texttt{Ex\allowbreak FDStab}
under its original contextual assumptions, which are not otherwise
needed in this construction.  The domain may have infinite Haar
mass at this stage.
\end{proof}
% END SVW7

% BEGIN SVW8
\paragraph{SVW8. Domain independence at the original almost-everywhere scope.}
\label{claim:SVW8}
Let $\Lambda$ be countable and let $\lambda$ be a left-invariant
measure on $K$.  If $D$ and $F$ are original fundamental domains,
then every nonnegative measurable left-$\Lambda$-invariant function
$f$ has
\[
                     \int_D f\,d\lambda=\int_F f\,d\lambda.
 \tag{SVW.15}
\]
The equality allows infinity.  In particular fundamental-domain
mass is independent of every original witness, including witnesses
which are only null measurable.

\begin{proof}
For a Borel strict domain, its translates partition the group.
For a supplied original domain, countable almost-everywhere
covering and pairwise almost-disjointness imply
\[
                  \sum_{\gamma\in\Lambda}1_F(\gamma x)=1
                         \quad\hbox{for almost every }x.
 \tag{SVW.16}
\]
To interpret indicators measurably, replace $F$ by a measurable
subset equal to it almost everywhere, and do the same for $D$.
Left invariance and countability ensure that this changes neither
the translated identities nor any of the integrals.

Insert (SVW.16) for $D$ in the integral over $F$ and use Tonelli.
For each summand change variable $y=\gamma x$.  Invariance of $f$
and $\lambda$ gives
\[
 \begin{split}
 \int_F f(x)\,d\lambda(x)
 &=\sum_{\gamma\in\Lambda}
      \int_K1_F(x)1_D(\gamma x)f(x)\,d\lambda(x)\\
 &=\int_D\sum_{\gamma\in\Lambda}1_F(\gamma^{-1}y)
                                     f(y)\,d\lambda(y)
 =\int_D f(y)\,d\lambda(y).
 \end{split}
 \tag{SVW.17}
\]
All summands are nonnegative, so no finite-volume or integrability
condition has entered.  Taking $f=1$ proves the mass assertion.
The same calculation applies to right domains and right-invariant
measures.  Countability holds for $\Gamma$ and $\Delta$ by SVW1.
\end{proof}
% END SVW8

% BEGIN SVW9
\paragraph{SVW9. Both actual adelic groups are unimodular.}
\label{claim:SVW9}
Every supplied Haar measure on $G$ and on the nondegenerate-tuple
stabilizer $H$ is invariant under right multiplication and inversion.
These give the full original \texttt{Unimod} and
\texttt{Unimod\allowbreak Stab} in the present source context.

\begin{proof}
GMR14 proves local unimodularity for arbitrary nondegenerate
unitary forms, with every Haar normalization.  The same trace-pairing
argument works at a real place: the real unitary Lie algebra has a
nondegenerate invariant real trace pairing, as can be checked after
complexification.  Conjugation therefore has absolute real
determinant one and preserves its Haar density.  Thus every local
factor of the adelic unitary group is unimodular.

Choose compact open subgroups at finite places and normalize their
local Haar masses to one.  Their restricted-product Haar measure,
together with the infinite product measure, is both left and right
invariant.  To check right invariance directly, an adelic group
element belongs to the chosen compact subgroups outside finitely
many places.  On any cylinder its right translation is therefore
tested in only finitely many unrestricted factors, where local
right invariance proves the equality; cylinder sets determine the
measure.  Every other Haar measure is a positive finite multiple,
so the assertion holds for the supplied ambient measure as well.

The span of $X_0$ is nondegenerate.  Orthogonal decomposition over
$E$ identifies its pointwise stabilizer with the unitary group of
the complementary nondegenerate form, over $F$ and over the adeles.
Restriction and extension by the identity are continuous inverse
maps, obtained from the rational orthogonal projection
$X_0T^{-1}X_0^*J$.  The preceding argument thus applies to $H$ and
to any of its supplied Haar measures, without imposing compactness
on this adelic group.

For a bi-invariant Haar measure, inversion sends it to a left Haar
measure with the same normalization.  For example, compare on a
symmetric compact neighborhood of positive finite mass, whose
inverse is itself; Haar uniqueness forces scalar one.  Inversion
therefore preserves both measures.  No quotient-volume statement
has been used.
\end{proof}
% END SVW9

% BEGIN SVW10
\paragraph{SVW10. Nonnegative unfolding before finite stabilizer mass.}
\label{claim:SVW10}
Take the actual $G,H,\Gamma,\Delta,X_0$ above, any supplied Haar
measures $\mu,\mu_H$, and any $\nu$ satisfying the original
\texttt{Is\allowbreak Compat\allowbreak T}.  For every nonnegative measurable $\phi$ on
tuple space and every pair of original fundamental domains
$\mathcal F,\mathcal F_H$, one has the extended-integral identity
\[
 \int_{\mathcal F}\sum_{\gamma\Delta\in\Gamma/\Delta}
               \phi(g^{-1}\gamma X_0)\,d\mu(g)
             =\mu_H(\mathcal F_H)\int_W\phi\,d\nu.
 \tag{SVW.18}
\]
Products and integrals here are nonnegative extended values,
with the usual convention $0\cdot\infty=0$.

\begin{proof}
First take strict Borel domains $D$ for $\Gamma\backslash G$ and
$E_H$ for $\Delta\backslash H$, whose existence was proved in
SVW7.  Choose one representative $\gamma$ for each right coset
$\gamma\Delta$ of $\Gamma$.  The inverse decomposition is
$\Gamma=\coprod_{\gamma\Delta}\Delta\gamma^{-1}$.  Hence
\[
          B=\coprod_{\gamma\Delta}\gamma^{-1}D
 \tag{SVW.19}
\]
is a strict Borel fundamental domain for the left $\Delta$ action
on $G$.  Countability of $\Gamma$ makes this a Borel disjoint union.
To verify the assertion explicitly, the sets
$\delta\gamma^{-1}D$, for all $\delta\in\Delta$ and the chosen
$\gamma$, are exactly the translates $\eta D$, $\eta\in\Gamma$,
and therefore partition $G$.

Tonelli, the left translation $g=\gamma b$ in each summand, and
inversion invariance from SVW9 give
\[
 \begin{split}
 \int_D\sum_{\gamma\Delta}\phi(g^{-1}\gamma X_0)\,d\mu(g)
     &=\int_B\phi(b^{-1}X_0)\,d\mu(b)\\
     &=\int_G1_{B^{-1}}(z)\phi(zX_0)\,d\mu(z).
 \end{split}
 \tag{SVW.20}
\]
The set $B^{-1}$ is a strict \emph{right} $\Delta$ fundamental
domain in $G$.  Apply (SVW.8) to the nonnegative measurable function
in the last line.  Choose a Borel section of the actual orbit map;
GMR10 supplies continuous local sections, whose countable
disjointification gives one.  Alternatively the original
all-section formula itself permits any section.  For a point $y$
on the orbit and any fixed representative $s(y)$, set
\[
            C_y=\{h\in H:s(y)h\in B^{-1}\}.
 \tag{SVW.21}
\]
This is Borel in $H$ and is a strict right $\Delta$ domain there:
the right $\Delta$ orbit of $s(y)h$ stays in $s(y)H$, and it has
exactly one representative in $B^{-1}$.  By the right-hand version
of SVW8 and inversion invariance,
\[
             \mu_H(C_y)=\mu_H(E_H^{-1})=\mu_H(E_H).
 \tag{SVW.22}
\]
No finiteness of this number is asserted or needed.  Since $H$
fixes $X_0$, the inner integral in (SVW.8) is thus
$\mu_H(E_H)\phi(y)$ for every $y$ on the orbit.  Its complement
has $\nu$-measure zero by compatibility.  Equations
(SVW.20)--(SVW.22) prove (SVW.18) for the strict domains.

Finally the orbit sum is a nonnegative measurable left-$\Gamma$-
invariant function: left translation reindexes the cosets in its
sum.  SVW8 changes $D$ to the supplied $\mathcal F$ without altering
its integral, and changes $E_H$ to the supplied $\mathcal F_H$
without altering its mass.  This proves (SVW.18) at the original
almost-everywhere fundamental-domain scope.  This direct argument
uses the actual locally compact groups; no stronger claim about
an abstract measurable-group unfolding theorem is required.
\end{proof}
% END SVW10

% BEGIN SVW11
\paragraph{SVW11. Finite mass at the correct stage: the full original \texttt{Unfold}.}
\label{claim:SVW11}
In the original context, let $\phi\geq0$ have admissible complex
extension, and let all of $\psi,\mu,\mathcal F,\nu,\mu_H,
\mathcal F_H$ be the supplied witnesses in \texttt{Unfold}.
If
\[
                  \int_{\mathcal F}\Theta_{T,\phi}\,d\mu\ne0,
 \tag{SVW.23}
\]
then $\mu_H(\mathcal F_H)<\infty$ and
\[
 \int_{\mathcal F}\Theta_{T,\phi}\,d\mu
          =\mu_H(\mathcal F_H)_{\mathbb R}
                                      \int_W\phi\,d\nu.
 \tag{SVW.24}
\]
The displayed finite integrals are regarded as complex numbers in
the source statement.

\begin{proof}
SVW2 gives a pointwise finite nonnegative measurable theta function.
Its complex Bochner integral is nonzero by (SVW.23), so it is
integrable: in the source definition a nonintegrable function has
zero integral.  Its nonnegative integral is therefore a finite
strictly positive real number $L$.  This is the precise use of
the nonzero-integral hypothesis, before any conversion of an
infinite measure to a real scalar.

The supplied \texttt{Is\allowbreak Om\allowbreak T} states that $\phi$ is integrable;
put $I=\int\phi\,d\nu$, a finite nonnegative real number.
By SVW6 the rational tuples in the theta sum are exactly the
cosets in SVW10.  That extended-integral identity says
\[
                       0<L=\mu_H(\mathcal F_H)I<\infty.
 \tag{SVW.25}
\]
It follows first that $I>0$ and then that
$\mu_H(\mathcal F_H)$ is finite.  In particular the otherwise
dangerous case $\infty\cdot I$ has been ruled out by a positive
factor, rather than by cancellation of real values of infinite
measures.  Taking real values of (SVW.25) now gives (SVW.24).
The argument applies to every supplied fundamental domain and
every supplied compatible Haar measure.  It is exactly the full
original \texttt{Unfold} conclusion.
\end{proof}
% END SVW11

\subsection{One positive test and the volume comparison}

% BEGIN SVW12
\paragraph{SVW12. An explicit admissible smooth test positive near the rational frame.}
\label{claim:SVW12}
There is a nonnegative real test $\phi$ with admissible smooth
complex extension and an open neighborhood $U_G$ of the identity
in $G$ such that $\phi(gX_0)>0$ for every $g\in U_G$.
This is the full original \texttt{Test\allowbreak Near}.

\begin{proof}
Start with an arbitrary $F$-basis $(b_i)$ of $E$.  There are only
finitely many coordinates of all columns of $X_0$.  Choose
$\beta\in\mathcal O_F\setminus\{0\}$ clearing their denominators
and replace the basis by $b_i/\beta$.  The coordinates of $X_0$
in this new basis are integral at every finite place.  Let $L_v$
be its tuple lattice in the local tensor space and put
$\phi_v=1_{L_v}$ at every finite place.  Each $L_v$ is compact
open, so each factor is a permitted Schwartz--Bruhat function.

In one additive real coordinate chart at infinity choose a smooth
nonnegative compact bump $\phi_\infty$ which is positive at
$(X_0)_\infty$.  For example a usual Euclidean bump supported in
a ball and equal to one on a smaller concentric ball does this.
The continuous-additive-change-of-coordinates argument in SVW2
shows that it satisfies the source's all-coordinate
\texttt{Is\allowbreak Smooth\allowbreak Cc}, not just smoothness in the chosen chart.
Define
\[
                  \phi(X)=\phi_\infty(X_\infty)
                                      \prod_v1_{L_v}(X_v).
 \tag{SVW.26}
\]
This is a nonnegative \texttt{Is\allowbreak Adm\allowbreak S} test with the literal
finite-factor product.

The finite-adelic lattice $\prod_v L_v$ is open in the restricted
product topology, not merely an intersection of infinitely many
unspecified open sets.  The orbit map is continuous by its matrix
formula.  Therefore the inverse image of this lattice, intersected
with the inverse image of $\{\phi_\infty>0\}$, is an open set
$U_G$ containing the identity.  On it all finite indicators are
one and the infinite factor is positive.  This proves the claim
without using any quotient-volume statement or a specially chosen
Haar normalization.
\end{proof}
% END SVW12

% BEGIN SVW13
\paragraph{SVW13. The same test has positive finite integral for every supplied compatible measure.}
\label{claim:SVW13}
For every original $\psi,\mu,\nu,\mu_H$ satisfying
\texttt{Is\allowbreak Om\allowbreak T} and \texttt{Is\allowbreak Compat\allowbreak T}, the test of SVW12 satisfies
\[
                         0<\int_W\phi\,d\nu<\infty.
 \tag{SVW.27}
\]
In particular the full original \texttt{Pos\allowbreak Test} holds.

\begin{proof}
Admissibility and the actual integrability clause of
\texttt{Is\allowbreak Om\allowbreak T} make the integral finite.  If it were zero,
nonnegativity would give $\phi=0$ almost everywhere for $\nu$.
Apply compatibility (SVW.8) to the measurable indicator $1_{U_G}$,
using any section $s$.  For almost every $y$, the point is on
the orbit and $\phi(y)=0$.  There is then no $h\in H$ with
$s(y)h\in U_G$, since such an element would satisfy
\[
                 0<\phi(s(y)hX_0)=\phi(y)=0.
\]
Thus the inner integral of $1_{U_G}$ is zero almost everywhere,
and compatibility would give $\mu(U_G)=0$.  This contradicts
positivity of Haar measure on the nonempty open set $U_G$.
The contradiction proves positivity.

This argument works with each supplied compatible measure and
does not presume that a quotient measure is already positive on
open sets.  Equivalently one could use the full-support property
of SVW3, but compatibility alone suffices.  Both arguments use
the same test as the density identity and the later unfolding.
\end{proof}
% END SVW13

% BEGIN SVW14
\paragraph{SVW14. All existence stages of the original nested volume statement.}
\label{claim:SVW14}
In the context (SVW.1)--(SVW.2), the original \texttt{Ex\allowbreak Char},
\texttt{Ex\allowbreak Om\allowbreak T}, \texttt{Ex\allowbreak Stab} and \texttt{Ex\allowbreak FDStab} hold at their
full scopes.  In particular, for \emph{every} permitted character,
every ambient Haar measure, and \emph{every supplied}
\texttt{Is\allowbreak Om\allowbreak T} measure, a compatible stabilizer Haar measure exists.
For every such Haar measure a stabilizer fundamental domain exists.

\begin{proof}
GMR1 supplies an allowed character, and GMR2--GMR9 supply an
\texttt{Is\allowbreak Om\allowbreak T} measure for every allowed character in the original
range $1\leq a$, $2a<M$.  The full product convergence and positive
good-place masses in that construction have already been proved;
no new restricted-product normalization is postulated here.

Now fix an arbitrary supplied $\nu$ satisfying \texttt{Is\allowbreak Om\allowbreak T},
not only the measure first constructed by GMR9.  SVW3 identifies
it with that canonical measure.  By GMR10 its fiber is the actual
adelic orbit, with the quotient topology and continuous local
sections, and it is invariant under $G$.  It is positive and
Radon.  GMR11 therefore supplies a unique stabilizer Haar
normalization satisfying (SVW.8) for the already supplied ambient
$\mu$ and this same $\nu$, with the full all-section conclusion.
This is the universal assertion required by \texttt{Ex\allowbreak Stab}.
Using only the existential pair in GMR12 would not by itself
establish this stage for an arbitrary supplied $\nu$; SVW3 is what
justifies that extension.

For every supplied $\mu_H$, SVW7 supplies a Borel domain for
$\Delta\backslash H$.  For every ambient Haar measure, the
retained \texttt{SWa1} supplies the ambient domain and its finite
mass; Haar existence is also part of that premise.  These facts
give all existence stages in their original order.  The
stabilizer-domain stage still carries no finite-mass assertion:
that will follow from SVW11 after applying the positive test.
\end{proof}
% END SVW14

% BEGIN SVW15
\paragraph{SVW15. Full original \texttt{Vol\allowbreak Stab\allowbreak From\allowbreak SWS} and scalar cancellation.}
\label{claim:SVW15}
For every original CM extension and nondegenerate Hermitian form,
the exact premises \texttt{SWa1} and
\texttt{Stabilizer\allowbreak First\allowbreak Chern}\allowbreak\texttt{Compatibility} imply the full original
\texttt{Vol\allowbreak Stab\allowbreak W}.  Thus for every tuple in (SVW.2), the statement
\texttt{Vol\allowbreak Stab\allowbreak At} holds, with its nested witness order
\[
 \psi,\quad\mu,\quad\nu,\quad\mu_H,\quad
                       \mathcal F,\quad\mathcal F_H,
 \tag{SVW.28}
\]
existence at each stage and universality over every witness there.
Its final assertion is the exact extended-real equality
\[
                         \mu_H(\mathcal F_H)=\mu(\mathcal F).
 \tag{SVW.29}
\]
This proves the original \texttt{vol\allowbreak Stab\allowbreak From\allowbreak SW\_\allowbreak of} assembly with
all seven of its named suppliers reconstructed at the needed scopes.

\begin{proof}
SVW14 gives all existence clauses.  Fix any witnesses in (SVW.28)
with their stated predicates.  The supplied ambient fundamental
domain has finite mass by \texttt{SWa1}; it has positive mass
because its countably many rational translates cover almost all
of the nonzero Haar measure.  Put $m=\mu(\mathcal F)_{\mathbb R}$,
so $0<m<\infty$.

Use the single test of SVW12--SVW13 and put
$I=\int\phi\,d\nu$.  By SVW13, $I$ is finite and strictly
positive.  The density identity SVW5 says
\[
                    m^{-1}\int_{\mathcal F}\Theta_{T,\phi}\,d\mu=I.
 \tag{SVW.30}
\]
In particular the theta integral is nonzero.  We may therefore
apply the full original unfolding statement SVW11.  It proves
\emph{now} that $\mu_H(\mathcal F_H)$ is finite and gives
\[
 \int_{\mathcal F}\Theta_{T,\phi}\,d\mu
                        =\mu_H(\mathcal F_H)_{\mathbb R}I.
 \tag{SVW.31}
\]
Substituting (SVW.31) into (SVW.30) and cancelling the positive
finite $I$ yields
$m^{-1}\mu_H(\mathcal F_H)_{\mathbb R}=1$.  Hence the two real
masses are equal.  Both extended masses are finite, so conversion
to their real values is injective and gives (SVW.29).

This is exactly the original cancellation argument, with its
analytic suppliers now proved: character and measure existence,
domain existence, same-measure theta density, unfolding and
positive test.  The argument did not choose a second character,
a new ambient Haar normalization, different self-dual measures
between the two identities, or a preferred fundamental domain
in place of the supplied ones.  Since the witnesses were arbitrary,
the full nested universality follows.  The retained premises have
not themselves been proved by this cancellation.
\end{proof}
% END SVW15

\subsection{The literal actual-setting volume join}

% BEGIN SVW16
\paragraph{SVW16. The original choice-based covolumes are equal.}
\label{claim:SVW16}
At every original tuple in the strict range (SVW.2), with every
allowed $\psi,\mu,\nu,\mu_H$ satisfying their original conditions,
\[
       \texttt{covol}(\mu_H)
           =\texttt{covol\allowbreak By}(\Gamma,\mu)>0.
 \tag{SVW.32}
\]
The equality alone is the full original
\texttt{covol\_\allowbreak eq\_\allowbreak covol\allowbreak By\_\allowbreak w}, even for an arbitrary supplied
\texttt{Vol\allowbreak Stab\allowbreak W}.  The added positivity in (SVW.32) uses the
retained \texttt{SWa1} and the conclusion of SVW15.

\begin{proof}
The source defines each covolume as the real mass of a chosen
fundamental domain when one exists, and as zero otherwise.
SVW14 supplies existence of both domains, so neither zero fallback
is used.  For the stabilizer, the subgroup called \texttt{rat\allowbreak Of}
is the inverse image of $\Gamma$ under $H\hookrightarrow G$.
The subgroup called \texttt{rat\allowbreak In} is $(\Gamma\cap H)$ regarded
inside $H$.  An element of $H$ lies in one exactly when it lies
in the other, so they agree literally.

Take the particular domains chosen by the two definitions.
SVW15 applies to \emph{every} pair of domains and proves equality
of their extended masses; these masses are finite and positive.
Taking real values proves (SVW.32).  This also shows that the
answer does not depend on the choice functions used in the
definitions.  For an arbitrary supplied \texttt{Vol\allowbreak Stab\allowbreak W}, its
nested clauses already give existence of both domains and equality
of their extended masses.  Applying the real-value function to
that equality proves equality of the choice-based covolumes even
without a separate finite-mass hypothesis.  Positivity is not
asserted in this more general transport statement.  This is the
exact scope of the original covolume transport theorem.
\end{proof}
% END SVW16

% BEGIN SVW17
\paragraph{SVW17. The full original \texttt{In\allowbreak Volumes\allowbreak At} implication.}
\label{claim:SVW17}
Let $X$ be any original supplied \texttt{Hida\allowbreak Setting}, with its
original bridge $h_b$, level-character witness, specified
$\mathrm{FdN},\mathrm{FdN}'$, words $u,u'$, local and analytic
records $L,A$, and auxiliary tuple $X_0$ of length $2r+2$, where
$r=D.R.rS$.  Suppose the exact two premises of SVW15 hold for
$X.B.I$.  Then the full original
\[
 \texttt{In\allowbreak Volumes\allowbreak At}(X,h_b,h_K,\mathrm{FdN},\mathrm{FdN}',
                                    u,u',X_0,L,A)
 \tag{SVW.33}
\]
holds.  Its two unipotent covolume-one antecedents remain part of
this generic implication, as in the source definition.

\begin{proof}
Write $n=X.B.I.n$.  The original rank constraints give
$m=2(r+1)$, $k\geq m+3$, and $n=k+m-1$.  Hence
$n\geq4r+6$, in particular the source's $4r+5\leq n$.
Both tuple lengths $2r+1$ and $2r+2$ are positive and satisfy
$2a<n$.  This proves the literal \texttt{step\allowbreak Size\allowbreak W} needed for
all three applications on the same setting.

For each word the original tuple is
$(e_1,\ldots,e_r,z_u,f_1,\ldots,f_r)$.  Its Gram matrix is
\[
 \begin{pmatrix}
   0&0&I_r\\0&-\operatorname{Tw}(u)&0\\I_r&0&0
 \end{pmatrix},
 \tag{SVW.34}
\]
in the actual coordinates of $X$.  The setting's basis and form
identity give these entries, as checked in GMR19 and AOS9.  Its
determinant is $(-1)^{r+1}\operatorname{Tw}(u)\ne0$; the same
calculation holds for $u'$.  This gives the actual
\texttt{Step\allowbreak Gram\allowbreak H} input, without choosing new frame coordinates.

Now take \emph{all} antecedents of (SVW.33), in their original
order: \texttt{Hyp\allowbreak X}, one $\psi$, the two word Gram-symmetry and
\texttt{Is\allowbreak Om\allowbreak T}/compatibility witnesses, the auxiliary Gram-symmetry,
nonzero determinant and corresponding witnesses, and finally
\[
 \texttt{covol}((X.\texttt{meas\allowbreak W}(u)).\mu_N)=1,
 \qquad
 \texttt{covol}((X.\texttt{meas\allowbreak W}(u')).\mu_N)=1.
 \tag{SVW.35}
\]
Their Haar hypotheses are supplied by the existing normalized
measure records of $X$.  Put
$c=\texttt{covol\allowbreak By}(\Gamma,X.\texttt{meas\allowbreak G}.\mu_G)$.
SVW16, with exactly the supplied orbit measures and the same
character, gives all three equalities
\[
 \begin{gathered}
 \texttt{covol}(X.\texttt{meas\allowbreak G}.\mu_B(X_0))=c,\\
 \texttt{covol}((X.\texttt{meas\allowbreak W}(u)).\mu_H)=c,
 \qquad
 \texttt{covol}((X.\texttt{meas\allowbreak W}(u')).\mu_H)=c.
 \end{gathered}
 \tag{SVW.36}
\]
The auxiliary tuple is arbitrary; only its explicitly supplied
nondegenerate Gram matrix has been used.

The first field of the source's \texttt{In\allowbreak Volumes} is
\[
 \texttt{covol}(\mu_B(X_0))\,c
       =\texttt{covol}(\mu_H(u))\,
                               \texttt{covol}(\mu_H(u')).
 \tag{SVW.37}
\]
Substitution of (SVW.36) gives $c^2=c^2$.  Its other two fields
are exactly (SVW.35), which are retained antecedents in this
generic statement.  This proves the full \texttt{In\allowbreak Volumes\allowbreak At}
contract.  If the original \texttt{SWXS} is supplied for all
settings, the argument is universal in $X$ and is precisely
\texttt{in\allowbreak Volumes\allowbreak At\_\allowbreak of\_\allowbreak SWS}.  None of $h_b$, the infinity
measures, the fundamental-domain sets, $L$, or $A$ has changed.
\end{proof}
% END SVW17

% BEGIN SVW18
\paragraph{SVW18. On the actual rescaled setting the volume fields have all their measure inputs.}
\label{claim:SVW18}
Retain the original \texttt{Siegel\allowbreak Weil\allowbreak Unitary\allowbreak Convergent} premise
and let $X'$ be the actual GMR19 rescaling of a supplied $X$,
with its same bridge and original infinity measures.  For every
pair of words, every auxiliary $X_0$ of length $2r+2$ with
nonzero Gram determinant, and every original supplied $h_K,L,A,
\mathrm{FdN},\mathrm{FdN}'$ satisfying the literal
\texttt{Hyp\allowbreak X} conditions for $X'$, the actual record
\[
   b=\texttt{b\allowbreak Of}(X',h_b,h_K,\mathrm{FdN},\mathrm{FdN}',
                                      u,u',X_0,L,A)
 \tag{SVW.38}
\]
satisfies its full original \texttt{In\allowbreak Volumes} field.

\begin{proof}
The unchanged Hermitian datum of $X'$ has the exact
\texttt{SWa1} and \texttt{Stabilizer\allowbreak First\allowbreak Chern}\allowbreak\texttt{Compatibility}
specializations of the retained original Siegel--Weil premise,
as recorded in GMR19.  Thus SVW17 supplies (SVW.33) for $X'$.
GMR19 selected one character before all words, finite places and
auxiliary tuples, and its resulting \texttt{Meas\allowbreak Compat\ensuremath\psi}
contains, for every word, an \texttt{Is\allowbreak Om\allowbreak T} measure compatible
with the actual $\mu_G$ and $\mu_H(u)$, together with the actual
equality $\texttt{covol}(\mu_N(u))=1$.  For every nondegenerate
$X_0$ it contains the corresponding compatible measure for
$\mu_B(X_0)$.  These are exactly the measure antecedents of
SVW17, now all supplied on this one setting.

Choose the two word witnesses and the auxiliary witness from that
same \texttt{Meas\allowbreak Compat\ensuremath\psi}.  Different proofs of the
Hermitian symmetry of the same matrix merely identify the same
element of the Hermitian subspace and do not change any measure
or equality.  Apply SVW17 with these witnesses, the common
character, the supplied nonzero determinant, and the two already
proved covolume-one equalities.  The result is $b.\texttt{In\allowbreak Volumes}$.
This is the literal \texttt{in\allowbreak Volumes\_\allowbreak of\_\allowbreak compat} return, with
its actual compatibility premise discharged by GMR19.

In particular no independent unipotent-volume supplier remains
for this join on the rescaled setting: AMD17 and GMR19 already
provided it.  The generic implication of SVW17 keeps those
antecedents because an arbitrary supplied setting need not have
that normalization.  The Haar measures at infinity, the bridge,
the rational frames, the auxiliary tuple and both words are all
the original objects.  This conclusion proves the three specified
volume fields, not the other fields of the mass-data record.
\end{proof}
% END SVW18

\subsection{The exact remaining analytic boundary}

The retained smooth Siegel--Weil compatibility theorem has been
used at its original scope.  This chapter does not prove that
theorem, nor the broader arithmetic theta identities contained in
\texttt{Siegel\allowbreak Weil\allowbreak Unitary\allowbreak Convergent}.  It proves the stabilizer
volume comparison \emph{from} that original premise and supplies
its full actual-setting volume join.

The completed AOS scalar is unchanged.  The distinct
\texttt{Real\allowbreak Bracket} calculation still needs its reference
unipotent covolume/discriminant calculation and its local
Jacobian, place-splitting and product-transfer producers; the
present arithmetic covolume equality is not a replacement for
those equations.  Likewise finite-place factors, the remaining
mass/counting fields, ambient \texttt{Forms0}, the all-family
transport and level match, Bessel uniqueness and spectral ties
remain separate.  No whole selected RH or Dirichlet-GRH
obligation is asserted closed here.


\section{Descent of the actual invariant Frobenius matrices}
\label{sec:invariant-frobenius-descent-v055}

\status{This chapter proves the full original inhomogeneous matrix
descent theorem, on the supplied Wach coefficient ring and the actual
Witt tensor ring. It also proves the abstract additive and matrix
bijectivity statements used in that descent. The approximation argument
uses the exact sharp-root coordinates of an invariant tilt element and
the previously proved completed-fixed-field theorem. Together with WIS18,
this leaves the actual Wach prism as the one remaining period-ring
supplier in the conditional all-layer return WLD26. The original tied
period data and crystalline-realization input are retained.}

There are two different approximations in this proof. Powers of the
cyclotomic parameter control the zeroth Witt coordinate; powers of $p$
control the successive Witt coordinates. Frobenius improves the first
approximation. Saturation then permits division by $p$, so that the
argument can be repeated. The final passage to the coefficient ring uses
its $p$-adic closed image, already proved in WIS19.

Keep the original prime $p$, period data $P,AI$, finite fields $K,LL$,
level $s$, and supplied records $\mathrm{Good}(K,s,LL)$,
$\Gamma,R,I,\mathrm{IoW}(I)$. In the notation of WIS put
\[
 O=\mathcal O_{LL},\quad S=\mathcal O_C^\flat,\quad F=\operatorname{Frac}(S),
 \quad B=O\otimes_{\mathbb Z_p}W(S),\quad W=O\otimes_{\mathbb Z_p}W(F).
\]
Write $c:B\hookrightarrow W$, $A=\prod_{i\in\mathbb Z/n\mathbb Z}O[[T]]$,
$E=R.\mathrm{EE}$, $a:A\to E$, $i:A\to B$ and $j=I.\mathrm{ioE}:E\to W$.
Thus $ja=ci$. The symbol $T$ in $A$ means the tuple whose every entry
is the power-series variable. Put
\[
 \tau_s=\epsilon_s-1\in S,\quad \mu_s=[\epsilon_s]-1\in W(S),\quad
 t=c(1\otimes\mu_s)=j(a(T)),\qquad
 \mathcal N_m=t^mc(B)+pW.                                      \tag{IFD.1}
\]
Here $t$ is a unit of $W$. Accordingly $\mathcal N_m$ must not be
replaced by the ideal $t^mW+pW$. The coefficient $c(B)$ in (IFD.1)
is indispensable. The fixed ring used at the endpoint is
\[
 W^H,\qquad H=G_K\cap\ker\chi.
                                                               \tag{IFD.2}
\]
The finite coefficient basis in WIS4 is kept throughout. Its coordinates
will be denoted by $w=\sum_\alpha b_\alpha\otimes w_\alpha$.
The formulas and completeness statements of WIS18--WIS25 are already
human proofs and are used at their stated scopes.

\subsection{The filtration and its closed coefficient image}

% BEGIN IFD1
\paragraph{IFD1: the elementary near-parameter calculus.}
\label{claim:IFD1}
More generally let $c:B_0\to W_0$ be a ring map, let
$t=c(\tau)$, and put $N_m=t^mc(B_0)+pW_0$ for $m\geq0$.
Then $N_m$ is an additive subgroup, $N_{m+1}\subseteq N_m$,
$N_mN_k\subseteq N_{m+k}$, and $c(B_0)N_m\subseteq N_m$.
If $t$ is a unit, then
\[
             t^k x\in N_{m+k}\quad\Longrightarrow\quad x\in N_m.
                                                               \tag{IFD.3}
\]
For the actual rings, membership in $\mathcal N_m$ is equivalent to
$(w_\alpha)_0\in\tau_s^m S$ for every $\alpha$, with $S$ embedded in $F$.

\begin{proof}
Addition and negation act on the two coefficients in a presentation
$x=t^mc(b)+pr$. Descending the exponent multiplies $b$ by $\tau$.
Multiplying two such presentations gives leading term
$t^{m+k}c(bb')$ and a remainder divisible by $p$; the assertion about
$c(B_0)$ is the same calculation. In (IFD.3), divide the presentation
$t^kx=t^{m+k}c(b)+pr$ by the unit $t^k$.
The last assertion is exactly WIS23: in the unchanged finite $O$-basis,
zeroth Witt coordinates detect the ideal $pW$, and the zeroth coordinate
of $\mu_s$ is $\tau_s$. This proves the equivalence with the precise
subgroup in (IFD.1).
\end{proof}
% END IFD1

% BEGIN IFD2
\paragraph{IFD2: Frobenius improves the filtration and detects constants.}
\label{claim:IFD2}
The original Frobenius $\phi$ satisfies
\[
 \phi(c(B))\subseteq c(B),\qquad
 \phi(\mathcal N_m)\subseteq\mathcal N_{pm},\qquad
 \phi(t)\in\mathcal N_2.                                      \tag{IFD.4}
\]
In particular it preserves every $\mathcal N_m$. If the image in $W$
of $o\in O$ belongs to $\mathcal N_1$, then $o\in pO$.
These are all four clauses of the original
\texttt{frobenius\_\allowbreak and\_\allowbreak constants} theorem.

\begin{proof}
Frobenius on the Witt factor is induced by $x\mapsto x^p$ on $S$ or
$F$, and it fixes the $O$ factor. It therefore preserves $c(B)$.
If $(w_\alpha)_0=\tau_s^m v_\alpha$, then
$(\phi w_\alpha)_0=\tau_s^{mp}v_\alpha^p$.
IFD1 proves the middle inclusion. Since $t\in\mathcal N_1$ and
$p\geq2$, it proves the third assertion and also preservation of
each subgroup. These arguments include $p=2$ and every $s\geq0$.

Expand $o=\sum_\alpha z_\alpha b_\alpha$ with $z_\alpha\in\mathbb Z_p$.
Its $\alpha$th Witt coordinate is the image of $z_\alpha$ in $W(F)$.
If $p\nmid z_\alpha$, then $z_\alpha$ is a unit of $\mathbb Z_p$,
so its image in $W(S)$, and then its zeroth coordinate in $S$, are
units. That coordinate cannot lie in $\tau_s S$: WIS22 implies
that $\tau_s$ is a nonunit, since otherwise $1$ would lie in every
$\tau_s^mS$. IFD1 thus forces $p\mid z_\alpha$ for every $\alpha$.
Dividing the finitely many coefficients gives $o\in pO$.
\end{proof}
% END IFD2

% BEGIN IFD3
\paragraph{IFD3: a near-divisible component idempotent is zero.}
\label{claim:IFD3}
Let $c:B_0\hookrightarrow W_0$ be injective, with $p$-saturated image,
where $W_0$ has no $p$-torsion and is $p$-adically separated.
Let $t=c(\tau)$ be a unit, and suppose $\bigcap_mN_m\subseteq pW_0$.
If $e,c_0,d,w\in B_0$ and $N\geq0$ satisfy
\[
 e^2=e,\qquad p=c_0d,\qquad d^{N+1}=pw,
                 \qquad c(c_0e)\in N_1,
                                                               \tag{IFD.5}
\]
then $e=0$.

\begin{proof}
Write $c(c_0e)=tc(b)+pr$. Saturation gives $r=c(r_0)$ and
injectivity gives $c_0e=\tau b+pr_0$ in $B_0$.
Multiplication by $d$ yields
$p(e-dr_0)=\tau db$. In $W_0$ this says
$c(db)=p t^{-1}c(e-dr_0)$. Saturation gives an $s\in B_0$ with
$c(s)=t^{-1}c(e-dr_0)$, and hence
\[
                              e=dr_0+\tau s.
\]
Raise this equality to the power $N+1$. Idempotence gives
$e^{N+1}=e$, the term $(dr_0)^{N+1}$ is divisible by $p$, and every
other term is divisible by $\tau$. Thus $e\in\tau B_0+pB_0$.
Raising this last expression to any positive power shows
$c(e)\in N_m$ for every $m$. The intersection hypothesis gives
$c(e)\in pW_0$, and saturation and injectivity give $e=pe_1$ in $B_0$.
For every $k$, idempotence now gives
$e=e^{k+1}=p^ke_1^ke$. Consequently $c(e)\in p^kW_0$ for all $k$.
Separatedness and then injectivity give $e=0$.
\end{proof}
% END IFD3

% BEGIN IFD4
\paragraph{IFD4: detection of every component's constant coefficient.}
\label{claim:IFD4}
In addition to IFD3, let $O_0$ be a discrete valuation domain with
$p\ne0$, let $B_0$ be an $O_0$-algebra, and let
$i_0:\prod_{i\in\mathbb Z/n\mathbb Z}O_0[[T]]\hookrightarrow B_0$
be an injective $O_0$-algebra map.
If a tuple of constants $u=(u_i)$ satisfies
$c(i_0(u))\in N_1$, then $p\mid u_i$ for every $i$.
This is the full generic component-constant implication
\texttt{Rational.\allowbreak const\_\allowbreak dvd}.

\begin{proof}
Suppose $p\nmid u_i$. Total divisibility in the valuation domain gives
$p=u_i d$, and $d$ is a nonunit: if it were a unit, $p$ would divide
$u_i$. If $p$ has valuation $e_p$, some power $d^{N+1}$ is divisible
by $p$. Let $\varepsilon_i$ be the nonzero idempotent which is $1$
in component $i$ and $0$ elsewhere. Multiplication of the assumed
near-relation by $i_0(\varepsilon_i)$ gives
$c(u_i i_0(\varepsilon_i))\in N_1$. Apply IFD3 with
$e=i_0(\varepsilon_i)$ and with the images of $u_i,d$ and
$d^{N+1}/p$. It gives $i_0(\varepsilon_i)=0$, contrary to injectivity.
The argument uses the actual component idempotent and does not require
the tensor ring $B_0$ to be a domain.
\end{proof}
% END IFD4

% BEGIN IFD5
\paragraph{IFD5: the precise integral-series normal form.}
\label{claim:IFD5}
Let $A_0=\prod_iO_0[[T]]$, with ring maps $i_0:A_0\to B_0$ and
$c:B_0\to W_0$. Suppose $t=c(i_0(T))$ is a unit and suppose the
constant-coefficient conclusion of IFD4 holds, without requiring its
particular proof hypotheses. Then
\[
 c(i_0(f))\in N_m\quad\Longrightarrow\quad
                   f=T^mf_1+pf_2\quad(f_1,f_2\in A_0).        \tag{IFD.6}
\]
For any sequence $b_k\in A_0$ there is also a single $\sigma\in A_0$
such that
\[
             \sigma-\sum_{k<m}T^kb_k\in T^mA_0
                                  \quad\hbox{for every }m.  \tag{IFD.7}
\]

\begin{proof}
For (IFD.6), induct on $m$. At $m=0$ take $f_1=f,f_2=0$.
At the next step the inductive expression is $f=T^mf_1+pf_2$.
Use IFD1 to cancel $t^m$ from the stronger assumption in $N_{m+1}$;
it gives $c(i_0(f_1))\in N_1$. Write $f_1=Tg+u$, where $u$ is
its tuple of constant coefficients. Since $c(i_0(Tg))\in N_1$,
so is $c(i_0(u))$. The assumed criterion gives $u=pu'$.
Thus $f=T^{m+1}g+p(f_2+T^mu')$, proving the induction.

For (IFD.7), define the coefficient of $T^r$ in component $i$ of
$\sigma$ to be the coefficient in $\sum_{k\leq r}T^kb_k(i)$.
This is a finite sum. For $r<m$ that coefficient equals the one in
$\sum_{k<m}T^kb_k(i)$, since the additional terms have degree at
least $m$. Their difference therefore has its first $m$ coefficients
zero, which is exactly divisibility by $T^m$ in every component.
No analytic convergence is used in this formal-series sum.
\end{proof}
% END IFD5

% BEGIN IFD6
\paragraph{IFD6: full near-parameter closedness modulo the prime.}
\label{claim:IFD6}
Keep the hypotheses of IFD5. Let $a_0:A_0\to E_0$ and
$j_0:E_0\to W_0$ be ring maps with $j_0a_0=ci_0$.
Assume that for every $e\in E_0$ there are $f\in A_0$, $k\geq0$
and $r\in E_0$ with
\[
                     e a_0(T)^k=a_0(f)+pr,                 \tag{IFD.8}
\]
and that $\bigcap_mN_m\subseteq pW_0$. Then
\[
 \bigl(\forall m\ \exists e_m\in E_0:\ x-j_0(e_m)\in N_m\bigr)
          \quad\Longrightarrow\quad x\in j_0(E_0)+pW_0.    \tag{IFD.9}
\]
This is the full original generic \texttt{near\_\allowbreak closed} statement.

\begin{proof}
First, if $j_0(e)\in N_m$, (IFD.8) implies
$c(i_0(f))\in N_{m+k}$. IFD5 gives
$f=T^{m+k}f_1+pf_2$. After applying the maps and dividing by the
unit $t^k$, we get the useful normal form
\[
                     j_0(e)=c(i_0(T^mf_1))+pr'.             \tag{IFD.10}
\]
Choose the approximants $e_m$ in (IFD.9). Their differences satisfy
$j_0(e_{m+1}-e_m)\in N_m$. Apply (IFD.10) to write
\[
                     j_0(e_{m+1}-e_m)=c(i_0(T^mb_m))+pr_m.
\]
Choose the single formal sum $\sigma$ from (IFD.7).
For every $M$, telescoping the displayed differences shows that
$x-j_0(e_0)-c(i_0(\sigma))$ is the sum of $x-j_0(e_M)$,
a term in $t^Mc(B_0)$, and a term in $pW_0$.
It therefore belongs to every $N_M$ and hence to $pW_0$.
Since $c(i_0(\sigma))=j_0(a_0(\sigma))$, this proves (IFD.9).
In particular this step has not assumed completeness of $E_0$
for the parameter $T$.
\end{proof}
% END IFD6

% BEGIN IFD7
\paragraph{IFD7: the actual image is fixed, has bounded poles, and is near-closed.}
\label{claim:IFD7}
For the original maps fixed above,
\[
 j(E)\subseteq W^H,\qquad \phi(j(E))\subseteq j(E),\qquad
 \forall e\in E\ \exists k:\ t^kj(e)\in\mathcal N_0.       \tag{IFD.11}
\]
Moreover
\[
 \forall m\ \exists e_m:\ x-j(e_m)\in\mathcal N_m
               \quad\Longrightarrow\quad x\in j(E)+pW.    \tag{IFD.12}
\]
All the original \texttt{io\allowbreak E\_\allowbreak fixed}, \texttt{bounded\_\allowbreak pole} and
\texttt{closed} conclusions hold.

\begin{proof}
For $h\in H$, the specified character identity and the injection
$O\hookrightarrow LL$ give $\chi_\Gamma(\operatorname{res}_\Gamma h)=1$.
Its stipulated injectivity gives $\operatorname{res}_\Gamma h=1$.
The original $\Gamma$-equivariance of $j$ therefore says $hj(e)=j(e)$.
Its Frobenius identity proves the second assertion.
The record $R.\mathrm{EE\_\allowbreak dense}$, at precision $p$, gives (IFD.8)
for the original $a:A\to E$. Applying $j$ proves the bounded-pole
assertion.

The map $i:A\to B$ is injective: if $i(f)=i(g)$, the identity
$ja=ci$ and WIS18 give $a(f)=a(g)$, and the supplied injection $a$
gives $f=g$. WIS5 gives injectivity and $p$-saturation of $c$,
and absence of $p$-torsion and $p$-adic separation in $W$.
The image of $T$ is a unit because $a(T)$ is a unit in the original
record. WIS24 gives $\bigcap_m\mathcal N_m=pW$.
Finally WIS1 makes $O$ a DVR with $p\ne0$.
Thus IFD4 supplies the component-constant criterion for these exact
maps, and IFD6 proves (IFD.12). The same reasoning proves the generic
\texttt{closed} implication with its original printed hypotheses:
each actual WIS input is simply replaced by the corresponding supplied
injectivity, saturation, torsion, separation and intersection premise.
\end{proof}
% END IFD7

\subsection{Invariant tilt approximation by the original generators}

For this subsection $K$ is any finite extension, with no
$\mathrm{Good}$ assumption. Let $K_0$ be its maximal unramified
subfield, $K_r=K_0(\zeta_{p^r})$ and
$K_\infty=\bigcup_rK_r$. Write $H_\infty=G_{K_\infty}$.
The roots here are the original compatible roots in $AI$.
Let $G_0\subseteq S$ be the subring generated by all $\epsilon_r$
and by the prime-to-$p$ root tilts in the source:
\[
 G_0=\mathbb Z[\,\epsilon_r\ (r\geq0),\ x_b:
  \operatorname{IsTiltOf}(u,m,x_b),\ u\in K_0,\ u^m=1,
                  \ m>0,\ p\nmid m\,].                    \tag{IFD.13}
\]
As a subring of a characteristic-$p$ ring, the indicated integer
scalars of course factor through $\mathbb F_p$.
No closure for a topology is part of this definition.

% BEGIN IFD8
\paragraph{IFD8: sharp-root coordinates preserve the exact fixed group.}
\label{claim:IFD8}
For $x\in S$ write $x^{\sharp,N}\in\mathcal O_C$ for its exact
$N$th sharp-root coordinate. For every original Galois element $g$,
\[
        (g x)^{\sharp,N}=g(x^{\sharp,N}).                    \tag{IFD.14}
\]
If $x$ is fixed by $H_\infty$, then each $x^{\sharp,N}$ belongs
to the completion $\widehat{K_\infty}$ and has norm at most one.
If the $N$th inverse-limit coordinate of $z\in S$ is zero, then
\[
                            |z|_\flat\leq |p|^{p^N}.        \tag{IFD.15}
\]

\begin{proof}
Choose arbitrary integral representatives $a_r$ of the inverse-limit
coordinates $x_r\in\mathcal O_C/p$. The coordinate formula is
\[
                 x^{\sharp,N}=\lim_{r\to\infty}a_{N+r}^{p^r}.
                                                               \tag{IFD.16}
\]
Here is the estimate which makes both the limit and its independence
of representatives precise. If $|u|,|v|\leq1$ and $|u-v|\leq|p|^k$
with $k\geq1$, the binomial expansion gives
$|u^p-v^p|\leq|p|^{k+1}$. The intermediate binomial coefficients
are divisible by $p$, and the last term has valuation at least
$pk\geq k+1$. Iterating gives
$|u^{p^r}-v^{p^r}|\leq|p|^{k+r}$.
Since $a_{r+1}^p\equiv a_r\pmod p$, successive terms of (IFD.16)
differ by at most $|p|^{r+1}$. They are Cauchy; changing any
representatives has the same vanishing bound. This includes $p=2$.

The original field \texttt{act\allowbreak T\_\allowbreak coeff} says that $g(a_r)$ is a
representative of $(gx)_r$. The original \texttt{act\allowbreak C\_\allowbreak norm} says
that $g$ is an isometry of $C$, hence continuous. Apply it to
(IFD.16) and pass through the limit to obtain (IFD.14).
This uses only the original action fields, not an additional choice
of action on a completion.

If $x$ is $H_\infty$-fixed, (IFD.14) fixes each sharp-root
coordinate. LS1, applied to this same isometric extension of the
action on $\Omega$, gives
$C^{H_\infty}=\widehat{\Omega^{H_\infty}}
=\widehat{K_\infty}$. The last equality is the algebraic
fixed-field correspondence for $H_\infty=G_{K_\infty}$.
The coordinates in (IFD.16) are integral, so the limit has norm at
most one. Finally $z^{\sharp,N}$ reduces to the $N$th coordinate
of $z$. If it is zero modulo $p$, its norm is at most $|p|$.
The exact relation $(z^{\sharp,N})^{p^N}=z^\sharp$ proves (IFD.15).
No additivity of the sharp map has been asserted or used.
\end{proof}
% END IFD8

% BEGIN IFD9
\paragraph{IFD9: integral coefficients in the finite cyclotomic layers.}
\label{claim:IFD9}
The nonzero norms in $K_0$ are integral powers of $p$.
Every integral element of $K_0$ is congruent modulo $p$ to zero or
to a root of unity in $K_0$ of order prime to $p$.
For every $r\geq1$,
\[
                 \mathcal O_{K_r}=\mathcal O_{K_0}[\zeta_{p^r}].
                                                               \tag{IFD.17}
\]

\begin{proof}
We recall the elementary unramified facts in the form needed here.
For a finite residue field $\mathbb F_{p^f}$, lift a monic
irreducible polynomial for one residue generator to $\mathbb Z_p[X]$.
Its unramified integral extension is finite free over $\mathbb Z_p$;
the residue is $\mathbb F_{p^f}$, and its maximal ideal is generated
by $p$. To see the last assertion and the norm statement directly,
expand an element in the lifted basis. The smallest common power
of $p$ in its coefficients leaves a nonzero residue. Such a residue
is a unit: lift its inverse and sum the geometric series for the
error in $p$. Completeness holds coordinatewise. Thus this ring is
a DVR with uniformizer $p$.

Every nonzero residue has a unique lift satisfying $u^{p^f-1}=1$.
For existence, Newton correction for $X^{p^f-1}-1$ is integral
and doubles the error valuation, because its derivative is a unit
at the chosen residue. The corrections converge by completeness.
For uniqueness, if two such roots have the same residue, division
of their power difference by their difference gives a unit
(its residue is the nonzero derivative), so the difference is zero.
These roots generate the unramified extension. Conversely any
prime-to-$p$ root of unity lies in one of these extensions: choose
$f$ with its order dividing $p^f-1$ and use the unique simple-root
lifts of all the roots in the finite residue field. The union of
these fields is the original maximal unramified extension.
Its finite subfield $K_0$ is therefore unramified, with the stated
norm group and residue representatives. Equivalently, finite
subfields correspond to the subfields of the finite residue fields;
the lifted residue basis gives the same degree and the same
uniformizer $p$.

Let $e=p^{r-1}(p-1)$ and $\pi=\zeta_{p^r}-1$.
The elementary cyclotomic product calculation in LS2 gives
$|\pi|=p^{-1/e}$: primitive conjugates have the same distance from
$1$ and their product at $1$ is $p$. For $a_i\in K_0$ the nonzero
terms $a_i\pi^i$, $0\leq i<e$, have distinct valuations modulo
$\mathbb Z$. Hence they cannot cancel. In particular these $e$
powers are linearly independent over $K_0$. The cyclotomic
polynomial gives the opposite degree bound, so they are a basis
of $K_r$ over $K_0$.

Write an integral $z\in K_r$ in this basis. The norm of its sum is
the largest of the distinct term norms, so
$|a_i|p^{-i/e}\leq1$ for every $i$. Since the valuation of $a_i$
is an integer and $0\leq i/e<1$, each $a_i$ is integral.
Thus $z\in\mathcal O_{K_0}[\pi]=\mathcal O_{K_0}[\zeta_{p^r}]$.
The reverse inclusion follows from integrality of the generator,
proving (IFD.17), including $p=2,r=1$ where $e=1$.
\end{proof}
% END IFD9

% BEGIN IFD10
\paragraph{IFD10: every completed integral element has a generator lift.}
\label{claim:IFD10}
For every $N\geq0$ and $z\in\widehat{K_\infty}$ with $|z|\leq1$,
there exists $w\in G_0$ such that
\[
                  w_N=z\bmod p\quad\hbox{in }\mathcal O_C/p.
                                                               \tag{IFD.18}
\]
The same coefficient index $N$ and the same unramified field $K_0$
are used throughout.

\begin{proof}
Let $U_N$ be the inverse image in $\mathcal O_C$ of the image of
$G_0$ under its $N$th coefficient map. It is a subring containing
$p\mathcal O_C$. It is therefore closed under replacing an element
by any other integral element congruent to it modulo $p$.

The $N$th coefficient of $\epsilon_j$ is
$\zeta_{p^{j+N}}\bmod p$. Raising to $p^N$ shows that
$\zeta_{p^j}\in U_N$ for every $j$.
If $u\in K_0$, $u^m=1$, $p\nmid m$, choose an integer
$a$ with $ap\equiv1\pmod m$. The compatible sequence
\[
                           u^{p^N a^k}\qquad(k\geq0)
\]
defines a tilt of $u^{p^N}$ of order dividing $m$ and has $N$th
coordinate $u$. It is one of the generators allowed in (IFD.13).
Thus $u\in U_N$. IFD9 and congruence stability now imply
$\mathcal O_{K_0}\subseteq U_N$ and then
$\mathcal O_{K_j}\subseteq U_N$ for every $j\geq1$.

Choose $z'\in K_\infty$ with $|z-z'|<|p|$. The ultrametric
inequality and $|z|\leq1$ give $|z'|\leq1$. It lies in some finite
layer, so $z'\in U_N$. Congruence stability gives $z\in U_N$.
This is exactly (IFD.18). The construction requires only membership
in the subring generated by the specified elements; it does not
silently replace that subring by its completion.
\end{proof}
% END IFD10

% BEGIN IFD11
\paragraph{IFD11: the full original invariant tilt approximation.}
\label{claim:IFD11}
If $x\in S$ is fixed by $H_\infty$, then for every $m\geq0$
there exists $w\in G_0$ such that
\[
                           (\epsilon-1)^m\mid x-w
                                      \quad\hbox{in }S.    \tag{IFD.19}
\]
This proves the original \texttt{Tilt\allowbreak Invariants.\allowbreak invariant\_\allowbreak approx}.

\begin{proof}
The element $\epsilon-1$ is nonzero. Indeed its sharp-root
coordinates reduce to $\zeta_{p^r}-1$, whose norm is greater than
$|p|$ for some $r$ by IFD9; their reductions cannot all be zero.
Set $b=(\epsilon-1)^m$, which is nonzero also for $m=0$.
Choose $N$ so large that $|p|^{p^N}\leq|b|_\flat$.
By IFD8 the exact root $x^{\sharp,N}$ belongs to
$\widehat{K_\infty}\cap\mathcal O_C$.
IFD10 gives $w\in G_0$ whose $N$th coefficient equals its
reduction modulo $p$, hence equals the $N$th coefficient of $x$.
The $N$th coefficient of the tilt difference $x-w$ is zero.
IFD8 therefore gives $|x-w|_\flat\leq|p|^{p^N}\leq|b|_\flat$.
The tilt valuation-ring property proved in LS4 says exactly that
this norm inequality is divisibility by $b$ in $S$.
This proves (IFD.19). The sharp map is used only for norms,
limits and multiplicative root relations, never as an additive map.
\end{proof}
% END IFD11

% BEGIN IFD12
\paragraph{IFD12: fractional invariant approximation with a common pole.}
\label{claim:IFD12}
Let $y\in F$ be fixed by $H_\infty$, and keep any $s,m\geq0$.
There is an $N_0$ such that for every $N\geq N_0$ one can find
$u\in G_0$, $v\in S$ with
\[
                         \tau_s^N y=u+\tau_s^{N+m}v.
                                                               \tag{IFD.20}
\]
This is the complete original \texttt{flat\_\allowbreak invariant\_\allowbreak approx} scope.

\begin{proof}
The element $\tau_s$ is nonzero, has norm less than one, and is
fixed by $H_\infty$: its inverse-limit coordinates are the
cyclotomic roots, all fixed by that group, so the coordinate action
fixes it. The norm and nonvanishing follow from IFD9 and the same
coordinate test as IFD11. Thus for all sufficiently large $N$,
$\tau_s^N y$ is integral; it is also fixed.
Apply IFD11 to it with exponent $N+m$.
In the characteristic-$p$ ring $S$,
$\epsilon-1=(\epsilon_s-1)^{p^s}=\tau_s^{p^s}$.
Consequently divisibility by $(\epsilon-1)^{N+m}$ implies
divisibility by $\tau_s^{N+m}$, giving (IFD.20).
The threshold $N_0$ depends only on $y$ and $s$, and every larger
$N$ is allowed. A single maximum therefore works for any finite
list of fractional coordinates, as required below.
\end{proof}
% END IFD12

% BEGIN IFD13
\paragraph{IFD13: the correct fixed coordinates in a good layer.}
\label{claim:IFD13}
Under the original $\mathrm{Good}(K,s,LL)$, if $y\in W^H$ then
every $(y_\alpha)_0\in F$ is fixed by $H_\infty$.

\begin{proof}
The original good-layer condition includes $K\subseteq K_s$.
Every element of $H_\infty$ fixes $K_s$, hence $K$, and fixes
all cyclotomic roots, hence has cyclotomic character $1$.
Therefore $H_\infty\subseteq H$.
The action on the $O$ factor is the identity. In the fixed finite
basis, applying an element of $H_\infty$ to the equality $hy=y$
gives equality of each Witt coordinate, and then of its zeroth
coordinate. These are the actual actions on $W(F)$ and $F$.
No replacement of $K$ by an unramified extension has occurred.
\end{proof}
% END IFD13

% BEGIN IFD14
\paragraph{IFD14: elements eventually in the coefficient image form a subring.}
\label{claim:IFD14}
Let $j_0:E_0\to W_0$ be a ring map and $\phi_0:W_0\to W_0$
a ring endomorphism with $\phi_0(j_0(E_0))\subseteq j_0(E_0)$.
Then
\[
 \mathcal P=\{x\in W_0:\ \exists r\geq0,\
                    \phi_0^r(x)\in j_0(E_0)+pW_0\}         \tag{IFD.21}
\]
is a subring. If its defining condition holds for $r$, it holds
for every larger exponent.

\begin{proof}
A ring endomorphism fixes the integer $p$. Applying $\phi_0$
to $j_0(e)+pw$ again gives an element of $j_0(E_0)+pW_0$,
which proves persistence. Zero, one and negatives satisfy (IFD.21).
For two elements satisfying it at exponents $r,r'$, pass to
$\max(r,r')$. Their images then lie in the subring
$j_0(E_0)+pW_0$, which is closed under addition and multiplication.
The sum and product therefore satisfy the same condition.
This verifies the exact original \texttt{perf} subring without
assuming that $j_0(E_0)$ is already perfect.
\end{proof}
% END IFD14

% BEGIN IFD15
\paragraph{IFD15: approximation inside any subring containing the generators.}
\label{claim:IFD15}
Let $T_0\subseteq W$ be a subring containing the image of $O$,
$t^{-1}$, and the elements $1\otimes[x]$ for every generator $x$
in (IFD.13). Under the original good-layer condition, for each
$y\in W^H$ and $m\geq0$ there is $y_0\in T_0$ with
$y-y_0\in\mathcal N_m$.

\begin{proof}
Consider the inverse image of $T_0$ in $W(S)$ under
$f_W:w\mapsto1\otimes w$, followed by its image under the
zeroth-coordinate ring map to $S$. This image is a subring of $S$
containing every generator of $G_0$. Thus every $u\in G_0$ has
a lift $w\in W(S)$ with $w_0=u$ and $f_W(w)\in T_0$.
This step does not claim that Teichmuller lifting is additive.

By IFD13 the finitely many coordinates $(y_\alpha)_0$ are
$H_\infty$-fixed. Choose a common $N$ above all the thresholds
in IFD12. It gives $u_\alpha\in G_0,v_\alpha\in S$ with
\[
       \tau_s^N(y_\alpha)_0
                      =u_\alpha+\tau_s^{N+m}v_\alpha.
\]
Choose lifts $w_\alpha$ in the preceding inverse image and put
$z=\sum_\alpha b_\alpha\otimes w_\alpha$ in $W$.
Each term is the product of the image of $b_\alpha\in O$ and
$f_W(w_\alpha)$, so $z\in T_0$. IFD1's exact coordinate
criterion gives $t^Ny-z\in\mathcal N_{N+m}$.
Set $y_0=t^{-N}z\in T_0$ and use (IFD.3) to obtain
$y-y_0\in\mathcal N_m$. This proves the original
\texttt{approx\_\allowbreak subring} implication for every supplied $T_0$.
\end{proof}
% END IFD15

% BEGIN IFD16
\paragraph{IFD16: the actual invariant approximation has eventual coefficient values.}
\label{claim:IFD16}
For every $y\in W^H$ and $m\geq0$ there are $r\geq0$,
$y_0\in W$, $e\in E$ and $v\in W$ with
\[
          y-y_0\in\mathcal N_m,\qquad \phi^r(y_0)=j(e)+pv.
                                                               \tag{IFD.22}
\]
All the original \texttt{Rational.\allowbreak approx} hypotheses and outputs
are retained.

\begin{proof}
Use the subring $\mathcal P$ of IFD14 for the actual $j,\phi$.
It contains $j(E)$, hence the image of $O$ and $t^{-1}$.
If $x=\epsilon_r$, then the Witt Frobenius identity gives
\[
 \phi^r(1\otimes[\epsilon_r])
       =1\otimes[\epsilon]
       =(1+t)^{p^s}=j\bigl((1+a(T))^{p^s}\bigr).
\]
Thus these generators belong to $\mathcal P$.
For a prime-to-$p$ root tilt $x_b$ of $u\in K_0$, each
$\Gamma.\mathrm{emb}_i(u)$ is integral in $LL$, since its positive
power is $1$. The original \texttt{Io\allowbreak W} identity identifies
$1\otimes[x_b]$ with the image under $i$ of the tuple of these
integral constants. Applying $c$ identifies it with an element of
$j(E)$. It too belongs to $\mathcal P$.
IFD15 now approximates $y$ in $\mathcal P$ at the prescribed
precision. Membership in $\mathcal P$ supplies $r,e,v$ in
(IFD.22). There is no assertion that $y_0$ itself is fixed; the
later argument requires only the fixedness of $y$ and the displayed
approximation.
\end{proof}
% END IFD16

\subsection{An additive contraction, followed by prime-adic lifting}

% BEGIN IFD17
\paragraph{IFD17: one prime-adic step for the kernel.}
\label{claim:IFD17}
We state the exact general setting used in IFD17--IFD20.
Let $N$ be an abelian group, $p$ a natural number,
$\Phi:N\to N$ additive, $D\subseteq M\subseteq N$ subgroups,
and $F_m\subseteq N$ a decreasing sequence of subgroups.
Assume the following properties:
\begin{enumerate}
\item $\Phi(M)\subseteq M$, $\Phi(D)\subseteq D$ and
      $\Phi(F_m)\subseteq F_{m+1}$.
\item $M$ and $D$ are $p$-saturated in $N$: if $px$ belongs
      to either subgroup, so does $x$.
\item If $f_i\in M\cap F_i$ for all $i$, there is $s\in M$
      with $s-\sum_{i<m}f_i\in F_m$ for every $m$.
\item If $x$ can be approximated by $D$ modulo every $F_m$,
      then $x\in D+pN$.
\item For each $y\in M$, there are $j\geq0$ and $d\in D$
      with $\Phi^jy-d\in F_0$.
\item Every sequence $x_i\in M$ has a sum $s\in M$ satisfying
      $s-\sum_{i<k}p^ix_i\in p^kN$ for all $k$.
\item If $x\in D+p^kN$ for every $k$, then $x\in D$.
\end{enumerate}
If $z\in M$ and $\Phi z-z\in D$, then there are $d\in D$,
$r\in M$ with
\[
                         z-d=pr,\qquad \Phi r-r\in D.      \tag{IFD.23}
\]

\begin{proof}
Put $\delta=\Phi z-z\in D$.
For every $j$, telescoping gives
$\Phi^jz-z=\sum_{i<j}\Phi^i\delta\in D$.
Property 5 therefore provides an approximation $z-d_0\in F_0$.
If $z-d_m\in F_m$, put $d_{m+1}=\Phi d_m-\delta\in D$.
Then $z-d_{m+1}=\Phi(z-d_m)\in F_{m+1}$.
Thus $z$ is approximable by $D$ at every filtration level.
Property 4 gives $z-d=pr$ for some $d\in D,r\in N$.
Since $pr\in M$, property 2 gives $r\in M$.
Finally
\[
                p(\Phi r-r)=(\Phi z-z)-(\Phi d-d)\in D.
\]
The saturation of $D$ gives the second part of (IFD.23).
Only the displayed hypotheses have been used; no topology on the
abstract group has been silently imposed.
\end{proof}
% END IFD17

% BEGIN IFD18
\paragraph{IFD18: injectivity on the quotient by the coefficient subgroup.}
\label{claim:IFD18}
Under the hypotheses of IFD17,
\[
                 z\in M,\quad \Phi z-z\in D
                                    \quad\Longrightarrow\quad z\in D.
                                                               \tag{IFD.24}
\]

\begin{proof}
Apply IFD17 repeatedly. Inductively one obtains
$z-d_k=p^kr_k$ with $d_k\in D$, $r_k\in M$ and
$\Phi r_k-r_k\in D$. At $k=0$ take $d_0=0,r_0=z$.
If $r_k-e_k=pr_{k+1}$ is the next step, set
$d_{k+1}=d_k+p^ke_k$. Then
$z-d_{k+1}=p^{k+1}r_{k+1}$ with the same required properties.
Property 7 of IFD17 now gives $z\in D$.
Equivalently, the map induced by $\Phi-1$ on $M/D$ is injective.
\end{proof}
% END IFD18

% BEGIN IFD19
\paragraph{IFD19: one prime-adic step for the cokernel.}
\label{claim:IFD19}
Under the same hypotheses, for every $y\in M$ there exist
$z\in M$, $d\in D$ and $r\in M$ such that
\[
                         y-(\Phi z-z)-d=pr.                \tag{IFD.25}
\]

\begin{proof}
Choose $j,d_0$ with $e=\Phi^jy-d_0\in F_0$ using property 5.
It also belongs to $M$. The terms $\Phi^ie$ belong to $M\cap F_i$.
Property 3 supplies $s\in M$ with
$s-S_m\in F_m$, where $S_m=\sum_{i<m}\Phi^ie$.
The finite telescoping identity
\[
                          \Phi S_m-S_m=\Phi^me-e
\]
gives
\[
 \Phi s-s+e=(\Phi(s-S_m)-(s-S_m))+\Phi^me\in F_m.
\]
Here $\Phi(F_m)\subseteq F_{m+1}\subseteq F_m$ has been used.
Property 4 applied with the zero approximant in $D$ gives
$\Phi s-s+e=d_1+pr$ for some $d_1\in D,r\in N$.
Set
\[
                     z=-\left(\sum_{i<j}\Phi^iy+s\right)\in M.
\]
A second telescoping calculation gives
$y-(\Phi z-z)=d_0+(\Phi s-s+e)=d_0+d_1+pr$.
This proves (IFD.25), with $d=d_0+d_1$.
Its left side belongs to $M$, so saturation puts $r$ in $M$.
The construction is the ordinary geometric-series argument, with
the sign determined by solving for $\Phi-1$.
\end{proof}
% END IFD19

% BEGIN IFD20
\paragraph{IFD20: the full original additive quotient bijectivity theorem.}
\label{claim:IFD20}
Under all the hypotheses printed in IFD17,
\[
                         \Phi-1:M/D\longrightarrow M/D
                                      \quad\hbox{is bijective}.       \tag{IFD.26}
\]
Explicitly, for every $y\in M$ there is $z\in M$ with
$y-(\Phi z-z)\in D$, and the kernel assertion is (IFD.24).
This is the full original
\texttt{sub\_\allowbreak one\_\allowbreak bijective\_\allowbreak quotient\_\allowbreak of\_\allowbreak top\allowbreak Nilpotent}
at its original additive-group scope.

\begin{proof}
Starting with $y_0=y$, repeatedly apply IFD19 to obtain
\[
                  y_i-(\Phi z_i-z_i)-d_i=p y_{i+1},
                   \qquad z_i,y_{i+1}\in M,\ d_i\in D.
\]
For every $k$ the finite sum satisfies
\[
 y-\sum_{i<k}p^i(\Phi z_i-z_i)-\sum_{i<k}p^id_i=p^ky_k.
\]
Property 6 gives $z\in M$ with
$z-\sum_{i<k}p^iz_i=p^kr_k$ for each $k$.
Since $\Phi$ is additive, it commutes with integer multiplication.
Subtracting $(\Phi-1)$ of this last equality from the preceding
one yields
\[
 y-(\Phi z-z)-\sum_{i<k}p^id_i
                         =p^k\bigl(y_k-(\Phi r_k-r_k)\bigr).
\]
The sum of the $d_i$ terms lies in $D$. Property 7 therefore gives
$y-(\Phi z-z)\in D$. This proves surjectivity; IFD18 proves
injectivity. Both statements descend to the quotient because
$\Phi(D)\subseteq D$.
\end{proof}
% END IFD20

\subsection{The original matrix theorem and its actual specialization}

% BEGIN IFD21
\paragraph{IFD21: a finite coefficient matrix has one contracting filtration.}
\label{claim:IFD21}
Let $c:B_0\to W_0$, $t=c(\tau)$ and $N_m=t^mc(B_0)+pW_0$
be as in IFD1. Let $\phi_0$ preserve $c(B_0)$ and satisfy
$\phi_0(t)\in N_2$. Let $j_0:E_0\to W_0$ have the property
that each $j_0(e)$ has a bounded pole:
$t^kj_0(e)\in N_0$ for some $k$.
For every finite square matrix $P\in M_d(E_0)$ there is a
single $c_0\geq0$ such that the additive operator
\[
                   \Phi(v)=j_0(P)\phi_0(v)
\]
satisfies
\[
 F_m=(N_{m+c_0+1})^d,\qquad \Phi(F_m)\subseteq F_{m+1}.
                                                               \tag{IFD.27}
\]

\begin{proof}
Choose $c_0$ above all pole bounds for the finitely many entries
of $P$; for the empty matrix take $c_0=0$.
Increasing a pole bound preserves it, by multiplying the integral
coefficient by a power of $\tau$.
From $\phi_0(t)\in N_2$ and multiplicativity one gets
$\phi_0(N_r)\subseteq N_{2r}$: apply $\phi_0$ to
$t^rc(b)+pw$, use its preservation of $c(B_0)$, and expand the
$r$th power of the relation in $N_2$.

We also need the following direct product calculation, which does
not require $t$ to be a unit. If $t^{c_0}q=c(b)+pu$ and
$v=t^{r+c_0}c(a)+pw$, then
\[
                         qv=t^rc(ba)+p(t^rc(a)u+qw)\in N_r.
\]
For $v\in F_m$, its Frobenius coordinates are in
$N_{2(m+c_0+1)}\subseteq N_{m+1+2c_0+1}$.
The product calculation lowers the exponent by $c_0$ when
multiplied by any entry of $j_0(P)$. Summing the finitely many
entries gives (IFD.27). In particular $\Phi$ preserves $F_0$.
\end{proof}
% END IFD21

% BEGIN IFD22
\paragraph{IFD22: eventual smallness uses one common iterate.}
\label{claim:IFD22}
Keep IFD21, and assume $\phi_0(j_0(E_0))\subseteq j_0(E_0)$.
Let $S_0\subseteq W_0$ be a subring containing $j_0(E_0)$.
Assume that for every $y\in S_0$ and $m$ there is $y_0$ with
$y-y_0\in N_m$ and $\phi_0^r(y_0)\in j_0(E_0)+pW_0$ for
some $r$. Then for every vector $y\in S_0^d$ there are $r$
and $d_0\in j_0(E_0)^d$ with
\[
                              \Phi^r y-d_0\in F_0.         \tag{IFD.28}
\]

\begin{proof}
Approximate each coordinate at precision $c_0+1$, and take the
maximum $r$ of their finitely many iteration exponents.
Persistence in IFD14 puts every $\phi_0^r((y_0)_i)$ in
$j_0(E_0)+pW_0$. There is a matrix $P_r$ over $E_0$ with
\[
                              \Phi^r(v)=j_0(P_r)\phi_0^r(v).
\]
Indeed $P_0=1$. For the induction, choose preimages in $E_0$ of
the finitely many entries of $\phi_0(j_0(P_r))$, and multiply
the resulting matrix on the left by $P$. This proves the formula
without requiring an inverse to $\phi_0$ or injectivity of $j_0$.
It follows that $\Phi^r(y_0)\in j_0(E_0)^d+pW_0^d$.
Meanwhile $y-y_0\in F_0$, and IFD21 shows
$\Phi^r(y-y_0)\in F_0$. Since $pW_0^d\subseteq F_0$, their
sum proves (IFD.28). Empty coordinate sets give the same statement
with $r=0$.
\end{proof}
% END IFD22

% BEGIN IFD23
\paragraph{IFD23: the full generic matrix quotient theorem.}
\label{claim:IFD23}
Take commutative rings $W_0,E_0,B_0$, maps $j_0,c$, a ring
endomorphism $\phi_0$, a subring $S_0\subseteq W_0$, and
$t=c(\tau)$, with the filtration $N_m$ of IFD1.
Assume precisely the following original matrix hypotheses:
\begin{enumerate}
\item $j_0(E_0)\subseteq S_0$, and $\phi_0$ preserves $S_0$,
      $j_0(E_0)$ and $c(B_0)$.
\item $\phi_0(t)\in N_2$, and every $j_0(e)$ has a bounded
      pole as in IFD21.
\item $S_0$ and $j_0(E_0)$ are $p$-saturated in $W_0$.
\item For every $k$ and sequence $f_i\in S_0\cap N_{i+k}$,
      there is $s\in S_0$ with
      $s-\sum_{i<m}f_i\in N_{m+k}$ for every $m$.
\item An element approximable by $j_0(E_0)$ modulo every
      $N_m$ belongs to $j_0(E_0)+pW_0$.
\item The eventual-coefficient approximation hypothesis of
      IFD22 holds for every element of $S_0$.
\item Every sequence in $S_0$ has a $p$-weighted sum with the
      congruences in property 6 of IFD17; and $j_0(E_0)$ is
      $p$-adically closed in the sense of property 7 there.
\end{enumerate}
For every $d\geq0$ and $P\in M_d(E_0)$, the map
\[
             j_0(P)\phi_0-1:S_0^d/j_0(E_0)^d
                                    \longrightarrow S_0^d/j_0(E_0)^d
                                                               \tag{IFD.29}
\]
is bijective as a map of additive groups. No invertibility of $P$
is required. This is the full original
\texttt{matrix\_\allowbreak frob\_\allowbreak sub\_\allowbreak one\_\allowbreak bijective} implication.

\begin{proof}
Apply IFD20 to the ambient additive group $N=W_0^d$,
$M=S_0^d$, $D=j_0(E_0)^d$, the operator
$\Phi=j_0(P)\phi_0$ and the subgroups $F_m$ in IFD21.
The first hypothesis gives preservation of $M,D$.
IFD21 gives contraction and decreasingness of $F_m$.
Saturation is coordinatewise. Hypothesis 4, with $k=c_0+1$,
supplies the required sums in $M\cap F_m$ one coordinate at a time.
If $x$ is approximable by $D$ modulo every $F_m$, it is
approximable modulo every $N_m$ in each coordinate, because
$N_{m+c_0+1}\subseteq N_m$. Hypothesis 5 gives $x\in D+pN$.
IFD22 gives eventual smallness. Finally the weighted $p$-sums
and closedness in hypothesis 7 apply coordinatewise.
These check every hypothesis of IFD17, so IFD20 proves (IFD.29).
All choices are finite coordinate choices, including the maximum
in IFD22. No extra uniformity over an infinite matrix has entered.
\end{proof}
% END IFD23

% BEGIN IFD24
\paragraph{IFD24: the actual fixed-vector descent and its surjectivity companion.}
\label{claim:IFD24}
For every original good layer and supplied $\Gamma,R,I,\mathrm{IoW}$,
and every $d\geq0$, $P_0\in M_d(E)$, the additive map
\[
              j(P_0)\phi-1:(W^H)^d/j(E)^d
                                      \longrightarrow(W^H)^d/j(E)^d
                                                               \tag{IFD.30}
\]
is bijective. In particular, if $z\in(W^H)^d$ and
$j(P_0)\phi(z)-z\in j(E)^d$, then every coordinate of $z$
belongs to $j(E)$. This includes the full original
\texttt{Rational.\allowbreak vec} theorem.

\begin{proof}
Use IFD23 with the unchanged actual rings and $S_0=W^H$.
IFD7 gives containment of $j(E)$, its Frobenius stability and
the bounded-pole property. WIS20 gives stability of $W^H$,
its $p$-saturation and its weighted $p$-sums.
IFD2 gives stability of $c(B)$ and $\phi(t)\in\mathcal N_2$.
WIS18 gives $p$-saturation of $j(E)$ in $W$: if $px=j(e)$,
reflected divisibility gives $e=pe'$, and absence of $p$-torsion
gives $x=j(e')$.

Since $t=j(a(T))$, IFD7 says $t\in W^H$.
WIS25 therefore supplies the full fixed-ring near-parameter sums
for every shift $k$, using this exact $H$ and this exact $t$.
IFD7 gives near-closedness modulo $p$ and IFD16 gives eventual
coefficient approximation. WIS19 gives $p$-adic closedness of
$j(E)$. These are precisely the seven groups of hypotheses in
IFD23. It proves (IFD.30) and hence the stated vector descent.
The bijectivity assertion is a valid additional consequence of the
same supplied data; the original matrix consumer uses only its
injectivity half.
\end{proof}
% END IFD24

% BEGIN IFD25
\paragraph{IFD25: full inhomogeneous invariant Frobenius-matrix descent.}
\label{claim:IFD25}
Keep every original prime, field, level and supplied record as above.
Let $a,b\geq0$, $A_0\in M_a(E)$, $A_1\in M_b(E)$,
$Y\in M_{a,b}(E)$ and $U\in M_{a,b}(W)$.
Assume $A_0,A_1$ are invertible, every entry of $U$ is $H$-fixed,
and
\[
                        j(A_0)\phi(U)=Uj(A_1)+j(Y).         \tag{IFD.31}
\]
Then every entry of $U$ belongs to $j(E)$.
This is the complete original \texttt{Rational.\allowbreak main}, equivalently
the exported fixed-cyclotomic-kernel Frobenius-solution descent
theorem, including its inhomogeneous term.

\begin{proof}
Choose $C\in M_b(E)$ with $A_1C=1$. Multiplication of (IFD.31)
on the right by $j(C)$ gives
\[
                   j(A_0)\phi(U)j(C)-U=j(YC).
\]
Index the $ab$ entries of $U$ by pairs $(i,j)$ and define
\[
               P_{(i,j),(k,\ell)}=(A_0)_{ik}C_{\ell j}.
\]
If $z_{(i,j)}=U_{ij}$, the preceding equation says exactly
$j(P)\phi(z)-z=j(YC)$ coordinatewise. IFD24 therefore puts
every $z_{(i,j)}$ in $j(E)$, as asserted.
The zero-dimensional cases are vacuous and are included.
The argument in fact only needs a right inverse for $A_1$;
the original invertibility assumption on $A_0$ is retained at
the stated endpoint and has not been traded for a new condition.
\end{proof}
% END IFD25

% BEGIN IFD26
\paragraph{IFD26: exactly one period-ring supplier remains before Wach existence.}
\label{claim:IFD26}
At the original tied local period data, retain the original
crystalline-lattice realization premise used by WLD26.
Suppose the full actual Wach-prism producer
\texttt{exists\_\allowbreak wach\_\allowbreak prism\_\allowbreak over\_\allowbreak integers} is supplied at every
original good layer, including its maps, scalar identity and both
finite coefficient-coordinate assertions.
Then the full original \texttt{Wach\allowbreak Ex} follows for every permitted
$\Gamma,R,I$ and every supplied original etale module.
The actual Wach prism is the only still-retained period-ring
supplier in this return.

\begin{proof}
WLD26 proves the full original all-layer return from three precisely
listed period-ring suppliers. WIS18 supplies its injectivity and
unit-reflection supplier on the actual $I.\mathrm{ioE}$.
IFD25 with $Y=0$ supplies its homogeneous fixed-matrix descent
for every pair of invertible coefficient matrices and every
rectangular fixed matrix. This is the exact input used in WLD15
and WLD25, with the same coefficient map, Frobenius and kernel.
The remaining prism supplier is the hypothesis printed above.
Thus all three entries of WLD26 are available at each permitted
layer, and its conclusion supplies every field of the original
$\mathrm{NS}$ for every supplied etale module.
No extra compatibility condition between distinct outer layers
is introduced: the original \texttt{Wach\allowbreak Ex} quantifies over them.

This conditional conclusion does not assert existence of the
retained prism, of the original tied period data, or of the
crystalline-realization theorem. It also does not supply the
separate \texttt{Faith\allowbreak H}, \texttt{Et\allowbreak Basis}, Herr class maps or
quantitative comparison needed for the full $\mathrm{HW/Fon/Cmp}$
package. The common quantitative family in
\texttt{GLOBAL\_\allowbreak DENOMINATOR} remains a later obligation.
\end{proof}
% END IFD26

\paragraph{Where the next proof must begin.}
The coefficient inclusion, both completeness structures and the
entire invariant matrix descent are now available on the original
objects. The next local-ring task is the actual Wach prism with
its prescribed maps and scalar identities. After that task, IFD26
returns the full all-layer lattice existence statement. The
comparison and common-family quantitative arguments remain separate
steps toward the global denominator obligation.

\endgroup
\end{document}
